\documentclass[11pt,a4paper]{article}
\usepackage[T1]{fontenc}
\usepackage[utf8]{inputenc}
\usepackage{lmodern,microtype}
\usepackage[a4paper,left=25mm,right=25mm,top=24mm,bottom=25mm,headheight=15pt]{geometry}
\usepackage{amsmath,amssymb,amsthm,mathtools}
\usepackage{booktabs,array,longtable,tabularx,enumitem}
\usepackage{xcolor,fancyhdr,listings,tikz}
\usetikzlibrary{positioning,calc}
\usepackage[unicode=true,hypertexnames=false,colorlinks=true,linkcolor=ink,citecolor=teal,urlcolor=teal]{hyperref}
\definecolor{ink}{HTML}{17374D}
\definecolor{teal}{HTML}{176B75}
\definecolor{quiet}{HTML}{586B78}
\pagestyle{fancy}
\fancyhf{}
\fancyhead[L]{\small\textsc{Prime products and graph sampling}}
\fancyhead[R]{\small Consolidated research study}
\fancyfoot[C]{\thepage}
\renewcommand{\headrulewidth}{0.3pt}
\setlength{\parindent}{1.2em}
\setlength{\parskip}{3pt}
\setlength{\emergencystretch}{2em}
\linespread{1.025}
\raggedbottom
\allowdisplaybreaks[2]
\setcounter{tocdepth}{2}
\setcounter{secnumdepth}{3}
\setlist{topsep=4pt,itemsep=3pt,parsep=0pt,leftmargin=*}
\newtheorem{theorem}{Theorem}[section]
\newtheorem{lemma}[theorem]{Lemma}
\newtheorem{proposition}[theorem]{Proposition}
\newtheorem{corollary}[theorem]{Corollary}
\theoremstyle{definition}
\newtheorem{definition}[theorem]{Definition}
\newtheorem{conjecture}[theorem]{Conjecture}
\newtheorem{hypothesis}[theorem]{Hypothesis}
\newtheorem{question}[theorem]{Question}
\newtheorem{algorithm}[theorem]{Algorithm}
\newtheorem{example}[theorem]{Example}
\theoremstyle{remark}
\newtheorem{remark}[theorem]{Remark}
\providecommand{\lcm}{\operatorname{lcm}}
\providecommand{\supp}{\operatorname{supp}}
\providecommand{\Stab}{\operatorname{Stab}}
\providecommand{\one}{\mathbf 1}
\providecommand{\E}{\mathbb E}
\providecommand{\Var}{\operatorname{Var}}
\providecommand{\ad}{\operatorname{ad}}
\providecommand{\dist}{\operatorname{dist}}
\lstset{basicstyle=\ttfamily\scriptsize,breaklines=true,columns=fullflexible,
  keepspaces=true,showstringspaces=false,frame=single,rulecolor=\color{quiet!30},
  xleftmargin=4pt,xrightmargin=4pt,aboveskip=8pt,belowskip=8pt,
  captionpos=b,tabsize=4}
\hypersetup{pdftitle={Sharp Bounds and Obstructions in Prime Products, Graph Sampling, and Zeta-Zero Correlations},
pdfauthor={Research manuscript prepared for Keio},
pdfsubject={Consolidated proofs; Dirichlet GRH; finite groups; Randić index; conditional prime gaps},
pdfkeywords={Dirichlet GRH, subset products, finite abelian groups, Randic index, random induced subgraphs, prime gaps}}
\begin{document}
\pagenumbering{roman}
\begin{titlepage}
\thispagestyle{empty}
\noindent{\small\sffamily\bfseries\color{teal} NUMBER THEORY AND EXTREMAL GRAPH THEORY}
\vspace{14pt}

\noindent{\color{ink}\LARGE\bfseries Sharp Bounds and Obstructions\par
in Prime Products, Graph Sampling,\par
and Zeta-Zero Correlations\par}
\vspace{10pt}
\noindent{\large A consolidated research study with complete proofs,\par
conditional implications, and unresolved extensions}
\vspace{13pt}

\noindent Research manuscript prepared for Keio\par
\noindent 22 November 2026
\vspace{12pt}\hrule\vspace{12pt}

\noindent\textbf{Abstract.}
This article consolidates the mathematical results developed in a research
programme starting from the Riemann hypothesis and Dirichlet GRH. For
small-prime pools, it proves a sharp asymptotic one-third dichotomy for
squarefree subset products, identifies finite-quotient obstructions at
arbitrary and vanishing retention densities, and gives first-order
optimal compression of the attainable support. Further consequences
include exact-cardinality estimates, short-window variants, a weighted
two-prime exceptional-set estimate, and induced quadratic-residue patterns.
The finite-group arguments are separated from their GRH applications.

An independent, unconditional part determines sharp constants for the sum
of generalized Randić indices over all vertex deletions, including every
equality case and the transition at exponent one half. The result extends
to arbitrary uniform vertex sampling: the extremal problem reduces exactly
to a one-variable hypergeometric maximum, attained by clique components.
The deletion inequality forces any smallest counterexample to the
Caporossi--Hansen additive conjecture to have a cut vertex; the conjecture
itself remains unresolved here.

A third part proves that an explicitly stated Gaussian zero-correlation
hypothesis, together with RH, implies prime gaps of order
$o(\sqrt p)$ and the expected asymptotic count between consecutive squares.
It also records a positive-kernel transfer to weighted prime Dirichlet series
and the off-diagonal estimate still missing from that approach. Exact
identities and numerical diagnostics explain the limitations of an attempted
irrationality proof for the Gieseking volume. Finite certificates and
reproducibility code are included. All assumptions and incomplete steps
are stated explicitly; bibliographic priority is not asserted.
\vfill
\noindent{\small\color{quiet} One self-contained source file. No external figure,
bibliography, or TeX input file is required.}
\end{titlepage}

\section*{Reader's guide and logical status}
\addcontentsline{toc}{section}{Reader's guide and logical status}
The three parts share a research objective---understanding what can be
deduced from balanced arithmetic data, positive kernels, and extremal
inequalities---but their conclusions have different logical status.
The graph theorems do not use RH. The prime-product applications use
Dirichlet GRH. The prime-gap implication uses an additional zero-correlation
hypothesis that is not proved in this article.

\begin{center}\small
\begin{tabularx}{\textwidth}{@{}>{\raggedright\arraybackslash}p{.29\textwidth}>{\raggedright\arraybackslash}X>{\raggedright\arraybackslash}p{.24\textwidth}@{}}
\toprule
Material & Result supplied & Assumptions and status\\
\midrule
Finite-group subset products & Obstruction structure, Fourier bounds,
and constructive compression proofs & Unconditional finite statements\\[4pt]
Small-prime applications & One-third dichotomy; fixed and vanishing
density classification; repair and compression & Dirichlet GRH;
asymptotic thresholds\\[4pt]
Two-prime products & A logarithmic exceptional set and one fixed
correction prime & Dirichlet GRH; full coverage not established\\[4pt]
Graph deletion and sampling & Exact optimal constants and complete
equality classifications & Unconditional\\[4pt]
Caporossi--Hansen & A cut-vertex reduction and exact tests of candidate
induction schemes & Reduction proved; full conjecture open here\\[4pt]
Primes in square-root intervals & Conditional short-interval asymptotics
and $o(\sqrt p)$ gaps & RH plus the stated Gaussian hypothesis\\[4pt]
Gieseking and related diagnostics & Exact algebraic identities, a norm
barrier, and exploratory determinant data & No irrationality conclusion\\
\bottomrule
\end{tabularx}
\end{center}

The mathematical premise is the stated meaning of the five zeta and
Dirichlet declarations supplied in the discussion. A record of those
declarations is retained in Part~\ref{r-part-prime}. They are not treated
as five independent analytic inputs. The arguments presented here are
ordinary mathematical proofs; this article does not claim a Lean
formalization or an external referee verification.

The full proofs of the finite and conditional results are retained.
Results already established in cited literature, including the
Mills-type transcendence consequence, are attributed as such. Numerical
experiments and failed transfer arguments are explicitly separated from
theorems. A search that did not locate an identical formulation is not a
proof of historical novelty.

\clearpage
\begingroup
\small
\setlength{\parskip}{0pt}
\linespread{1}\selectfont
\tableofcontents
\endgroup
\clearpage
\pagenumbering{arabic}

\clearpage
\part{Small-prime products under Dirichlet GRH}\label{r-part-prime}


% Consolidated prime chapter. All local identifiers use the p- prefix.
\providecommand{\pSP}{\operatorname{SP}}
\providecommand{\pGq}{(\mathbb Z/q\mathbb Z)^\times}
\providecommand{\pli}{\operatorname{li}}
\providecommand{\psupp}{\operatorname{supp}}
\providecommand{\pone}{\mathbf 1}
\providecommand{\pGRHD}{\mathrm{GRH}_{\mathrm D}}
\providecommand{\pStab}{\operatorname{Stab}}

\section{Scope, conventions, and main theorem}

Throughout, $\log$ is the natural logarithm and $\log_2 x=\log x/\log2$.
We write $\pGRHD$ for GRH for every primitive Dirichlet $L$-function,
including the Riemann hypothesis through the principal primitive
character. Standard passage to imprimitive characters is understood.
This is the standing premise for the prime results below. The supplied
zeta and Dirichlet critical-line statements are used through these
mathematical meanings; they do not supply five independent analytic
hypotheses. The present proofs are ordinary mathematical proofs, not
a Lean formalization.

For a finite set $R$ of primes not dividing $q$, define
\[
 \pSP_q(R)=\left\{\prod_{p\in U}p\pmod q:U\subseteq R\right\}.
\]
The empty product is included. Each prime can appear at most once in
each product. This differs from ordinary generation of a group, which
allows repeated powers and inverses. Different subsets may represent
the same residue; all resulting integers are squarefree. In the
finite-group sections we write $\Sigma(R)$ for the same construction,
and in prime applications $\Sigma_q(R)=\pSP_q(R)$.

\begin{theorem}[One-third principle and optimal compression]\label{p-thm:main}
Assume $\pGRHD$. Fix
\[
 B>2,\qquad 0<c<\min\{B/2-1,1\}.
\]
For every sufficiently large integer $q$, put
\[
 \ell=\log q,\quad
 P_q=\{p\le\ell^B:p\text{ prime},\ p\nmid q\},\quad d=|P_q|.
\]
Let $R\subseteq P_q$ satisfy
\begin{equation}\label{p-eq:main-density}
 |R|\ge(1/3+\ell^{-c})d.
\end{equation}
Then, writing $H=\langle R\rangle\le\pGq$, the following hold uniformly
over the choice of $R$.
\begin{enumerate}[label=(\roman*),leftmargin=*,itemsep=4pt]
\item $\pSP_q(R)=H$ and $[\pGq:H]\in\{1,2\}$.
\item If $H\ne\pGq$, then $H=\ker\chi$ for a unique nonprincipal
quadratic character $\chi\bmod q$, and $\chi(p)=1$ for every $p\in R$.
\item There is a set $T\subseteq R$ of distinct primes such that
\begin{equation}\label{p-eq:main-compression}
 \pSP_q(T)=H,\qquad
 |T|\le\log_2|H|+
 O_{B,c}\!\left(\frac{\log q}{\log\log q}\right).
\end{equation}
\item If $H\ne\pGq$, any one prime $p\in P_q\setminus H$ restores full
coverage: $\pSP_q(T\cup\{p\})=\pGq$. Such primes exist.
\end{enumerate}
The implied constants and the sufficiently-large threshold depend only
on $B,c$.
\end{theorem}

For a prime modulus $q$, there is just one nonprincipal quadratic
character, the Legendre symbol. Thus, in the range of the theorem,
\emph{one quadratic nonresidue among the retained primes is enough to
rule out every obstruction to squarefree coverage}. For composite $q$,
the equivalent test is that the images of $R$ span $\pGq/(\pGq)^2$.

The theorem permits adversarial deletion: the removed primes can be
chosen with complete knowledge of their residues. The small covering
set $T$ is chosen \emph{after} the deletion and may depend on both $q$
and $R$. It is not a claim that deleting nearly two thirds of an already
compressed $T$ preserves coverage.

The bound is optimal at leading order because any $t$ primes have only
$2^t$ subsets, so covering $H$ requires
$t\ge\lceil\log_2|H|\rceil$. The threshold $1/3$ is also sharp in an
asymptotic sense, proved in Section~\ref{p-sec:sharpness}. No exact finite
cutoff at $d/3$ is asserted. In particular, the positive margin in
\eqref{p-eq:main-density} cannot simply be erased.

\subsection{A single squarefree integer and simultaneous congruences}
If $M=\prod_{p\in T}p$, then the divisors of $M$ represent exactly $H$,
and
\[
 \tau(M)=2^{|T|}=|H|^{1+o(1)},\qquad
 \log M\le\left(\frac B{\log2}+o(1)\right)
       (\log q)\log\log q.
\]
Here $\tau(M)$ is the number of positive divisors. The result optimizes
the number of shared prime factors, not the numerical size of each
representative. If $q=\prod_j q_j$ with pairwise coprime moduli and
$H=\pGq$, the same divisor set realizes every tuple of unit residues
modulo the $q_j$ by the Chinese remainder theorem. This is a direct
corollary, not a separate novelty claim.




\subsection{Context and the scope of the claims}

Jao--Miller--Venkatesan~\cite{p-JMV} derive expansion from primes of
size $(\log q)^B$, $B>2$, under GRH. Booker--Pomerance~\cite{p-BP}
and Munsch--Shparlinski~\cite{p-MS} establish squarefree smooth
representations in other size regimes. Optimizing an individual
representative differs from choosing one shared list of prime factors
whose divisors cover many residue classes.

Klurman--Shparlinski--Ter\"av\"ainen~\cite{p-KST} construct a common
set of polylogarithmically many small primes for almost all moduli,
and record a GRH Boolean-covering consequence. The primary abstract
of Hughes~\cite{p-Hughes} states a first-order optimal
$(1/\log2+o(1))\log Q$ count in an almost-all-moduli setting.
Thus a logarithmic number of shared primes alone is not claimed as
a new discovery. Nor is a conditional modulus-dependent construction
being compared as a blanket improvement over an unconditional
common-list theorem. The structure of incomplete sets in abelian
groups is an established subject; see Vu~\cite{p-Vu}. Here relative
density is measured inside a sparse character-balanced pool.

The results proved in this part comprise the sharp one-third
quadratic dichotomy, the obstruction classification at every fixed
positive density, its quantitative vanishing-density extension,
exact-support compression, correction by additional primes, and
several refinements. All finite-group arguments are included.
No assertion of publication priority is made for their combined
formulation. The subsequent two-prime results leave the
Erd\H{o}s--Odlyzko--S\'ark\"ozy problem open.



\section{The analytic input}
Under $\pGRHD$, the standard prime-character estimate is
\begin{equation}\label{p-eq:prime-character}
 \sum_{p\le y}\psi(p)
 =\pone_{\psi=\psi_0}\pli(y)
  +O\bigl(\sqrt y\log(qy)\bigr),\qquad y\ge2,
\end{equation}
uniformly for Dirichlet characters modulo $q$. Characters are zero on
nonunits. The complex-valued weighted form, the correction for
imprimitive characters, and partial summation are given in
\cite[Section 2, equation (2.6)]{p-JMV}; see also
\cite[Theorem 13.7]{p-MV}. Only the rational-field Dirichlet case is used.

Let $y=\ell^B$, $\ell=\log q$, and $d=|P_q|$. The prime number theorem
and the elementary bound $\omega(q)=O(\ell)$ give
\begin{equation}\label{p-eq:d-asymptotic}
 d\sim\frac{\ell^B}{B\log\ell}.
\end{equation}
For every nonprincipal character,
\begin{equation}\label{p-eq:prime-bias}
 \frac1d\left|\sum_{p\in P_q}\psi(p)\right|
 \ll_B(\log\ell)\ell^{-\kappa},\qquad \kappa=B/2-1>0.
\end{equation}
This small character average is the only analytic input for the main
theorem. In particular, all remaining group-theoretic lemmas are
unconditional.



\section{Finite Fourier bounds and quadratic obstructions}\label{p-sec:inverse}
Let $G$ be a finite abelian group of order $N\ge2$. A pool $P$ consists
of $d$ distinct labels carrying elements of $G$; repeated group elements
are allowed but a label can be used only once. Call the pool
$\varepsilon$-balanced if
\begin{equation}\label{p-eq:inverse-bias}
 \left|\sum_{p\in P}\psi(p)\right|\le\varepsilon d
 \quad(1\ne\psi\in\widehat G),\qquad 0\le\varepsilon<1.
\end{equation}
For $R\subseteq P$, let $r=|R|$, let $\Sigma(R)$ denote its subset-product
set, and let $r_R(a)$ count its subsets with product $a$. Define
\[
 K=3r-(1+2\varepsilon)d.
\]

\begin{theorem}[Complete obstruction classification above one third]\label{p-thm:finite-inverse}
If
\begin{equation}\label{p-eq:inverse-threshold}
 K>16\log(N-1),
\end{equation}
then
\[
 \Sigma(R)=\langle R\rangle,
 \qquad [G:\langle R\rangle]\in\{1,2\}.
\]
More precisely, either $\Sigma(R)=G$, or there is a unique nonprincipal
quadratic character $\chi$ such that
\[
 R\subseteq\ker\chi,
 \qquad \Sigma(R)=\ker\chi.
\]
In particular, the following are equivalent:
\begin{enumerate}
\item $\Sigma(R)=G$;
\item $R$ generates $G$ as a group;
\item for every nonprincipal quadratic character $\chi$, at least one
label $p\in R$ satisfies $\chi(p)=-1$.
\end{enumerate}
If $H=\langle R\rangle$ and $h=[G:H]\in\{1,2\}$, then uniformly for
$a\in H$,
\begin{equation}\label{p-eq:inverse-equidistribution}
 \left|\frac{|H|\,r_R(a)}{2^r}-1\right|
 \leq\frac{N-h}{h}\,e^{-K/16}.
\end{equation}
For $a\notin H$, $r_R(a)=0$.
\end{theorem}

\begin{proof}
For a character $\psi$, put
\[
 F_R(\psi)=\prod_{p\in R}\frac{1+\psi(p)}2,
 \qquad
 D_R(\psi)=\sum_{p\in R}(1-\operatorname{Re}\psi(p)).
\]
Since $|z|=1$ implies
\[
 \left|\frac{1+z}{2}\right|^2
 =1-\frac{1-\operatorname{Re}z}{2}
 \leq e^{-(1-\operatorname{Re}z)/2},
\]
we have
\begin{equation}\label{p-eq:inverse-fourier-product}
 |F_R(\psi)|\leq e^{-D_R(\psi)/4}.
\end{equation}

Suppose first that $\psi^2\ne1$. The degree-two Fej\'er polynomial
\[
 F(z)=|1+z+z^2|^2
     =3+4\operatorname{Re}z+2\operatorname{Re}z^2
\]
is nonnegative for $|z|=1$, and
\begin{equation}\label{p-eq:fejer-lower}
 F(z)\geq9-12(1-\operatorname{Re}z).
\end{equation}
Indeed, writing $u=\operatorname{Re}z$, the difference between the two
sides is $4(u-1)^2$. Nonnegativity and \eqref{p-eq:inverse-bias} give
\[
 9r-12D_R(\psi)
 \leq\sum_{p\in R}F(\psi(p))
 \leq\sum_{p\in P}F(\psi(p))
 \leq(3+6\varepsilon)d.
\]
Both $\psi$ and $\psi^2$ are nonprincipal here. It follows that
\[
 D_R(\psi)\geq K/4,
 \qquad |F_R(\psi)|\leq e^{-K/16}.
\]

For a nonprincipal quadratic character $\chi$, the coefficient
$F_R(\chi)$ is exactly $1$ if $R\subseteq\ker\chi$, and exactly $0$
otherwise. There can be at most one such quadratic kernel containing
$R$. Otherwise distinct quadratic characters $\chi_1,\chi_2$ would give
\[
 r\leq |P\cap\ker\chi_1\cap\ker\chi_2|
 =\frac14\sum_{p\in P}(1+\chi_1(p))(1+\chi_2(p))
 \leq\frac{1+3\varepsilon}{4}d.
\]
This contradicts $K>0$, since
$(1+2\varepsilon)/3-(1+3\varepsilon)/4=(1-\varepsilon)/12>0$.

If no quadratic kernel contains $R$, put $H=G$ and $h=1$. If one does,
put $H=\ker\chi$ and $h=2$. Character orthogonality now yields
\[
 \frac{N r_R(a)}{2^r}
 =h\,1_H(a)+E(a),
 \qquad |E(a)|\leq (N-h)e^{-K/16}.
\]
The error comes only from characters of order at least three; every
other nonprincipal quadratic coefficient vanishes. Under
\eqref{p-eq:inverse-threshold} the error has absolute value less than $1$.
Thus every $a\in H$ has a representation. Conversely, all elements of
$R$ belong to $H$, so every subset product belongs to $H$.
Consequently $\Sigma(R)=H=\langle R\rangle$. Division by $h$ proves
\eqref{p-eq:inverse-equidistribution}. The stated equivalences follow.
\end{proof}

\begin{corollary}[One-prime repair]
Under the preceding hypotheses, if $\Sigma(R)\ne G$, any label
$p\in P$ outside its unique quadratic kernel satisfies
\[
 \Sigma(R\cup\{p\})=G.
\]
Such labels exist, since their number is at least $(1-\varepsilon)d/2$.
\end{corollary}
\begin{proof}
Write $\Sigma(R)=H$ with $[G:H]=2$. For $p\notin H$,
$\Sigma(R\cup\{p\})=H\cup pH=G$.
\end{proof}




\section{Near-minimal compression in a finite group}\label{p-sec:compression}
The next argument extracts a small common collection, rather than
retaining the entire large set of available labels.

\begin{lemma}[Positive overlap estimate]\label{p-lem:overlap}
For every $A\subseteq G$, $a=|A|$, a balanced pool satisfies
\[
 \sum_{p\in P}|A\cap pA|
 \le d\left(\frac{a^2}{N}+\varepsilon a(1-a/N)\right).
\]
Consequently, for any $R'\subseteq P$ with $|R'|\ge\beta d$, some
$p\in R'$ has
\begin{equation}\label{p-eq:overlap-positive}
 |A\cap pA|\le\frac{a^2/N+\varepsilon a}{\beta}.
\end{equation}
\end{lemma}
\begin{proof}
With $\widehat{\pone_A}(\psi)=\sum_{a\in A}\overline{\psi(a)}$,
Fourier inversion writes the sum as
\[
 \frac1N\sum_{\psi\in\widehat G}|\widehat{\pone_A}(\psi)|^2
                   \sum_{p\in P}\overline{\psi(p)}.
\]
The principal contribution is $da^2/N$. Parseval gives nonprincipal
mass $Na-a^2$, yielding the first assertion. All individual overlaps
are nonnegative, so their sum over $R'$ is no greater than their sum
over $P$. Averaging gives the second assertion. No symmetry of $P$
under inversion is needed.
\end{proof}
For a probability measure $\nu$ on $G$, use the normalised-character
transform $\widehat\nu(\chi)=\sum_{g\in G}\nu(g)\chi(g)$ and the energy
\[
 E(\nu)=\sum_{\chi\ne1}|\widehat\nu(\chi)|^2
       =N\sum_{g\in G}\nu(g)^2-1.
\]
For a label $p$, put $b_p=(\delta_1+\delta_p)/2$.

\begin{lemma}[Weighted positive restriction]
If $R'\subseteq P$ contains at least $\beta d$ labels, some $p\in R'$
satisfies
\begin{equation}\label{p-eq:weighted-positive}
 E(\nu*b_p)\leq
 \frac12\left(1+\frac\varepsilon\beta\right)E(\nu)
 +\frac12\left(\frac1\beta-1\right).
\end{equation}
\end{lemma}
\begin{proof}
For $C_p=\langle\nu,\nu(p^{-1}\cdot)\rangle$ one has $C_p\geq0$.
Fourier inversion and the bias hypothesis give
\[
 \sum_{p\in P}C_p\leq\frac dN(1+\varepsilon E(\nu)).
\]
Restricting this sum to $R'$ and averaging gives an available label with
$NC_p\leq(1+\varepsilon E(\nu))/\beta$.
Since
\[
 \|\nu*b_p\|_2^2=\tfrac12\|\nu\|_2^2+\tfrac12 C_p,
\]
the asserted estimate follows.
\end{proof}

\begin{theorem}[Finite compression above one third]\label{p-thm:finite-compression}
Suppose $P$ is $\varepsilon$-biased, and $R\subseteq P$ generates $G$.
Put $L=\log N$ and $H=\log(1/\varepsilon)$, and assume
\[
 16\leq H\leq L,\qquad
 0<\eta\leq\tfrac23,\qquad
 \varepsilon\leq\eta/4,\qquad
 |R|\geq(1/3+\eta)d.
\]
Define
\[
 M=\left\lceil
 \frac L{\log2}+\frac{16L}{H}
 +\frac{8\log(16H)}{\eta}+H+5
 \right\rceil.
\]
If $\eta d\geq3M$, then $R$ contains distinct labels $T$ with
$\Sigma(T)=G$ and $|T|\leq M$.
\end{theorem}
\begin{proof}
Choose $U\subseteq R$ whose images are a basis of the elementary
abelian group $G/G^2$, and put $s=|U|$.
The $2^s$ subsets of $U$ give distinct images in $G/G^2$, and hence
distinct elements of $G$. The probability measure
$\nu_U=\mathop{*}_{p\in U}b_p$ is uniform on these $2^s$ elements,
so
\[
 E(\nu_U)=N/2^s-1.
\]
Every nonprincipal quadratic Fourier coefficient of $\nu_U$ is zero:
its character is negative on at least one of the selected basis labels.
These zero coefficients remain zero after all subsequent convolutions.

All labels below are chosen outside those already used.
As long as at most $M\leq\eta d/3$ labels have been used, the remaining
pool $R'$ has at least $(1/3+2\eta/3)d$ labels, in particular at least
$d/3$. Moreover
\begin{equation}\label{p-eq:unused-third-gap}
 3|R'|-(1+2\varepsilon)d
 \geq(2\eta-2\varepsilon)d\geq\eta d.
\end{equation}

First, while $E(\nu)>H$, use \eqref{p-eq:weighted-positive} with
$\beta=1/3$ to choose a new label satisfying
\[
 E(\nu*b_p)
 \leq\tfrac12(1+3\varepsilon)E(\nu)+1
 \leq\left(\tfrac12+\frac2H\right)E(\nu).
\]
Here $\varepsilon=e^{-H}$ and $H\geq16$.
If this phase is nonempty, it uses at most
\[
 \frac{\log(N/2^s)}{\log(1/(1/2+2/H))}+1
 \leq\frac L{\log2}-s+\frac{14L}{H}+1
\]
labels. If it is empty, $|U|=s\leq\log_2N$ already.
Thus in either case the total number of labels used so far is at most
$L/\log2+14L/H+1$.

Next, for every character $\chi$ of order at least three, the Fej\'er
estimate and \eqref{p-eq:unused-third-gap} give
\[
 D_{R'}(\chi)
 :=\sum_{p\in R'}(1-\operatorname{Re}\chi(p))\geq\eta d/4.
\]
Consequently
\[
 \frac1{|R'|}\sum_{p\in R'}|\widehat b_p(\chi)|^2
 =1-\frac{D_{R'}(\chi)}{2|R'|}\leq1-\eta/8.
\]
Quadratic coefficients are already zero, so averaging the energy over
available labels gives one for which
$E(\nu*b_p)\leq(1-\eta/8)E(\nu)$.
At most $8\log(16H)/\eta+1$ additional labels make $E(\nu)\leq1/16$.

At every stage $\nu$ is a probability measure supported on
$\Sigma(T)$ for the distinct labels used so far. Cauchy--Schwarz gives
\[
 |\operatorname{supp}\nu|\geq\frac N{1+E(\nu)}\geq\frac{16N}{17}.
\]
Therefore the actual subset-product set $S=\Sigma(T)$ misses at most
$N/17$ elements. For its missing set $B=G\setminus S$, the unweighted
overlap estimate, restricted to at least $d/3$ unused labels, supplies
a label with
\[
 |B\cap pB|\leq3\left(\frac{|B|^2}{N}+\varepsilon|B|\right).
\]
While $\varepsilon N<|B|\leq N/17$, this contracts $|B|$ by a factor
at most $3/17+3\varepsilon<1/4$. At most $H+1$ further labels give
$|B|\leq\varepsilon N$. Subsequently the contraction factor is at most
$6\varepsilon$, and at most $2L/H+2$ further labels make $|B|<1$.
Since $|B|$ is an integer, $S=G$.

The total count is at most $M$. The assumed inequality
$\eta d\geq3M$ justifies the unused-pool bounds throughout, so the
argument is non-circular. No prime, or other label, has been used twice.
\end{proof}




\section{Proof and sharpness of the prime theorem}
\subsection{Uniform specialization of the finite bounds}
Write $\ell=\log q$, $\kappa=B/2-1$, and choose
\begin{equation}\label{p-eq:parameters}
 \eta=\ell^{-c},\qquad
 \varepsilon=\ell^{-(c+\kappa)/2},\qquad
 H_0=\log(1/\varepsilon)=\frac{c+\kappa}{2}\log\ell.
\end{equation}
Since $c<\kappa$, equation~\eqref{p-eq:prime-bias} is bounded by
$\varepsilon$ for all sufficiently large $q$, and
$\varepsilon=o(\eta)$. Furthermore,
\[
 \frac{\eta d}{\ell}\longrightarrow\infty,\qquad
 \eta^{-1}\log H_0=O_{B,c}(\ell^c\log\log\ell)
       =o(\ell/\log\ell),
\]
where the last assertion uses $c<1$. Also
$\log\varphi(q)\asymp\ell$, for example from
$\varphi(q)\ge\sqrt{q/2}$.

For $r=|R|\ge(1/3+\eta)d$,
\[
 K=3r-(1+2\varepsilon)d\ge(3\eta-2\varepsilon)d
       \gg\eta d\gg\log\varphi(q).
\]
Theorem~\ref{p-thm:finite-inverse} therefore proves the asserted support
dichotomy and identifies $H=\langle R\rangle$.

If $H=\pGq$, apply Theorem~\ref{p-thm:finite-compression}. Its bound is
\[
 M=\log_2\varphi(q)+O_{B,c}(\ell/\log\ell),
\]
and its reserve condition $\eta d\ge3M$ follows from the preceding
estimates. This proves the desired compression in the full group.

If $H=\ker\chi$ has index two, put $P_H=P_q\cap H$, $d_H=|P_H|$.
Then
\[
 d_H=\tfrac12(d+O(\varepsilon d)).
\]
Every nonprincipal character $\rho$ of $H$ extends to a character
$\psi$ of $\pGq$. Both $\psi$ and $\psi\chi$ are nonprincipal, so
\[
 \left|\sum_{p\in P_H}\rho(p)\right|
 =\left|\frac12\sum_{p\in P_q}(\psi(p)+\psi(p)\chi(p))\right|
 \le\varepsilon d.
\]
Thus $P_H$ is $4\varepsilon$-balanced for large $q$. Also $R$ generates
$H$ and $|R|/d_H\ge(2/3+2\eta)/(1+\varepsilon)>1/2$ eventually.
Apply the finite compression theorem inside $H$ with a fixed margin
$\eta_H=1/6$ and $\varepsilon_H=4\varepsilon$. All its hypotheses
hold for large $q$. This gives
\[
 |T|\le\log_2|H|+O_{B,c}(\ell/\log\ell).
\]
Finally, any $p\in P_q\setminus H$ gives
$\Sigma(T\cup\{p\})=H\cup pH=\pGq$. There are at least
$(1-\varepsilon)d/2$ such available primes. This proves every assertion
of Theorem~\ref{p-thm:main}.

\subsection{Uniform multiplicities on the attainable subgroup}
Let $h=[\pGq:H]\in\{1,2\}$ and $r=|R|$. The finite theorem yields
\begin{equation}\label{p-eq:prime-counts}
 \left|\frac{|H|r_R(a)}{2^r}-1\right|
 \le\frac{\varphi(q)-h}{h}
 \exp\left(-\frac{(3\eta-2\varepsilon)d}{16}\right)
 \quad(a\in H).
\end{equation}
For every fixed $A>0$, the right-hand side is $o(q^{-A})$, uniformly
over all permitted $R,a$. Thus the conclusion includes extremely
uniform counts over the full surviving pool. The compressed $T$ is
only guaranteed to cover $H$; this estimate is not transferred to $T$.

\subsection{Sharpness from cubic characters}\label{p-sec:sharpness}
There are arbitrarily large prime moduli $q\equiv1\pmod3$. Fix such a
$q$ and a character $\chi$ of order three. Put $J=\ker\chi$ and
$P_0=P_q\cap J$. Orthogonality gives
\[
 |P_0|=(1/3+O(\varepsilon))d.
\]
Every nonprincipal character $\rho$ of $J$ extends to a character
$\psi$ of $\pGq$, and each $\psi\chi^j$ ($j=0,1,2$) is
nonprincipal. Hence
\[
 \left|\sum_{p\in P_0}\rho(p)\right|
 =\left|\frac13\sum_{j=0}^2\sum_{p\in P_q}\psi(p)\chi(p)^j\right|
 \le\varepsilon d.
\]
Since $|P_0|\ge(1-2\varepsilon)d/3$, every nonprincipal $\rho$ has
\[
 D_{P_0}(\rho)\ge(1-5\varepsilon)d/3\gg\log|J|.
\]
The Fourier product bound from Section~\ref{p-sec:inverse} implies
$\pSP_q(P_0)=J$ for all sufficiently large $q$.

Choose $p_1\in P_q$ with $\chi(p_1)=\omega=e^{2\pi i/3}$;
that coset contains $(1/3+O(\varepsilon))d$ available primes. Then
\[
 R=P_0\cup\{p_1\},\quad |R|=(1/3+o(1))d,\quad
 \langle R\rangle=\pGq,
\]
but
\[
 \pSP_q(R)=J\cup p_1J\ne\pGq.
\]
The third coset $p_1^2J$ is absent. This proves asymptotic sharpness
of the one-third threshold for generation implying squarefree coverage.
The simpler retained pool $P_0$ also shows sharpness for one-prime
repair: adding any one prime can cover at most two of the three cosets.

\subsection{The half-density corollary without an obstruction test}
Under the parameters of Theorem~\ref{p-thm:main}, if
$|R|\ge(1/2+\ell^{-c})d$, no quadratic kernel can contain $R$, because
each such kernel contains at most $(1+\varepsilon)d/2$ candidate primes.
Therefore $H=\pGq$ automatically, with no prime needing restoration.
Here ``automatically'' concerns the absence of a further algebraic
condition; the statement still assumes $\pGRHD$.
The half-density threshold for full coverage without any obstruction
test is asymptotically optimal by retaining only the primes in a
quadratic kernel.



\section{All positive densities: bounded quotients and exact compression}

We work with a family of finite abelian groups $G$, of orders $N\to\infty$,
and labelled pools $P$, of sizes $d$. Different labels may carry the same
group element. Assume
\begin{equation}\label{p-eq:ad-balanced}
 \max_{\chi\ne1}\left|\sum_{p\in P}\chi(p)\right|\leq\varepsilon d,
 \qquad \varepsilon\longrightarrow0,
 \qquad \frac d{\log N}\longrightarrow\infty.
\end{equation}
For labels $R\subseteq P$, write $\Sigma(R)$ for the set of products of
all subsets of $R$, including the empty product. All assertions below
are uniform in $R$ of a fixed positive relative density.

\begin{lemma}[Restriction to a subgroup]\label{p-lem:subgroup-restriction}
For any subgroup $J\leq G$ of index $j$, put $P_J=P\cap J$ and
$d_J=|P_J|$. Then
\[
 d_J=\frac dj+O(\varepsilon d),
 \qquad
 \left|\sum_{p\in P_J}\rho(p)\right|\leq\varepsilon d
 \quad(1\ne\rho\in\widehat J).
\]
The constants in these two statements do not depend on $j$.
\end{lemma}
\begin{proof}
The annihilator $J^\perp$ has $j$ characters, and
$1_J(g)=j^{-1}\sum_{\psi\in J^\perp}\psi(g)$. This gives the first
assertion. Extend a nonprincipal $\rho\in\widehat J$ to a character
$\rho_0$ of $G$. All characters in $\rho_0J^\perp$ are nonprincipal,
and averaging their prime sums proves the second assertion.
\end{proof}

\begin{lemma}[A higher Fej\'er estimate]\label{p-lem:higher-fejer}
Let $J\leq G$, let $R_J\subseteq P_J$ have $r_J$ labels, and let $\chi$
be a character of $J$ of order at least $m$, where $m\geq2$. Then
\begin{equation}\label{p-eq:ad-fejer}
 r_J\leq\frac{d_J}{m}
       +\left(1-\frac1m\right)\varepsilon d
       +\frac{m^2-1}{6}D_{R_J}(\chi),
 \qquad
 D_{R_J}(\chi)=\sum_{p\in R_J}(1-\operatorname{Re}\chi(p)).
\end{equation}
\end{lemma}
\begin{proof}
The nonnegative polynomial
\[
 F_m(z)=|1+z+\cdots+z^{m-1}|^2
       =m+2\sum_{a=1}^{m-1}(m-a)\operatorname{Re}z^a
\]
satisfies, on $|z|=1$,
\[
 F_m(z)\geq m^2-\frac{m^2(m^2-1)}6(1-\operatorname{Re}z).
\]
Indeed, $1-\operatorname{Re}z^a\leq a^2(1-\operatorname{Re}z)$ and
$2\sum_{a=1}^{m-1}(m-a)a^2=m^2(m^2-1)/6$.
The characters $\chi^a$, $1\leq a<m$, are all nonprincipal. Thus
\[
 m^2r_J-\frac{m^2(m^2-1)}6 D_{R_J}(\chi)
 \leq\sum_{p\in R_J}F_m(\chi(p))
 \leq\sum_{p\in P_J}F_m(\chi(p))
 \leq m d_J+m(m-1)\varepsilon d.
\]
Division by $m^2$ gives the claim.
\end{proof}

\begin{theorem}[All-density obstruction theorem]\label{p-thm:all-density}
Fix $0<\alpha\leq1$. Under \eqref{p-eq:ad-balanced}, for all sufficiently
large members of the family the following holds for every
$R\subseteq P$ with $|R|\geq\alpha d$.
Put $S=\Sigma(R)$. If $S\ne G$, let
\[
 H=\operatorname{Stab}(S)=\{g\in G:gS=S\},\qquad h=[G:H].
\]
Then
\begin{equation}\label{p-eq:ad-classification}
 2\leq h\leq\lfloor1/\alpha\rfloor,
 \qquad |R\setminus H|\leq h-2,
 \qquad
 S=\pi^{-1}\bigl(\Sigma(\pi(R\setminus H))\bigr),
\end{equation}
where $\pi:G\to G/H$ is the quotient map. In particular, every failure
of subset-product coverage at a fixed positive density is controlled
by a quotient of bounded order and at most $h-2$ exceptional labels.
\end{theorem}
\begin{proof}
We first find a subgroup $J$ such that $\Sigma(R\cap J)=J$, with
$[G:J]\leq1/\alpha+o(1)$.

Begin with $J=G$. If $\Sigma(R\cap J)\ne J$, character orthogonality
at a missing element yields a nonprincipal character $\chi$ of $J$
with
\[
 \left|\prod_{p\in R\cap J}\frac{1+\chi(p)}2\right|
 \geq\frac1{|J|-1}.
\]
Since this product is at most $e^{-D_{R\cap J}(\chi)/4}$, we have
$D_{R\cap J}(\chi)\leq4\log(|J|-1)\leq4\log N$.

Choose a fixed integer $m>8/\alpha$. As long as only a bounded number
of the following steps have been performed, the current set $R\cap J$
has $\alpha d-O_\alpha(\log N)$ labels. If $\chi$ had order at least
$m$, \eqref{p-eq:ad-fejer} would give
\[
 |R\cap J|\leq d/m+\varepsilon d+O_\alpha(\log N)
 <\alpha d/2
\]
for sufficiently large members of the family, a contradiction.
Consequently $\chi$ has order $u$ with $2\leq u<m$.

Replace $J$ by $J\cap\ker\chi$, where the kernel is understood inside
the current $J$. This multiplies $[G:J]$ by $u\geq2$ and discards at
most $O_\alpha(\log N)$ labels of $R$: every discarded label contributes
at least
$\min_{2\leq u<m}(1-\cos(2\pi/u))>0$ to $D_{R\cap J}(\chi)$.

The process must stop after a bounded number of steps. Otherwise, after
a fixed $k$ with $2^k>4/\alpha$, the current subgroup would have
index at least $2^k$, and hence contain at most
$d/2^k+\varepsilon d<\alpha d/2$ pool labels, whereas it still contains
$\alpha d-O_\alpha(\log N)>\alpha d/2$ labels of $R$.
This contradiction proves termination. At its end,
\[
 \Sigma(R\cap J)=J,
 \qquad |R\cap J|\geq\alpha d-O_\alpha(\log N).
\]
The restriction lemma now gives
$\alpha\leq1/[G:J]+o(1)$. Since the index is a bounded integer, it is
at most $\lfloor1/\alpha\rfloor$ eventually.

The identity
$S=J\,\Sigma(R\setminus J)$ shows that $J\leq H$.
Thus $h=[G:H]\leq[G:J]\leq\lfloor1/\alpha\rfloor$.
The image $\bar S\subseteq G/H$ is proper and aperiodic. Scan the
nonidentity images of the labels in $R$ in any order. Every one must
strictly enlarge the current subset-product set: if a nonidentity
$g$ did not, it would fix that set, and by commutativity would fix all
later unions of its translates. It would then be a nontrivial period
of $\bar S$, which is impossible. Starting from a set of size one, this
gives
\[
 1+|R\setminus H|\leq|\bar S|\leq h-1.
\]
Therefore $|R\setminus H|\leq h-2$. Finally, $H$-periodicity and the
definition of the quotient imply the last identity in
\eqref{p-eq:ad-classification}.
\end{proof}


\subsection{The positive-overlap estimate used in compression}

We record the averaging calculation explicitly. Let $J\le G$, let
$n=|J|$, and let $R'\subseteq P\cap J$ contain at least $\beta d$
unused labels. For a probability measure $\nu$ on $J$, set
$E=n\sum_x\nu(x)^2-1$. The correlations
$C_p=\sum_x\nu(x)\nu(p^{-1}x)$ are nonnegative. Fourier inversion and
the subgroup restriction lemma give
\[
 \sum_{p\in P\cap J}C_p
 \le\frac{|P\cap J|+\varepsilon dE}{n}.
\]
Restriction to $R'$ is legitimate by nonnegativity. Since
\[
 E(\nu*(\delta_1+\delta_p)/2)=(E-1)/2+nC_p/2,
\]
some unused label satisfies
\begin{equation}\label{p-eq:positive-energy}
 E(\nu*b_p)\le
 \frac12(1+\varepsilon/\beta)E
 +\frac12(1/\beta-1).
\end{equation}
Likewise, for a missing set $B\subseteq J$, the nonnegative correlations
$|B\cap pB|$ have full-pool average bounded by
$|P\cap J||B|^2/n+\varepsilon d|B|$. Hence some unused label gives
\begin{equation}\label{p-eq:positive-holes}
 |B\cap pB|\le\beta^{-1}
       \bigl(|B|^2/n+\varepsilon|B|\bigr).
\end{equation}
Inserting that label into the subset-product list changes the missing
set to precisely $B\cap pB$. All energies in this calculation are
normalized by $|J|$, not by the order of the original ambient group.

\subsection{A core with a uniform character gap}

For compression we use a different stopping rule. Fix $0<\alpha\leq1$.
There is a constant $\theta=\theta(\alpha)>0$ such that, under
\eqref{p-eq:ad-balanced}, every $R$ with $|R|\geq\alpha d$ contains a
core $R_J=R\cap J$ in a subgroup $J$ satisfying
\begin{equation}\label{p-eq:ad-core}
 [G:J]\leq2/\alpha,
 \qquad |R_J|\geq3\alpha d/4,
 \qquad
 D_{R_J}(\chi)>\theta d
 \quad(1\ne\chi\in\widehat J).
\end{equation}

Here is a fully quantitative choice of constants for the argument.
Choose $m>8/\alpha$,
$k=\lceil\log_2(4/\alpha)\rceil+1$, and put
\[
 c_* =\min_{2\leq u<m}(1-\cos(2\pi/u)),\qquad
 C_*=(m^2-1)/6,
 \qquad
 \theta=\min\left\{\frac{\alpha}{16C_*},
                         \frac{\alpha c_*}{4k}\right\}.
\]
For sufficiently large members of the family, $\varepsilon\leq\alpha/16$.
Begin with $J=G$ and restrict to the kernel of any nonprincipal
$\chi\in\widehat J$ with $D_{R\cap J}(\chi)\leq\theta d$.
As before, \eqref{p-eq:ad-fejer} forces the order of $\chi$ to be less
than $m$, provided $|R\cap J|\geq3\alpha d/4$.
Each step loses at most $\theta d/c_*$ labels, so $k$ steps lose at
most $\alpha d/4$. But $k$ steps would make the index too large for a
subgroup containing $3\alpha d/4$ pool labels. Therefore the procedure
stops before then. At termination the character gap is as stated,
and the restriction lemma bounds the index by
$1/(3\alpha/4-\varepsilon)<2/\alpha$.
In particular $\Sigma(R_J)=J$, since
$(|J|-1)e^{-\theta d/4}<1$ eventually.

\begin{theorem}[Exact compression at every positive density]\label{p-thm:all-compression}
Fix $0<\alpha\leq1$. Under \eqref{p-eq:ad-balanced}, uniformly for
$R\subseteq P$ with $|R|\geq\alpha d$, there is $T\subseteq R$ with
\[
 \Sigma(T)=\Sigma(R),\qquad
 |T|\leq\log_2|\Sigma(R)|+o_\alpha(\log N).
\]
More quantitatively, if an upper bound $\varepsilon_*>0$ for the bias
is chosen so that
\[
 v=\log(1/\varepsilon_*)\longrightarrow\infty,
 \qquad v\leq(\log N)^{1/3},
\]
then the remainder is
$O_\alpha(\log N/v+v+\log v+1)$.
\end{theorem}
\begin{proof}
Use the subgroup $J$ and core $R_J$ from \eqref{p-eq:ad-core}.
Put $n=|J|$. We first select $T_0\subseteq R_J$ whose subset products
cover $J$, using
\begin{equation}\label{p-eq:ad-core-count}
 |T_0|\leq\log_2n+
 O_\alpha(\log N/v+v+\log v+1)
\end{equation}
labels. The case $J=\{1\}$ is immediate.

For a probability measure $\nu$ on $J$, define
$E(\nu)=n\sum_x\nu(x)^2-1$ and $b_p=(\delta_1+\delta_p)/2$.
At all times only $O_\alpha(\log N)$ labels will be used. Since
$d/\log N\to\infty$, at least $\alpha d/2$ unused core labels remain,
and the gap of every nonprincipal character on the unused core is
at least $\theta d/2$.

The restriction lemma and nonnegativity of
$\langle\nu,\nu(p^{-1}\cdot)\rangle$ give an unused core label with
\[
 E(\nu*b_p)
 \leq\tfrac12(1+2\varepsilon_*/\alpha)E(\nu)+1/\alpha.
\]
Start with $\nu=\delta_1$. While $E(\nu)>v$, this nearly halves the
energy, so at most $\log_2n+O_\alpha(\log N/v+1)$ labels give
$E(\nu)\leq v$.

For each nonprincipal character, averaging
$|\widehat b_p(\chi)|^2=(1+\operatorname{Re}\chi(p))/2$
over unused core labels gives a value at most $1-\theta/4$.
Averaging the whole Fourier energy therefore supplies an unused label
that contracts it by this factor. At most $O_\alpha(\log v+1)$
additional labels give $E(\nu)\leq\alpha/16$.
The support of $\nu$ then has size at least $n/(1+\alpha/16)$.
All measures constructed are supported on subset products of the
distinct labels selected so far.

Let $B$ be the missing set of the actual subset-product support.
It has at most $\alpha n/16$ elements. Positive restriction of the
overlap estimate to at least $\alpha d/2$ unused core labels supplies
an unused $p$ with
\[
 |B\cap pB|
 \leq\frac2\alpha\left(\frac{|B|^2}{n}
                         +\varepsilon_*|B|\right).
\]
While $\varepsilon_*n<|B|\leq\alpha n/16$, this contracts $|B|$
by a factor less than $1/4$ for sufficiently large members of the
family. After $O(v+1)$ labels, $|B|\leq\varepsilon_*n$.
The contraction factor is now at most $4\varepsilon_*/\alpha$, and
$O_\alpha(\log N/v+1)$ more labels make the integer $|B|$ equal zero.
This proves \eqref{p-eq:ad-core-count}. The stated bounds on $v$ ensure
that the total is $O_\alpha(\log N)$, validating the unused-core
estimates throughout.

The quotient $Q=G/J$ has bounded order $j\leq2/\alpha$.
Scan the labels in $R\setminus J$, retaining a label only if its image
strictly enlarges the current subset-product set in $Q$. If a label
does not enlarge that set, its image fixes the set and continues to
fix every future union of translates, so omitting it is harmless.
At most $j-1$ labels are retained, and their quotient subset products
are exactly those of $R\setminus J$. Call this retained set $U$.
Since $\Sigma(T_0)=J$,
\[
 \Sigma(T_0\cup U)
 =J\Sigma(U)
 =J\Sigma(R\setminus J)
 =\Sigma(R).
\]
Finally, $\log_2|\Sigma(R)|=\log_2|J|+O_\alpha(1)$.
Combining this with \eqref{p-eq:ad-core-count} proves the theorem.
For the qualitative statement, take
$\varepsilon_* =\max\{\varepsilon,e^{-(\log N)^{1/3}}\}$.
\end{proof}

\subsection{Prime application under Dirichlet GRH}

Fix $B>2$ and $0<\alpha\leq1$. For sufficiently large $q$, let
\[
 G=(\mathbb Z/q\mathbb Z)^\times,
 \qquad P_q=\{p\leq(\log q)^B:p\text{ prime},\ p\nmid q\}.
\]
Dirichlet GRH gives\[
 d\sim\frac{(\log q)^B}{B\log\log q},\qquad
 \varepsilon\ll_B(\log\log q)(\log q)^{1-B/2}.
\]
Thus the preceding theorems apply simultaneously to every
$R\subseteq P_q$ of cardinality at least $\alpha|P_q|$.
If its subset products fail to cover $G$, the stabilizer has index
$2\leq h\leq\lfloor1/\alpha\rfloor$, and at most $h-2$ primes of
$R$ lie outside that stabilizer. Its whole subset-product support can
be reproduced by a set $T\subseteq R$ satisfying
\[
 |T|\leq\log_2|\Sigma(R)|+
 O_{\alpha,B}\left(\frac{\log q}{\log\log q}\right).
\]

There is also a bounded-prime repair. If $H$ is the actual stabilizer,
write $G/H\cong C_{n_1}\times\cdots\times C_{n_t}$. Each quotient
coset contains $(1/h+o(1))d$ candidate primes. Choose available primes
lifting the binary powers of a generator in each cyclic factor.
Their subset products cover $G/H$, so adding at most
\[
 b(G/H)=\sum_{i=1}^t\lceil\log_2 n_i\rceil
 \leq h-1\leq\lfloor1/\alpha\rfloor-1
\]
primes repairs full coverage. The inequality follows from
$\lceil\log_2n\rceil\leq n-1$ for $n\geq2$ and
$\sum_i(n_i-1)\leq\prod_i n_i-1$. There are plenty of unused candidates
in each needed nonidentity coset, because $R$ has at most $h-2$ labels
outside $H$. No optimality claim for $b(G/H)$ is made.

The structural exception bound is sharp: for a cyclic quotient of
order $h$, retain the whole kernel pool together with $h-2$ distinct
prime labels in one generating coset. For large $q$ the kernel pool
covers its subgroup, while the quotient subset products are exactly
$0,g,\ldots,(h-2)g$, missing one coset. Its stabilizer is precisely
the kernel and there are exactly $h-2$ exceptional primes. Such
examples have relative density $1/h+o(1)$. To realize every fixed
$h\geq2$ arithmetically, take prime moduli $q\equiv1\pmod h$;
their unit groups have a cyclic quotient of order $h$. The restriction
lemma makes every nonprincipal Fourier coefficient of the kernel
pool $O(\varepsilon d)$, whereas its size is $d/h+O(\varepsilon d)$.
The product estimate \eqref{p-eq:inverse-fourier-product} and
$d/\log q\to\infty$ therefore show that this pool covers the kernel.
The generating coset contains $d/h+O(\varepsilon d)$ available primes,
so the $h-2$ exceptional labels can be chosen distinctly.






\section{Keeping a vanishing fraction of the prime pool}

\subsection{Uniform obstruction structure}

\begin{theorem}[Uniform obstruction structure]\label{p-thm:vanishing-structure}
The all-density obstruction theorem admits a uniform extension. Let
$\alpha=\alpha_N\in(0,1]$ vary, and assume
\begin{equation}\label{p-eq:vd-assumptions}
 \varepsilon=o(\alpha),\qquad
 \frac{\alpha^3 d}{\log N}\longrightarrow\infty.
\end{equation}
For every $R\subseteq P$ with $|R|\geq\alpha d$, if
$S=\Sigma(R)\ne G$, its actual stabilizer $H$ satisfies
\begin{equation}\label{p-eq:vd-structure}
 [G:H]\leq
 \left(\alpha-\varepsilon-
       \frac{C\log N}{\alpha^2d}\right)^{-1}
 =\frac{1+o(1)}{\alpha},
 \qquad
 |R\setminus H|\leq[G:H]-2,
\end{equation}
where $C$ is an absolute constant. As before, $S$ is the full preimage
of the subset-product set of $R\setminus H$ in $G/H$.
\end{theorem}

\begin{proof}
Repeat the missing-element character argument from the all-density
theorem, but at a current subgroup of index $j$ choose
$m_j=\lceil8/(\alpha j)\rceil+1$. Until termination, the current
index is at most $2/\alpha$ and the retained core has
$\alpha d-o(\alpha d)$ labels; these bounds will be justified below.
A missing element supplies $\chi$ with $D(\chi)\leq4\log N$.
If its order were at least $m_j$, the higher Fej\'er estimate would give
\[
 |R\cap J|
 \leq \frac{d/j+\varepsilon d}{m_j}+\varepsilon d
       +O\left(\frac{\log N}{\alpha^2j^2}\right)
 \leq\frac{\alpha d}{8}+o(\alpha d),
\]
contradicting the size of the core. Thus the order $u$ is less than
$m_j=O(1/(\alpha j))$.

For $u\geq2$, the elementary bound
$1-\cos(2\pi/u)\geq8/u^2$ shows that restricting to its kernel loses
at most
\[
 \frac{u^2}{8}\,4\log N
 \ll\frac{\log N}{\alpha^2j^2}
\]
labels. The indices at successive restrictions at least double.
Therefore the total number discarded is at most
$C\log N/\alpha^2=o(\alpha d)$, by summing a geometric series.
A restriction cannot produce an index larger than $2/\alpha$, since
the resulting subgroup would then have at most
$\alpha d/2+\varepsilon d$ pool labels while still containing
$\alpha d-o(\alpha d)$ labels of $R$.
The process consequently terminates after $O(\log(1/\alpha)+1)$ steps,
at a subgroup $J$ with $\Sigma(R\cap J)=J$ and
\[
 |R\cap J|\geq\alpha d-C\log N/\alpha^2.
\]
The restriction estimate $|P\cap J|\leq d/[G:J]+\varepsilon d$
gives the quantitative index bound in \eqref{p-eq:vd-structure}.
Since $J$ is a period of $S$, the actual stabilizer $H$ contains $J$.
The elementary aperiodicity argument in $G/H$ gives
$|R\setminus H|\leq[G:H]-2$ exactly as before.
\end{proof}

\subsection{Uniform prime compression: the exponent one half}

\begin{theorem}[Vanishing retention, exact structure, and compression]\label{p-thm:vanishing-compression}
Assume Dirichlet GRH. Fix
\[
 B>2,\qquad
 0<c<\min\{B/2-1,\tfrac12\}.
\]
Put $\ell=\log q$, $\alpha=\ell^{-c}$, and
\[
 P_q=\{p\leq\ell^B:p\text{ prime},\ p\nmid q\},\qquad d=|P_q|.
\]
For every sufficiently large $q$ and every subset
$R\subseteq P_q$ with $|R|\geq\alpha d$, there is $T\subseteq R$ with
\begin{equation}\label{p-eq:vd-compression}
 \Sigma(T)=\Sigma(R),\qquad
 |T|\leq\log_2|\Sigma(R)|+
 O_{B,c}\left(\frac{\log q}{\log\log q}\right).
\end{equation}
If $\Sigma(R)\ne G=(\mathbb Z/q\mathbb Z)^\times$, then its stabilizer
has index $h\leq(1+o(1))(\log q)^c$, with at most $h-2$ labels
of $R$ outside $H$. With $\pi:G\to G/H$, its support is exactly
\[
 \Sigma(R)=\pi^{-1}\bigl(\Sigma(\pi(R\setminus H))\bigr).
\]
There is $U\subseteq P_q\setminus R$ with
\[
 |U|\leq2c\log_2\log q+O(1),\qquad\Sigma(R\cup U)=G.
\]
The augmented set contains a covering subcollection of at most
$\log_2\varphi(q)+O_{B,c}(\log q/\log\log q)$ primes.
If $\Sigma(R)=G$, no restoration is needed.
\end{theorem}

\begin{proof}
Write $\kappa=B/2-1$ and fix
$\sigma=(c+\kappa)/2$, so $c<\sigma<\kappa$.
The GRH prime-character estimate permits the bias upper bound
$\varepsilon=\ell^{-\sigma}$ for sufficiently large $q$.
Also $d\sim\ell^B/(B\log\ell)$.
In particular \eqref{p-eq:vd-assumptions} holds, because
$B-3c>2-c>1$. This proves the structural assertions.

For compression, first select an arbitrary $R_0\subseteq R$ with
$|R_0|=\lceil\alpha d\rceil$. This upper bound on the initial core
size is essential for the uniform estimate below.

Apply the fixed-threshold core construction, with its constants now
depending explicitly on $\alpha$. Take
\[
 m=\lceil8/\alpha\rceil+1,\qquad
 k=\lceil\log_2(4/\alpha)\rceil+1,
\]
and
\[
 \theta=\min\left\{
 \frac{\alpha}{16((m^2-1)/6)},
 \frac{\alpha\min_{2\leq u<m}(1-\cos(2\pi/u))}{4k}
 \right\}.
\]
Thus $\theta\asymp\alpha^3/k$, with absolute implied constants.
Every restriction has character order less than $m$, loses at most
$\alpha d/(4k)$ core labels, and multiplies the subgroup index by at
least two. Consequently it stops before $k$ steps, at $J\leq G$ with
\begin{equation}\label{p-eq:vd-core-bounds}
 [G:J]\leq2/\alpha,
 \qquad 3\alpha d/4\leq|R_0\cap J|\leq2\alpha d,
 \qquad
 D_{R_0\cap J}(\chi)>\theta d
 \quad(\chi\ne1).
\end{equation}
The estimates in the fixed-density proof are uniform here: its Fej\'er
contradiction uses $\varepsilon=o(\alpha)$, and its total loss is at
most $\alpha d/4$ by the displayed choice of $\theta$.

Fix in advance a core-label budget $M_0=\lceil4\ell\rceil$.
For sufficiently large $q$,
\[
 M_0<\alpha d/4,\qquad 2M_0<\theta d/2,
\]
because $\alpha d/\ell\to\infty$ and
$\theta d\asymp_{B,c}\ell^{B-3c}/(\log\ell)^2\gg\ell$.
Thus, at every stage before the budget is exhausted, at least
$\alpha d/2$ unused core labels remain, their number is at most
$2\alpha d$, and every nonprincipal character retains gap at least
$\theta d/2$. The phase counts below will show that the construction
finishes strictly before $M_0$ labels have been used.

Put $n=|J|$ and
\[
 v=\log(\alpha/\varepsilon)
   =(\sigma-c)\log\ell\asymp_{B,c}\log\ell,
 \qquad V=16v/\alpha.
\]
For the probability-measure energy $E(\nu)=n\|\nu\|_2^2-1$, positive
restriction to at least $\alpha d/2$ unused core labels supplies
\[
 E(\nu*b_p)\leq\tfrac12(1+2\varepsilon/\alpha)E(\nu)+1/\alpha.
\]
While $E(\nu)>V$, this contracts energy by a factor
$1/2+O(1/v)$. Starting at $\delta_1$, at most
$\log_2n+O_{B,c}(\ell/v)$ labels give $E(\nu)\leq V$.

Within the predetermined budget, unused core labels retain character
gap at least $\theta d/2$.
Their number is at most $2\alpha d$, by the deliberate initial
choice of $R_0$. Averaging the squared Fourier multiplier consequently
gives energy contraction at least
\[
 1-\frac{\theta d/2}{2(2\alpha d)}
 =1-\frac{\theta}{8\alpha}.
\]
This is the improvement over the weaker contraction $1-O(\theta)$.
At most
\[
 O\left(\frac{\alpha}{\theta}
               \log\frac{32V}{\alpha}\right)
 =O_{B,c}\bigl(\ell^{2c}(\log\ell)^2\bigr)
 =o_{B,c}\left(\frac\ell{\log\ell}\right)
\]
additional labels give $E(\nu)\leq\alpha/32$.

The missing set of the actual subset-product support now has size
at most $\alpha n/32$. Its positive-restriction update satisfies
\[
 |B\cap pB|\leq\frac2\alpha
       \left(\frac{|B|^2}{n}+\varepsilon|B|\right).
\]
It contracts by a factor less than $1/4$ until $|B|\leq\varepsilon n$,
which costs $O(v+1)$ labels. Thereafter the factor is at most
$4\varepsilon/\alpha=4e^{-v}$, so
$O_{B,c}(\ell/v)$ more labels eliminate all holes.
We have obtained a set $T_0\subseteq R_0\cap J$ with
\[
 \Sigma(T_0)=J,
 \qquad |T_0|\leq\log_2n+
 O_{B,c}(\ell/\log\ell).
\]
The first phase uses at most $\ell/\log2+o_{B,c}(\ell)$ labels, and
all remaining phases together use $o_{B,c}(\ell)$ labels. Their sum
is strictly less than $4\ell$ for sufficiently large $q$. Hence the
construction finishes before its predetermined budget is exhausted,
justifying all unused-core estimates without circularity.

Finally, in $G/J$ scan all labels of the original $R\setminus J$ and
retain a label only when it enlarges the quotient subset-product set.
At most $[G:J]-1\leq2/\alpha=O(\ell^c)$ labels are retained.
Together with $T_0$ they reproduce $\Sigma(R)$ exactly. Their number
is as in \eqref{p-eq:vd-compression}, since
$\ell^c=o(\ell/\log\ell)$ and
$\log_2n\leq\log_2|\Sigma(R)|$.

For repair, use the actual stabilizer $H$. Every quotient coset
contains $(1/[G:H]+O(\varepsilon))d\gg\alpha d$ pool labels.
There are at most $[G:H]-2$ retained labels outside $H$, so the needed
nonidentity cosets contain unused prime labels. If
$G/H\cong\prod_i C_{n_i}$, binary powers in each cyclic factor can be
lifted to distinct available primes, using at most
\[
 \sum_i\lceil\log_2n_i\rceil
 \leq2\log_2[G:H]
 \leq2c\log_2\ell+O(1)=O_{B,c}(\log\ell)
\]
labels. Their subset products cover $G/H$, repairing full coverage.
The augmented set still has at least $\alpha d$ labels, so the same
compression theorem applies again.
\end{proof}

\begin{remark}[Range and limitations]
The structural conclusion and bounded-quotient repair do not require
$c<1/2$. They hold in the larger range
$0<c<\min\{B/2-1,(B-1)/3\}$, which is exactly the range ensuring
$\varepsilon=o(\alpha)$ and $\alpha^3d/\log q\to\infty$ with the
present estimates. The restriction $c<1/2$ is used for the quantitative
near-optimal compression, through its middle energy-contraction stage.
No optimality claim for this exponent is made.
\end{remark}




\begin{corollary}[Size of the attainable support]\label{p-cor:vanishing-support}
Under Theorem~\ref{p-thm:vanishing-compression},
\[
 |\Sigma(R)|\geq(1-o(1))
 \frac{\varphi(q)}{(\log q)^c},
 \qquad
 \log_2|\Sigma(R)|
 \geq\log_2\varphi(q)-c\log_2\log q-o(1).
\]
\end{corollary}
\begin{proof}
If the support is incomplete, the identity lies in it and its stabilizer
is a period, so the support contains that stabilizer. The index bound
therefore proves both inequalities. They also hold when the support
is the full group.
\end{proof}

For example, $B=3$ and $c=2/5$ are admissible. A retained fraction
$(\log q)^{-2/5}$ can be repaired with at most
$(4/5)\log_2\log q+O(1)$ primes and then compressed to first-order
minimum size. The number of attainable residue classes loses only
$O(\log\log q)$ in its base-two logarithm, although the deleted
fraction tends to one. The compressed list is chosen after the
deletion. This statement concerns residue coverage and the number
of shared factors, not the numerical size of the integer products.



\section{Products with a prescribed number of distinct prime factors}
The one-third theorem concerns subsets of all cardinalities together.
An exact cardinality has an additional obstruction: if every retained
prime is a quadratic nonresidue, products of exactly $k$ of them all
have the fixed character value $(-1)^k$. The next theorem therefore
uses a stronger, explicit condition.

\begin{theorem}[Exact-cardinality equidistribution]\label{p-thm:exact-k}
Let $P$ be an $\varepsilon$-balanced pool of $d$ labels in a group of
order $N$, and let $R\subseteq P$ have $r$ labels. Put
\[
 \Delta=2r-(1+\varepsilon)d>0.
\]
For $1\le k\le r-1$, let $C_{R,k}(a)$ count the $k$-element subsets
of $R$ with product $a$. Then
\begin{equation}\label{p-eq:exact-k}
 \left|\frac{N C_{R,k}(a)}{\binom rk}-1\right|
 \le (N-1)(r+1)
 \exp\left(-\frac{k(r-k)}{r^2}\Delta\right).
\end{equation}
In particular, if the right-hand side is less than one, every element
has a representation by exactly $k$ distinct labels.
\end{theorem}
\begin{proof}
For a nonprincipal $\psi$,
\[
 \left|\sum_{p\in R}\psi(p)\right|
 \le\varepsilon d+d-r=r-\Delta.
\]
Write $\theta=k/r$ and $z=\theta/(1-\theta)>0$. For $|w|=1$,
\[
 |1+zw|\le(1+z)
 \exp\left(-\frac{z}{(1+z)^2}(1-\Re w)\right).
\]
This follows by squaring and using $1-u\le e^{-u}$.
Uniformly in real $t$, it implies
\[
 \prod_{p\in R}|1+ze^{it}\psi(p)|
 \le(1+z)^r e^{-\theta(1-\theta)\Delta}.
\]
The Cauchy coefficient estimate for
$\prod_{p\in R}(1+Z\psi(p))$ bounds its $k$th coefficient by
$z^{-k}(1+z)^r e^{-\theta(1-\theta)\Delta}$.
The binomial probability
$\binom rk\theta^k(1-\theta)^{r-k}$ is a mode of the
$\operatorname{Bin}(r,\theta)$ distribution, so it is at least
$1/(r+1)$. Consequently the coefficient divided by $\binom rk$ has
modulus at most $(r+1)e^{-\theta(1-\theta)\Delta}$.
Character orthogonality, with principal term $\binom rk/N$, proves
\eqref{p-eq:exact-k}.
\end{proof}

For the prime pool, retention $r\ge(1/2+\eta)d$ gives
$\Delta\ge(2\eta-\varepsilon)d$. Every fixed central range
$\theta\in[\theta_0,1-\theta_0]$, $\theta_0>0$, then has uniform
exact-cardinality counts with relative error smaller than every fixed
negative power of $q$. More generally, equation~\eqref{p-eq:exact-k}
itself supplies a finite sufficient condition. No such conclusion is
asserted for all fixed cardinalities at one-third retention.

\section{A version using primes in one narrow window}
\begin{corollary}\label{p-cor:window}
Assume $\pGRHD$, and fix
\[
 B>2,\qquad B/2+1<D\le B,\qquad
 0<c<\min\{D-B/2-1,1\}.
\]
Replace $P_q$ in Theorem~\ref{p-thm:main} by
\[
 P_q^{\mathrm{win}}=
 \{p:\ell^B<p\le\ell^B+\ell^D,\ p\text{ prime},\ p\nmid q\}.
\]
All conclusions of that theorem remain valid, including one-third
retention, the quadratic dichotomy, one-prime repair, and near-minimal
compression.
\end{corollary}
\begin{proof}
Subtracting \eqref{p-eq:prime-character} at the two endpoints gives a
nonprincipal sum $O_B(\ell^{B/2+1})$. The prime number theorem on RH
gives
$d\sim\ell^D/(B\log\ell)$; deleting prime divisors of $q$ changes
this by only $O(\ell)$. The normalized bias is therefore
$O_B((\log\ell)\ell^{-(D-B/2-1)})$.
Use this positive exponent in place of $\kappa$ in
Section~\ref{p-sec:compression} and the prime specialization.
\end{proof}
For example, the window
$(\ell^6,\ell^6+\ell^5]$ has relative width $1/\ell$ and permits
every $c<1$. Thus all selected prime factors may be constrained to a
single shrinking relative-width interval.



\section{A weighted refinement near the Erd\H{o}s--Odlyzko--S\'ark\"ozy problem}

Throughout this section, assume GRH for all primitive Dirichlet
$L$-functions, including the Riemann zeta function. All logarithms are
natural. For an integer $q\geq3$, put
\[
 G_q=(\mathbb Z/q\mathbb Z)^\times,
 \qquad n=\varphi(q),
 \qquad
 E_2(q)=\{p_1p_2\bmod q:p_1,p_2\leq q\text{ prime},\ p_1p_2\perp q\}.
\]
Here $p_1p_2\perp q$ means $\gcd(p_1p_2,q)=1$; the two primes need not
be distinct. The original Erd\H{o}s--Odlyzko--S\'ark\"ozy conjecture asks
for $E_2(q)=G_q$ for all sufficiently large prime $q$.
The estimates below also apply when $q$ is composite. The original
formulation is in~\cite{p-EOS1987}, and the general-modulus setting is
discussed in~\cite{p-Szabo2024}. The primary abstract of
Banks~\cite{p-Banks2026} explicitly records that full two-prime
coverage remains open even under GRH. No theorem below asserts it.

\begin{lemma}[Triangularly weighted prime sums]\label{p-lem:eos-smoothed}
Let $\chi$ be a Dirichlet character modulo $q$, induced by a primitive
character $\chi^*$ of conductor $r\mid q$. Write $\delta_\chi=1$ when
$\chi$ is principal and $\delta_\chi=0$ otherwise. Uniformly for $x\geq2$,
\begin{equation}\label{p-eq:eos-smoothed}
 S_\chi(x):=\sum_{p\leq x}\left(1-\frac px\right)(\log p)\chi(p)
 =\delta_\chi\frac x2
  +O\bigl(\sqrt{x}\log(2r)+\log(2q)\bigr).
\end{equation}
The implied constant is absolute. In particular, at $x=q$,
\begin{equation}\label{p-eq:eos-at-q}
 S_{\chi_0}(q)=\frac q2+O(\sqrt q),
 \qquad
 \max_{\chi\ne\chi_0}|S_\chi(q)|\ll\sqrt q\log q.
\end{equation}
\end{lemma}

\begin{proof}
We include the explicit-formula argument to keep track of the smoothing
and the dependence on the conductor. This is the standard first-order
smoothed estimate; see Montgomery~\cite[\S13, (13.19)]{p-Montgomery1971}
and, for its quadratic-character form,
Shparlinski--Sutherland~\cite[\S3.1]{p-SS2017}.

For a primitive character $\chi^*$, Mellin inversion gives
\[
 \Psi_1(x,\chi^*)
 :=\sum_{m\leq x}\left(1-\frac mx\right)\Lambda(m)\chi^*(m)
 =\frac1{2\pi i}\int_{(2)}
   -\frac{L'}{L}(s,\chi^*)\frac{x^s}{s(s+1)}\,ds.
\]
Move the contour to $\operatorname{Re}s=-1/2$, taking horizontal
heights that stay away from zero ordinates and then tending to infinity.
The factor $1/[s(s+1)]$ makes the resulting zero sum absolutely
convergent. The residues at the nontrivial zeros and at the pole $s=1$
give
\begin{equation}\label{p-eq:eos-explicit}
 \Psi_1(x,\chi^*)
 =\delta_\chi\frac x2
   -\sum_{\rho}\frac{x^\rho}{\rho(\rho+1)}
   +O\bigl(\log(2rx)\bigr),
\end{equation}
where zeros are counted with multiplicity.

Here are bounds for the remaining terms in this contour calculation.
The functional equation and the absolutely convergent Dirichlet series
on $\operatorname{Re}s=3/2$ give
\[
 \left|\frac{L'}{L}(-1/2+it,\chi^*)\right|
 \ll\log\bigl(2r(|t|+2)\bigr).
\]
Thus the integral on the new vertical line is
$O(x^{-1/2}\log(2r))$. For a nonprincipal primitive character, the
regular part of $L'/L$ at $s=0$ is $O(\log(2r))$ under GRH. To see
this directly, the logarithmic derivative of the completed $L$-function
at $s=1$ differs from its value at $s=2$ by the absolutely convergent
sum
\[
 \sum_\rho\left(\frac1{1-\rho}-\frac1{2-\rho}\right)
 \ll\log(2r).
\]
Indeed, under GRH its summands are $O((1+|\operatorname{Im}\rho|^2)^{-1})$,
and the standard local zero count is
\[
 \#\{\rho:T\leq|\operatorname{Im}\rho|<T+1\}
 \ll\log\bigl(2r(T+2)\bigr).
\]
The value at $s=2$ is $O(\log(2r))$ by the Euler product; the
functional equation transfers the bound to the regular part at $s=0$.
Consequently the residue at $s=0$ is $O(\log(2r)+\log(2x))$, including
the possible double pole for even $\chi^*$. For the primitive principal
character, $L(s,\chi^*)=\zeta(s)$ and this residue is bounded absolutely.
These observations justify the remainder in \eqref{p-eq:eos-explicit}.

Since $\rho=1/2+i\gamma$ under GRH, the same local zero count gives
\[
 \sum_\rho\frac1{|\rho(\rho+1)|}\ll\log(2r).
\]
Equation~\eqref{p-eq:eos-explicit} therefore implies
\[
 \Psi_1(x,\chi^*)
 =\delta_\chi\frac x2+O\bigl(\sqrt{x}\log(2r)\bigr),
\]
where $\log(2x)\ll\sqrt x$ absorbs the remaining logarithm.

Removing the higher prime powers costs at most
\[
 \sum_{\substack{k\geq2\\p^k\leq x}}\log p
 \leq\sum_{2\leq k\leq\log_2x}\vartheta(x^{1/k})
 \ll\sqrt x+x^{1/3}\log(2x)\ll\sqrt x,
\]
by the elementary Chebyshev bound $\vartheta(t)\ll t$.
Finally, on primes the induced and primitive characters differ only
at primes dividing $q$ but not $r$. Their total contribution is at most
\[
 \sum_{p\mid q}\log p\leq\log q.
\]
This proves \eqref{p-eq:eos-smoothed} without an extra logarithm from
imprimitive Euler factors. For comparison, if one instead compares
the full von Mangoldt sums before removing prime powers, the discrepancy
is bounded by $\omega(q)\log x$; that bound also suffices at $x=q$.
For the principal character $r=1$, which gives the sharper first
estimate in \eqref{p-eq:eos-at-q}.
\end{proof}

\begin{theorem}[At most a logarithmic exceptional set]\label{p-thm:eos-log3}
For every sufficiently large integer $q$,
\begin{equation}\label{p-eq:eos-log3}
 |G_q\setminus E_2(q)|
 \ll\left(\frac{\varphi(q)}q\right)^2(\log q)^3
 \ll(\log q)^3.
\end{equation}
All implied constants are absolute.
\end{theorem}

\begin{proof}
Put $w(p)=(1-p/q)\log p$ for primes $p<q$ with $p\nmid q$, and
view these weights as a function $w:G_q\to[0,\infty)$, zero elsewhere.
The primes in question occupy distinct residue classes, also for
composite $q$. With the unnormalized Fourier transform
\[
 \widehat w(\chi)=\sum_{a\in G_q}w(a)\chi(a),
 \qquad W=\widehat w(\chi_0),
\]
Lemma~\ref{p-lem:eos-smoothed} and Chebyshev's bound give
\[
 W=\frac q2+O(\sqrt q),\qquad
 B:=\max_{\chi\ne\chi_0}|\widehat w(\chi)|\ll\sqrt q\log q,
 \qquad
 \sum_{a\in G_q}w(a)^2
 \leq(\log q)\vartheta(q)\ll q\log q.
\]
Use the unnormalized multiplicative convolution
$(w*w)(a)=\sum_{b\in G_q}w(b)w(ab^{-1})$.
Parseval and the convolution identity yield
\begin{align*}
 \sum_{a\in G_q}\left|(w*w)(a)-\frac{W^2}{n}\right|^2
 &=\frac1n\sum_{\chi\ne\chi_0}|\widehat w(\chi)|^4\\
 &\leq\frac{B^2}{n}\sum_\chi|\widehat w(\chi)|^2
  =B^2\sum_{a\in G_q}w(a)^2
  \ll q^2(\log q)^3.
\end{align*}
If $a\notin E_2(q)$, then $(w*w)(a)=0$. Thus, writing
$K=|G_q\setminus E_2(q)|$,
\[
 K\frac{W^4}{n^2}\ll q^2(\log q)^3.
\]
Since $W\gg q$, this is precisely \eqref{p-eq:eos-log3}.
\end{proof}

\begin{corollary}[One fixed correction prime]\label{p-cor:eos-correction}
For every sufficiently large integer $q$, there is one prime
$\ell\nmid q$ such that
\begin{equation}\label{p-eq:eos-correction}
 \ell\ll(\log q)^6\log\log q,
 \qquad
 G_q=E_2(q)\cup\ell E_2(q).
\end{equation}
Equivalently, the same prime $\ell$, chosen once for $q$, lets every
invertible residue be represented either by two primes at most $q$
or by $\ell$ times two primes at most $q$. Repeated prime factors
are permitted.
\end{corollary}

\begin{proof}
Let $E=G_q\setminus E_2(q)$ and $K=|E|$.
The set $EE^{-1}$ has at most $K^2$ elements. Among the first
\[
 M=K^2+\omega(q)+1
\]
primes, at most $\omega(q)$ divide $q$. Since $K\ll\log^3q$,
the standard bound for the $M$th prime gives
\[
 p_M\ll M\log(2M)\ll(\log q)^6\log\log q<q
\]
for sufficiently large $q$. The remaining primes therefore represent
distinct elements of $G_q$. At most $K^2$ of them can belong to
$EE^{-1}$, so some prime $\ell\leq p_M$ satisfies
$\ell\nmid q$ and $\ell\bmod q\notin EE^{-1}$.
It follows that $E\cap\ell E=\varnothing$. Since the complement of
$\ell E_2(q)$ is $\ell E$, this proves \eqref{p-eq:eos-correction}.
The argument includes $E=\varnothing$.
\end{proof}

\paragraph{Relation to earlier work and the remaining problem.}
For prime moduli, Erd\H{o}s, Odlyzko and S\'ark\"ozy proved in
1987 that under their quasi-Riemann hypothesis the exceptional set
has size $O(q^{2\theta-1}\log^5q)$; GRH gives the special case
$O(\log^5q)$~\cite{p-EOS1987}. Theorem~\ref{p-thm:eos-log3} removes two
logarithms from that displayed historical bound by using a triangular
weight before applying Parseval. This is an elementary conditional
refinement assembled here, not a verified claim of priority: the
literature search has not established whether this smoothing argument
or a stronger bound already appears elsewhere. Neither the exceptional
set estimate nor the fixed-prime correction proves $E_2(q)=G_q$.
They leave the possible existence of even one exceptional residue
unresolved.





\section{An inversion-pair obstruction to scalar Fourier arguments}

\begin{theorem}\label{p-thm:inversion-countermodel}
Let $q\ge257$ be a prime and $a\in\mathbb F_q^\times$.
There exists a set $A\subseteq\mathbb F_q^\times$ such that
\[
 |A|=\left\lfloor\frac{q-1}{\log q}\right\rfloor,
 \qquad a\notin A\cdot A,
 \qquad
 \max_{\chi\ne1}\left|\sum_{x\in A}\chi(x)\right|
 \le12\sqrt q.
\]
Here the maximum is over the nonprincipal multiplicative characters
of $\mathbb F_q^\times$, and repeated factors are allowed in $A\cdot A$.
\end{theorem}

\begin{proof}
The map $\tau(x)=a/x$ is an involution.  It has $f\in\{0,2\}$ fixed
points, namely the solutions of $x^2=a$.  Discard those fixed points and
partition the other elements into $m=(q-1-f)/2$ two-element orbits.
Independently for every orbit $\{x,y\}$, select no element with probability
$1-2/\log q$, select $x$ with probability $1/\log q$, and select $y$
with probability $1/\log q$.  Denote the resulting random set by $A_0$.
At most one point is selected from each orbit, so
$a\notin A_0\cdot A_0$ deterministically.

The random size $Z=|A_0|$ has mean
\[
 \mu=\frac{q-1-f}{\log q}.
\]
For a nonprincipal character $\chi$, put $S_\chi=\sum_{x\in A_0}\chi(x)$.
Character orthogonality gives
\[
 \mathbb E S_\chi
 =-\frac1{\log q}\sum_{\tau(x)=x}\chi(x),
 \qquad |\mathbb E S_\chi|\le\frac2{\log q}.
\]
Both the real and imaginary parts of $S_\chi$ are sums of independent
orbit contributions.  Each centred contribution has absolute value at
most $2$, and the sum of variances is at most
$\mu\le q/\log q$.  The same bounds apply to $Z-\mu$.

Bernstein's inequality, with $t=4\sqrt q$, shows for each such real
sum $W$ that
\[
 \mathbb P(|W-\mathbb EW|>4\sqrt q)
 \le2\exp\!\left(-\frac{16q}{2(q/\log q+(8/3)\sqrt q)}\right)
 \le2q^{-4}.
\]
For the last inequality use
$\log q\le(3/8)\sqrt q$, valid for $q\ge257$.
The union bound over the real and imaginary parts of all $q-2$
nonprincipal characters and over $Z$ has failure probability at most
\[
 (4(q-2)+2)q^{-4}<1.
\]
There is therefore a realization satisfying
\[
 |Z-\mu|\le4\sqrt q,
 \qquad
 |S_\chi|\le4\sqrt{2q}+\frac2{\log q}
 \quad(\chi\ne1).
\]

Let $k=\lfloor(q-1)/\log q\rfloor$.  Since $k\le m$, we may obtain
a set $A$ of exactly $k$ elements by deleting elements if $Z>k$, or by
adding one element from each of sufficiently many unused two-element
orbits if $Z<k$.  The number of changes is at most
\[
 |Z-k|\le4\sqrt q+2.
\]
These changes preserve $a\notin A\cdot A$, and change any character sum
by at most their number.  Thus
\[
 \left|\sum_{x\in A}\chi(x)\right|
 \le(4\sqrt2+4)\sqrt q+2+\frac2{\log q}
 \le12\sqrt q.
\]
\end{proof}

\paragraph{Scope.}
This is not a counterexample to the Erd\H{o}s--Odlyzko--S\'ark\"ozy
conjecture: the constructed sets are not sets of primes.  It proves only
that prime-sized cardinality together with uniform scalar multiplicative
character bounds---even $O(\sqrt q)$ bounds, stronger than the usual
GRH bound $O(\sqrt q\log q)$ for primes up to $q$---does not logically
force every residue to be a product of two set elements.  The construction
is a standard independent-orbit probabilistic argument, and no novelty
claim is made for it.





\subsection{The exact two-prime term that remains uncontrolled}

For prime $q$, let $\mathcal P_q$ be the primes below $q$ and put
$F_q(\chi)=\sum_{p\in\mathcal P_q}\chi(p)$. The unweighted number
of ordered two-prime representations of a unit $a$ is exactly
\[
 \frac1{q-1}\left(
 |\mathcal P_q|^2+
 \sum_{\chi\ne1}\overline{\chi(a)}F_q(\chi)^2\right).
\]
The unresolved issue is controlling the signed, cross-character sum
well enough at every individual $a$. The GRH scalar upper bounds do
not establish the required positivity. The inversion-pair construction
above proves that this logical gap is real for arbitrary prime-sized
sets. It does not prove that GRH itself cannot imply EOS by some other,
prime-specific argument.

This distinction also appears in another major direction:
Kandhil--Languasco--Moree~\cite{p-KLM2026} derive a corrected Montgomery
conjecture and Elliott--Halberstam from GRH together with an additional
cross-character zero-pair-correlation conjecture. Their extra
correlation assumption is not supplied by the finite-pool arguments
in the present manuscript.



\section{A second universality result: prime quadratic-residue graphs}\label{p-sec:prime-graphs}
This section supplies a separate consequence of the same hypothesis.
The qualitative graph on all primes $p\equiv1\pmod4$, with an edge when
$(p/r)=1$, is already Cameron's prime realization of the countable random
graph~\cite{p-Cameron}. We make no novelty claim for that construction.
The result below is a uniform quantitative, deletion-resilient
formulation for short intervals and progressions, derived from standard
analytic inputs.

\begin{theorem}[Induced patterns in short progressions]\label{p-thm:graphs}
Assume $\pGRHD$. Fix
\[
 0<e\le\tfrac12,\quad 0\le u<e,\quad
 0<v<e-u,\quad 0<\alpha\le1.
\]
Set $Y=X^{1/2+e}$. For every sufficiently large $X$, simultaneously
for every integer $q\le X^u$ and every reduced residue $a\bmod q$
with $a\equiv1\pmod{\gcd(q,4)}$, let
\[
 V=\{p:X<p\le X+Y,\ p\text{ prime},\ p\equiv a\pmod q,
                  \ p\equiv1\pmod4\}.
\]
Every $W\subseteq V$ with $|W|\ge\alpha|V|$ contains an induced copy
of every simple graph on at most
\[
 \left\lfloor(e-u-v)\log_2X\right\rfloor
\]
vertices, where distinct primes are adjacent precisely when their
Legendre symbol is $+1$.
\end{theorem}

\begin{lemma}[An explicit spectral embedding bound]\label{p-lem:embedding}
Let $L$ be a real symmetric matrix indexed by $V$, with diagonal zero,
off-diagonal entries in $\{-1,1\}$, and $\|L\|_{\mathrm{op}}\le\lambda$,
$\lambda\ge1$. If $W\subseteq V$ has size
$m>8n^{3/2}2^n\lambda$, then every simple graph on $n\ge2$ vertices
has an induced embedding in $W$, interpreting $+1$ as adjacency.
\end{lemma}
\begin{proof}
Put $\beta=(1-1/n)/2$. Maintain a candidate set $R_j\subseteq W$ for
each unchosen target vertex, initially $W$; the sets may overlap.
For a future set $T$ and required sign $\sigma$, the available companions
of a vertex $x$ number
\[
 N_\sigma(x,T)=
 \frac{|T|-\pone_T(x)+\sigma(L\pone_T)_x}{2}.
\]
If this is smaller than $\beta|T|$ and $|T|\ge2n$, then
$|(L\pone_T)_x|>|T|/(2n)$. The norm bound therefore permits at most
$4n^2\lambda^2/|T|$ such bad vertices.

Let $t=m\beta^{n-1}$. Since
$(1-1/n)^{n-1}\ge e^{-1}$, the assumed bound gives
$t>4n^{3/2}\lambda\ge2n$. Throughout the selection, each candidate
set has size at least $t$ if every update retains a $\beta$-fraction.
The union of bad vertices for fewer than $n$ future sets has size at
most $4n^3\lambda^2/t<t$. Thus the current candidate set contains a
vertex good for all future requirements. Choose it, intersect the
future sets with its required neighbors or nonneighbors, and exclude
that vertex. Induction completes an injective induced embedding.
\end{proof}

\begin{proof}[Proof of Theorem~\ref{p-thm:graphs}]
Put $Q=\operatorname{lcm}(4,q)$. Compatibility gives one reduced residue modulo $Q$.
The GRH prime-count estimate from \eqref{p-eq:prime-character} yields,
uniformly,
\[
 |V|=\frac{\pli(X+Y)-\pli(X)}{\varphi(Q)}+O(\sqrt X\log X)
 \sim\frac{Y}{\varphi(Q)\log X}
 \gg\frac{X^{1/2+e-u}}{\log X}.
\]
The relative error is $O(X^{u-e}\log^2X)=o(1)$.

For primes in $V$, put $L_{pr}=(r/p)$. Quadratic reciprocity makes
$L$ symmetric. Heath-Brown's quadratic large sieve~\cite[Theorem 1]{p-HB}
states, for every $t>0$,
\[
 \sum_{\substack{m\le M\\m\text{ odd squarefree}}}
 \left|\sum_{\substack{n\le N\\n\text{ odd squarefree}}}
 a_n\left(\frac nm\right)\right|^2
 \ll_t (MN)^t(M+N)
       \sum_{\substack{n\le N\\n\text{ odd squarefree}}}|a_n|^2.
\]
Take $M=N=\lceil2X\rceil$, support the coefficients on $V$, and
restrict the outer sum to $V$. This gives the unconditional bound
$\|L\|_{\mathrm{op}}\ll_t X^{1/2+t}$.
Choose $t=v/2$. For $n\le(e-u-v)\log_2X$,
\[
 8n^{3/2}2^n\|L\|_{\mathrm{op}}
 \ll_v(\log X)^{3/2}X^{1/2+e-u-v/2}
 =o\!\left(\frac{X^{1/2+e-u}}{\log X}\right).
\]
Lemma~\ref{p-lem:embedding} applies to every allowed $W$.
\end{proof}

Any fixed surviving fraction is allowed, for example $\alpha=0.01$,
with the sufficiently-large threshold depending on $\alpha$. The
embedding is chosen after deletion and may differ for different target
graphs. Arbitrary preselected seed vertices need not retain all possible
extensions. For $q=1$ and $Y=X$, this graph consequence is already
unconditional by the prime number theorem in the fixed progression
$1\pmod4$; in short intervals the displayed conditional theorem supplies
the necessary prime count. This distinction prevents attributing every
part of the graph argument to GRH.

\subsection{A verified Petersen graph inside one prime interval}
The following ordered primes all lie in $(10^6,10^6+10^4]$ and are
$1\pmod4$:
\begin{center}
\begin{tabular}{rrrrr}
\toprule
$p_0$ & $p_1$ & $p_2$ & $p_3$ & $p_4$\\
1000033 & 1000037 & 1000213 & 1000537 & 1000429\\
\midrule
$p_5$ & $p_6$ & $p_7$ & $p_8$ & $p_9$\\
1000289 & 1002109 & 1005313 & 1000117 & 1009097\\
\bottomrule
\end{tabular}
\end{center}
Their $45$ pairwise Legendre symbols give exactly the Petersen graph.
All ten primalities were checked by trial division and every pairwise
sign by Euler's criterion. This is a finite unconditional certificate,
not a verification of the asymptotic threshold in the theorem.


With indices $0,\ldots,4$ interpreted modulo five, the edge set is
exactly
\[
 \{\{p_i,p_{i+1}\}:0\leq i<5\}
 \ \cup\ \{\{p_i,p_{i+5}\}:0\leq i<5\}
 \ \cup\ \{\{p_{i+5},p_{5+(i+2\bmod5)}\}:0\leq i<5\}.
\]
This describes the outer pentagon, the five spokes, and the inner
pentagram without relying on a drawing.




\section{Finite certificates and what they verify}\label{p-sec:certificates}
The following examples were computed by exact modular arithmetic.
Their coverage is independently verifiable without GRH; they illustrate
the definitions and conclusions but do not establish the unknown
numerical onset of the asymptotic theorems.

For each prime $q$, take all primes $p\le\lfloor(\log q)^3\rfloor$.
Retain $r=\lceil0.4d\rceil$ primes by taking the smallest $r-1$
quadratic-residue primes and the smallest quadratic-nonresidue prime.
The displayed $t$ is the size of a greedily extracted covering set;
the lower bound is $\lceil\log_2(q-1)\rceil$.
\begin{center}
\begin{tabular}{rrrrrr}
\toprule
$q$ & pool $d$ & retained $r$ & covered & $t$ & lower bound\\
\midrule
101 & 25 & 10 & 100 & 8 & 7\\
1,009 & 66 & 27 & 1,008 & 12 & 10\\
10,007 & 137 & 55 & 10,006 & 16 & 14\\
100,003 & 241 & 97 & 100,002 & 20 & 17\\
\bottomrule
\end{tabular}
\end{center}
For $q=101$, the eight selected primes are
\[
 (5,13,17,23,2,31,43,47).
\]
The successive numbers of represented nonzero residues are
\[
 1,2,4,8,16,32,56,88,100.
\]
Retaining only the smallest $r$ quadratic-residue primes instead gives
exactly $(q-1)/2$ represented residues in each of these four examples.

The cubic obstruction is already visible at $q=1009$. Among the
$66$ primes at most $330$, exactly $22$ belong to a cubic-character
kernel. They squarefreely cover its $336$ residues. Adding the prime
$2$, whose cubic-character value is a nontrivial cube root, produces
a $23$-prime set that generates all $1008$ units but has only $672$
squarefree-product residues. At $q=100003$, the analogous figures are
$241$ candidate primes, $80$ kernel primes, and $66668$ represented
residues after adding one prime, although the generated group has
$100002$ elements. These finite examples also explain why an
unqualified statement with merely $r>d/3$ would be too strong.

\subsection{Constructive meaning and computational limits}
The proofs prescribe finite choices: form a quadratic-basis set,
select unused primes decreasing Fourier energy, and then select primes
that shrink the missing set. All selections can be implemented by
enumerating the group and comparing exact correlations. The compression
proof guarantees existence and a finite construction; it does not give
a polynomial-in-$\log q$ decoder or extraction algorithm. The complete
representation table can have $\varphi(q)$ entries.



\section{Self-contained verification code for the finite certificates}
The following Python 3 code uses only the standard library. It verifies
the displayed compressed covers, the quadratic half-group examples,
both cubic obstructions, and the Petersen certificate. It is not part of
the asymptotic proofs.
\begin{lstlisting}[language=Python,basicstyle=\ttfamily\footnotesize,breaklines=true,columns=fullflexible,keepspaces=true,showstringspaces=false]
from math import isqrt, log, ceil

def prime(n):
    return n >= 2 and all(n % d for d in range(2,isqrt(n)+1))

def primes_up_to(n):
    flags = bytearray(b"\x01") * (n+1)
    flags[:2] = b"\x00\x00"
    for p in range(2,isqrt(n)+1):
        if flags[p]:
            flags[p*p:n+1:p] = b"\x00" * ((n-p*p)//p+1)
    return [i for i in range(2,n+1) if flags[i]]

def products(q, ps):
    out = {1}
    for p in ps:
        out |= {p*a % q for a in list(out)}
    return out

certificates = [
    (101, [5, 13, 17, 23, 2, 31, 43, 47]),
    (1009, [2, 3, 5, 7, 37, 67, 151, 11, 191, 157, 199, 29]),
    (10007, [2, 3, 11, 13, 19, 47, 53, 61, 79, 5, 139, 433, 271, 211, 421, 71]),
    (100003, [11, 17, 23, 41, 43, 71, 79, 83, 97, 101, 107, 131, 2, 911, 751, 857, 173, 733, 263, 113]),
]
for q, selected in certificates:
    assert prime(q)
    pool = primes_up_to(int(log(q)**3))
    positive = [p for p in pool if pow(p,(q-1)//2,q)==1]
    negative = [p for p in pool if pow(p,(q-1)//2,q)==q-1]
    r = ceil(0.4*len(pool))
    retained = positive[:r-1] + negative[:1]
    assert len(set(selected)) == len(selected)
    assert set(selected) <= set(retained)
    assert len(products(q, selected)) == q-1
    assert len(products(q, positive[:r])) == (q-1)//2

for q, d, k, symbol in [
        (1009, 66, 22, 374), (100003, 241, 80, 7120)]:
    pool = primes_up_to(int(log(q)**3))
    exponent = (q-1)//3
    kernel = [p for p in pool if pow(p, exponent, q)==1]
    assert len(pool)==d and len(kernel)==k
    assert pow(2, exponent, q)==symbol and symbol!=1
    assert len(products(q, kernel))==exponent
    assert len(products(q, kernel+[2]))==2*exponent
    # The kernel is covered and 2 lies outside it.
    # Therefore the retained labels generate the full group.

ps = [1000033,1000037,1000213,1000537,1000429,
      1000289,1002109,1005313,1000117,1009097]
assert all(prime(p) and p%4==1 and 10**6<p<=1010000
           for p in ps)
edges = set()
for i,p in enumerate(ps):
    for j in range(i):
        r = ps[j]
        a = 1 if pow(p,(r-1)//2,r)==1 else -1
        b = 1 if pow(r,(p-1)//2,p)==1 else -1
        assert a==b
        if a==1: edges.add((j,i))
expected = {tuple(sorted((i,(i+1)%5))) for i in range(5)}
expected |= {(i,i+5) for i in range(5)}
expected |= {tuple(sorted((i+5,5+(i+2)%5))) for i in range(5)}
assert edges==expected and len(edges)==15
print("All finite certificates verified.")
\end{lstlisting}




\section{A standard transcendence consequence and an independence barrier}
\label{p-sec:transcendence-scope}

The discussion of prime-representing constants contributes a standard
consequence of an external theorem, not an additional discovery.
Let $\xi_3$ denote the ordinary Mills constant: the least real
$A>1$ for which $\lfloor A^{3^k}\rfloor$ is prime for every integer
$k\geq1$. Saito's theorem~\cite[Theorem 1.8]{p-Saito2026} proves that
$\xi_3$ is transcendental under the Density Hypothesis for the
Riemann zeta function.

\begin{corollary}[Saito's consequence under RH]\label{p-cor:mills}
Assuming RH, the ordinary Mills constant is transcendental.
\end{corollary}
\begin{proof}
Write $N(\sigma,T)$ for the number of nontrivial zeta zeros with
real part at least $\sigma$ and imaginary part between $0$ and $T$.
The Density Hypothesis asks, for each fixed $\delta>0$, for
\[
 N(\sigma,T)\ll_\delta T^{\,2(1-\sigma)+\delta}
 \qquad(1/2\leq\sigma\leq1).
\]
Under RH this count is zero for $\sigma>1/2$, and at $\sigma=1/2$
the usual zero-counting formula gives
$N(1/2,T)\ll T\log(2T)\ll_\delta T^{1+\delta}$.
Thus RH implies the Density Hypothesis. Saito's theorem applies.
\end{proof}

We also considered algebraic independence of $\xi_3$ and the analogous
quartic constant $\xi_4$, defined using exponents $4^k$.
Saito--Takeda~\cite{p-ST2022} prove algebraic independence for families
with sufficiently rapidly increasing exponent ratios; that theorem's
hypotheses do not cover the two constant-ratio sequences considered
here. A polynomial relation between $\xi_3$ and $\xi_4$ does not imply
that either constant is algebraic over $\mathbb Q$. Consequently a
one-variable argument whose hypothesis is algebraicity cannot be
applied to such a relation.

The direct norm method has the following additional quantitative
limitation. Put
\[
 p_n=\lfloor\xi_3^{3^n}\rfloor,\quad
 r_m=\lfloor\xi_4^{4^m}\rfloor,\quad
 a_n=p_n^{1/3^n},\quad b_m=r_m^{1/4^m}.
\]
For $K=3^n$ choose $m$ so that $K\leq4^m<4K$.
The mean value theorem applied to $t^{1/K}$, together with
$0\leq\xi_3^K-p_n<1$, gives
\[
 |\xi_3-a_n|\ll K^{-1}e^{-c_1K}.
\]
The analogous estimate gives
$|\xi_4-b_m|\ll K^{-1}e^{-c_2K}$, with positive constants depending
only on the two fixed constants.

If a fixed nonzero polynomial $F\in\mathbb Z[X,Y]$ satisfied
$F(\xi_3,\xi_4)=0$, its bounded derivatives on a compact neighborhood
would therefore imply
\[
 |F(a_n,b_m)|\ll_F e^{-cK}
\]
for some $c>0$. On the other hand, $a_n,b_m$ are algebraic integers,
and
$[\mathbb Q(a_n,b_m):\mathbb Q]\leq3^n4^m<4K^2$.
Every conjugate of $a_n$ has modulus $a_n$, and every conjugate of
$b_m$ has modulus $b_m$, so every conjugate of $F(a_n,b_m)$ has
modulus at most a fixed $C_F\geq2$. Whenever $F(a_n,b_m)\ne0$,
its nonzero integral norm consequently gives only
\[
 |F(a_n,b_m)|\geq C_F^{-4K^2}
 =e^{-O_F(K^2)}.
\]
These lower and upper bounds are compatible. Moreover, the required
nonvanishing has not been established by this comparison. Thus the
argument does not prove algebraic independence; it isolates a need
for additional simultaneous arithmetic separation.



\section{Dependency record for the five supplied declarations}
The input declarations were taken as assumptions with the meanings
specified in the request. The two zeta declarations supply RH; the
Dirichlet nonvanishing statement, the all-character critical-line
statement, and the primitive-character version supply the Dirichlet
GRH setting. The following mathematical dependency chain is all that
is used here:
\[
\begin{aligned}
 \pGRHD&\ \Longrightarrow\ \text{uniform prime-character sums}\\
       &\ \Longrightarrow\ \text{balanced finite prime pools}\\
       &\ \Longrightarrow\ \text{the proved finite-group lemmas}.
\end{aligned}
\]
No stronger GRH for general number fields, Artin $L$-functions, or
prime-tuple correlation conjecture is assumed. Section~\ref{p-sec:prime-graphs} additionally
uses the unconditional quadratic large sieve. The exact Lean definitions,
the proofs of the five declarations, and a machine formalization of this
paper are outside the present artifact.



All asymptotic statements above leave their numerical starting
moduli unspecified. The finite lemmas record explicit sufficient
inequalities, but no verified finite starting range is claimed.
The algorithms obtained by enumerating a unit group are finite
constructions, not polynomial-in-$\log q$ extraction or decoding
algorithms. The two-prime conjecture and the algebraic independence
question just discussed remain unresolved here.

\clearpage
\part{Sharp graph deletion and vertex sampling}\label{r-part-graph}

\section{Sharp deletion and sampling inequalities for Randi\'c indices}
\label{g-sec-main}

This part is independent of the Riemann-hypothesis assumptions used elsewhere
in this manuscript.  Its main results are sharp inequalities for deleting a
vertex or sampling a prescribed number of vertices from a finite graph.  The
constants and all equality cases are determined.  The application to the
Caporossi--Hansen conjecture is a reduction of a smallest counterexample to a
graph with a cut vertex.  A stronger deletion statement is formulated below
as a conjecture, and the unsuccessful gluing arguments are recorded as exact
counterexamples to the proposed arguments, not as counterexamples to
Caporossi--Hansen.  No assertion of historical priority is made for the
sampling inequalities.

\subsection{Conventions}
\label{g-sub-conventions}

All graphs are finite, simple, and undirected.  Unless connectivity is
explicitly required, isolated vertices are allowed.  For a graph $G$ and
$\alpha\in\mathbb R$, put
\[
 R_\alpha(G)=\sum_{uv\in E(G)}(d_G(u)d_G(v))^{-\alpha},
 \qquad R(G)=R_{1/2}(G).
\]
Only edges occur in the sum, so no convention involving a zero degree raised
to a negative power is needed.  The degrees in $R_\alpha(G[S])$ are the
degrees in the induced graph $G[S]$, not the degrees in $G$.  In particular,
$R_0(G)=|E(G)|$.  We write $G-v$ for $G[V(G)\setminus\{v\}]$.

For a connected graph of order $n\ge2$, define
\[
 W(G)=\sum_{\{x,y\}\subseteq V(G)}d_G(x,y),\qquad
 \mu(G)=\frac{2W(G)}{n(n-1)},\qquad
 \tau_G(v)=\sum_{x\ne v}d_G(v,x).
\]
Thus $W$ is the Wiener index, $\mu$ is the average distance over distinct
unordered pairs, and $\tau_G(v)$ is the transmission of $v$.  We also use
\[
 M(G)=\sum_{uv\in E(G)}\frac1{\max\{d_G(u),d_G(v)\}}.
\]
The symbols $K_n$, $P_n$, and $K_{1,n-1}$ denote the complete graph, path,
and star, respectively.

\subsection{The sharp one-vertex deletion inequality}
\label{g-sub-deletion}

\begin{theorem}[Sharp deletion inequality]
\label{g-thm-deletion}
Let $G$ have $n\ge3$ vertices and let $\alpha\ge0$.  Then
\begin{equation}
 \sum_{v\in V(G)}R_\alpha(G-v)
 \le C_n(\alpha)R_\alpha(G),
 \label{g-eq-deletion}
\end{equation}
where
\begin{equation}
 C_n(\alpha)=
 \begin{cases}
 (n-1)^{2\alpha}(n-2)^{1-2\alpha},&0\le\alpha\le\tfrac12,\\[2pt]
 n-3+2^{2\alpha},&\alpha\ge\tfrac12.
 \end{cases}
 \label{g-eq-deletion-constant}
\end{equation}
The constant is best possible for every $n$ and $\alpha$.  For a graph with
at least one edge, equality holds precisely in the following cases:
\begin{enumerate}
 \item $\alpha=0$: every graph;
 \item $0<\alpha<\tfrac12$: $G=K_n$;
 \item $\alpha=\tfrac12$: every nontrivial component is a clique of order
 at least three;
 \item $\alpha>\tfrac12$: every nontrivial component is a triangle.
\end{enumerate}
Isolated vertices are permitted in the last two cases.  Edgeless graphs
satisfy equality for every $\alpha$.
\end{theorem}

\begin{proof}
Fix an edge $uw$, and write $a=d_G(u)$, $b=d_G(w)$.  For a deleted vertex
$z\notin\{u,w\}$, the contribution of this edge to $R_\alpha(G-z)$ is
$(ab)^{-\alpha}A_zB_z$, where
\[
 A_z=
 \begin{cases}
 (a/(a-1))^\alpha,&z\sim u,\\
 1,&z\not\sim u,
 \end{cases}
\]
and $B_z$ is defined in the same way with $w,b$ in place of $u,a$.
If $a=1$, no eligible $z$ is adjacent to $u$, and $A_z=1$ for all eligible
$z$; thus the displayed quotient is never evaluated at $a=1$.

Set $q=2\alpha$ and, for positive integers $j$, let
\[
 f_q(j)=j\bigl((1+1/j)^q-1\bigr),\qquad f_q(0)=0.
\]
There are $a-1$ eligible neighbors of $u$, and therefore
\[
 \sum_{z\ne u,w}A_z^2=n-2+f_q(a-1).
\]
Cauchy--Schwarz gives
\begin{equation}
 \sum_{z\ne u,w}A_zB_z
 \le\sqrt{\bigl(n-2+f_q(a-1)\bigr)
              \bigl(n-2+f_q(b-1)\bigr)}.
 \label{g-eq-edge-cauchy}
\end{equation}
For $q>0$, the identity
\[
 f_q(j)=q\int_0^1(1+t/j)^{q-1}\,dt
\]
shows that $f_q(j)$ is strictly increasing in $j\ge1$ when $0<q<1$,
strictly decreasing when $q>1$, and equal to one when $q=1$.
For $q=0$, it is zero.  Consequently, if $0<q<1$, the right side of
\eqref{g-eq-edge-cauchy} is at most
\[
 n-2+f_q(n-2)=(n-1)^q(n-2)^{1-q}.
\]
If $q>1$, its upper bound is instead
\[
 n-2+f_q(1)=n-3+2^q.
\]
At $q=1$ the bound is $n-1$, and at $q=0$ it is $n-2$.
Multiplying by $(ab)^{-\alpha}$ and summing over edges proves
\eqref{g-eq-deletion}.

We next check equality.  Every original edge has positive weight, so equality
in the sum requires equality in each of the edgewise bounds.  For $0<q<1$,
strict monotonicity forces $a=b=n-1$ on every edge.  The presence of one
such edge forces every vertex to be incident with an edge, and then every
vertex is universal.  Hence $G=K_n$.

For $q>1$, both endpoint degrees must be two.  The two vectors in
\eqref{g-eq-edge-cauchy} have the same squared norm, so equality in
Cauchy--Schwarz requires $A_z=B_z$ for every eligible $z$.  Each endpoint
has a unique remaining neighbor, and these neighbors must coincide.  Each
edge therefore belongs to a triangle whose vertices have degree two; all
nontrivial components are triangles.

For $q=1$, both endpoint degrees must be at least two.  The squared norms
again agree, and equality requires $A=B$.  Their exceptional entries have
values $\sqrt{a/(a-1)}$ and $\sqrt{b/(b-1)}$.  These values are strictly
decreasing functions of the degree.  It follows that $a=b$ and
\begin{equation}
 N_G(u)\setminus\{w\}=N_G(w)\setminus\{u\}.
 \label{g-eq-true-twins}
\end{equation}
In a connected component satisfying \eqref{g-eq-true-twins} on every edge,
the endpoints of every path of length two are adjacent.  A shortest path
between any two vertices thus has length at most one.  The component is a
clique, and its order is at least three because its degrees are at least
two.  Conversely, all the stated graphs give equality.

Finally,
\[
 R_\alpha(K_n)=\frac n2(n-1)^{1-2\alpha},
\]
so $K_n$ gives exactly the first ratio in
\eqref{g-eq-deletion-constant}.  A triangle together with $n-3$ isolated
vertices has $R_\alpha=3\cdot4^{-\alpha}$; deleting an isolated vertex
preserves that value, and deleting a triangle vertex leaves a single edge
of weight one.  Its ratio is $n-3+2^{2\alpha}$.  This proves sharpness.
\end{proof}

\begin{corollary}[The ordinary index]
\label{g-cor-ordinary-deletion}
For every graph of order $n\ge3$,
\begin{equation}
 \sum_vR(G-v)\le(n-1)R(G),\qquad
 \sum_v\bigl(R(G)-R(G-v)\bigr)\ge R(G).
 \label{g-eq-ordinary-deletion}
\end{equation}
For a connected graph of order at least three, equality holds exactly for
a complete graph.  A component $K_2$ makes the inequality strict.
\end{corollary}

\begin{remark}
The last qualification matters: for $G=K_2$ the left side of the first
inequality in \eqref{g-eq-ordinary-deletion} is zero, whereas
$(n-1)R(G)=1$.  For general graphs the equality components are cliques of
order at least three, not arbitrary nontrivial cliques.
\end{remark}

\begin{proposition}[Negative exponents]
\label{g-prop-negative-deletion}
If $\alpha<0$, then
\[
 \sum_vR_\alpha(G-v)\le(n-2)R_\alpha(G).
\]
The coefficient is sharp.  A nonempty graph gives equality exactly when
its nontrivial components are single edges.
\end{proposition}

\begin{proof}
The edge weight is now a positive power of the product of its endpoint
degrees.  Deleting a third vertex can only decrease it, and an edge survives
exactly $n-2$ deletions.  If an endpoint has another neighbor, deleting that
neighbor strictly decreases the surviving edge weight.  Thus equality on
every edge is equivalent to every endpoint having degree one.  A matching
with isolated vertices realizes the coefficient.
\end{proof}

\subsection{The exact envelope for uniform induced subgraphs}
\label{g-sub-sampling}

Let $2\le k<n$ and let $S$ be uniform among all $k$-element subsets of
$V(G)$.  For $1\le d\le n-1$, define a hypergeometric variable $H_d$ by
\begin{equation}
 \Pr(H_d=h)=
 \frac{\binom{d-1}{h}\binom{n-1-d}{k-2-h}}
      {\binom{n-2}{k-2}}.
 \label{g-eq-hypergeom}
\end{equation}
As usual, a binomial coefficient is zero when its lower argument lies
outside the range from zero to its nonnegative upper argument.  Thus
$H_d$ counts marked elements when $k-2$ elements are sampled from a
population of $n-2$ containing $d-1$ marked elements.  Put
\begin{equation}
 F_{n,k,q}(d)=d^q\mathbb E(1+H_d)^{-q},\qquad
 \mathcal D_{n,k,q}=\operatorname*{arg\,max}_{1\le d\le n-1}
                    F_{n,k,q}(d).
 \label{g-eq-sampling-F}
\end{equation}

\begin{theorem}[Sharp uniform-sampling inequality]
\label{g-thm-sampling}
For every $\alpha\ge0$ and every graph $G$ of order $n$,
\begin{equation}
 \mathbb E R_\alpha(G[S])\le
 \frac{k(k-1)}{n(n-1)}
 \max_{1\le d\le n-1}F_{n,k,2\alpha}(d)\,R_\alpha(G).
 \label{g-eq-sampling-main}
\end{equation}
The coefficient is sharp.  For each maximizing degree $d$, the graph
$K_{d+1}$ together with $n-d-1$ isolated vertices attains equality.

If $\alpha>0$ and $G$ has at least one edge, equality holds if and only if
every nontrivial component is a clique whose degree belongs to
$\mathcal D_{n,k,2\alpha}$.  Isolated vertices are allowed.  At $\alpha=0$,
every graph gives equality.
\end{theorem}

\begin{proof}
An edge $uv$ is retained with probability
\[
 p_{n,k}=\frac{k(k-1)}{n(n-1)}.
\]
Condition on retaining its endpoints.  The remaining selected vertices
form a uniform $(k-2)$-subset $T$ of the other $n-2$ vertices.  Write
$a=d_G(u)$ and $b=d_G(v)$, and set
\[
 A_T=\left(\frac a{1+|T\cap(N_G(u)\setminus\{v\})|}\right)^\alpha,
 \quad
 B_T=\left(\frac b{1+|T\cap(N_G(v)\setminus\{u\})|}\right)^\alpha.
\]
The retained edge has weight $(ab)^{-\alpha}A_TB_T$.  The two squared
moments are $F_{n,k,2\alpha}(a)$ and $F_{n,k,2\alpha}(b)$, so
\begin{equation}
 \mathbb E(A_TB_T)
 \le\sqrt{F_{n,k,2\alpha}(a)F_{n,k,2\alpha}(b)}
 \le\max_dF_{n,k,2\alpha}(d).
 \label{g-eq-sampling-cauchy}
\end{equation}
Multiplication by the original edge weight and by $p_{n,k}$, followed by
summation, proves \eqref{g-eq-sampling-main}.

For a clique component, both conditional endpoint degrees are identical
for every $T$.  The first inequality in \eqref{g-eq-sampling-cauchy} is
therefore an equality.  If its original degree is maximizing, the second
inequality is also exact.  This proves sharpness and the sufficiency of
the stated equality condition.

For necessity assume $\alpha>0$ and equality in the graph inequality.
Every edge must attain both bounds in
\eqref{g-eq-sampling-cauchy}.  In particular, its endpoints have
$F(a)=F(b)=\max F$.  Equality in Cauchy--Schwarz, together with these equal
squared norms, implies $A_T=B_T$ for every conditional sample $T$.
Since $\alpha>0$, this is equivalent to
\begin{equation}
 \frac a{1+|T\cap(N_G(u)\setminus\{v\})|}
 =
 \frac b{1+|T\cap(N_G(v)\setminus\{u\})|}.
 \label{g-eq-sample-ratio-equality}
\end{equation}
Taking reciprocals and expectations is useful here: for any degree $d$,
\begin{equation}
 \mathbb E\frac{1+H_d}{d}
 =\frac{k-2}{n-2}
  +\left(1-\frac{k-2}{n-2}\right)\frac1d.
 \label{g-eq-reciprocal-mean}
\end{equation}
Because $k<n$, the coefficient of $1/d$ is positive.  Equality in
\eqref{g-eq-sample-ratio-equality} consequently forces $a=b$.

Suppose first that $3\le k<n$.  The two remaining neighborhoods then
have equal intersection counts with every $(k-2)$-subset of a set of
size $n-2$.  If two subsets $A,B$ have that property for sample size
$s$ with $1\le s<N$, the quantities
$\mathbf1_A(z)-\mathbf1_B(z)$ are equal for all population elements:
exchange $z$ and $w$ in a common sample of $s-1$ other elements.
Here $|A|=|B|=a-1$, so the common value is zero.  Hence $A=B$, giving
\eqref{g-eq-true-twins}.  The component is a clique by the argument in
Theorem~\ref{g-thm-deletion}, and its degree must maximize $F$.

If $k=2$, the conditional sample is empty and $F_{n,2,q}(d)=d^q$.
For $q>0$ the unique maximizing degree is $n-1$.  Every endpoint of every
edge is universal, and the nonempty graph is $K_n$.  This agrees with
the stated classification.  Finally, at $\alpha=0$ the assertion is just
the exact expected edge count $\mathbb E|E(G[S])|=p_{n,k}|E(G)|$.
\end{proof}

\begin{remark}[Boundary cases]
The restriction $k<n$ is essential to the equality classification.
At $k=n$, the sample is $G$ itself and every graph gives equality.  If
$k=0$ or $1$, the sampled index is zero.  An edgeless graph gives equality
for every admissible parameter.  No additional regular-graph equality
case appears at $k=2$.
\end{remark}

\subsubsection{A closed form up to the ordinary Randi\'c exponent}

\begin{lemma}[The reciprocal hypergeometric moment]
\label{g-lem-hypergeom-one}
For $1\le d\le n-1$,
\begin{equation}
 F_{n,k,1}(d)=\frac{n-1}{k-1}
 \left(1-
 \frac{\binom{n-1-d}{k-1}}{\binom{n-1}{k-1}}\right).
 \label{g-eq-hypergeom-one}
\end{equation}
Its maximizing degrees are exactly $d\ge n-k+1$.
\end{lemma}

\begin{proof}
The identity
\[
 \frac d{h+1}\binom{d-1}{h}=\binom d{h+1}
\]
and Vandermonde's identity give
\begin{align*}
 F_{n,k,1}(d)
 &=\frac{\sum_{j\ge1}\binom d j
                   \binom{n-1-d}{k-1-j}}
          {\binom{n-2}{k-2}}\\
 &=\frac{\binom{n-1}{k-1}-\binom{n-1-d}{k-1}}
          {\binom{n-2}{k-2}}.
\end{align*}
This is \eqref{g-eq-hypergeom-one}.  Its upper bound $(n-1)/(k-1)$
is attained precisely when $n-1-d<k-1$, or $d\ge n-k+1$.
\end{proof}

\begin{corollary}[Closed sampling constant]
\label{g-cor-closed-sampling}
For $0\le\alpha\le\tfrac12$,
\begin{equation}
 \mathbb E R_\alpha(G[S])
 \le \frac{k}{n}
       \left(\frac{k-1}{n-1}\right)^{1-2\alpha}R_\alpha(G).
 \label{g-eq-closed-sampling}
\end{equation}
For $0<\alpha<\tfrac12$, the only nonempty equality graph is $K_n$.
For $\alpha=\tfrac12$, the nonempty equality graphs are precisely the
disjoint unions of cliques of order at least $n-k+2$, with isolated
vertices allowed.  At $\alpha=0$ every graph gives equality.
\end{corollary}

\begin{proof}
Put $q=2\alpha$.  If $0<q<1$, Jensen's inequality for the strictly
concave function $x^q$ yields
\[
 F_{n,k,q}(d)
 =\mathbb E\left(\frac d{1+H_d}\right)^q
 \le F_{n,k,1}(d)^q
 \le\left(\frac{n-1}{k-1}\right)^q.
\]
Substitution in Theorem~\ref{g-thm-sampling} proves the coefficient.
For $3\le k<n$ and $2\le d\le n-2$, the hypergeometric variable $H_d$
is nonconstant: one can exchange a marked and an unmarked element in
a sample of size between one and $n-3$.  Jensen is therefore strict.
The degree $d=1$ does not maximize the last bound, leaving only $d=n-1$.
For $k=2$, uniqueness follows from $F_{n,2,q}(d)=d^q$.  The equality
classification in Theorem~\ref{g-thm-sampling} now gives $K_n$.

At $q=1$, Lemma~\ref{g-lem-hypergeom-one} gives the exact maximizing
degrees.  A clique with degree at least $n-k+1$ has order at least
$n-k+2$, proving the critical equality statement.  The case $q=0$
again follows by counting edges.
\end{proof}

\begin{corollary}[Normalized expectations and random deletion]
\label{g-cor-normalized}
For a graph $H$ of order $m\ge2$, define
\[
 \mathcal R_\alpha(H)=
 \frac{R_\alpha(H)}{R_\alpha(K_m)}
 =\frac{2R_\alpha(H)}{m(m-1)^{1-2\alpha}}.
\]
If $0\le\alpha\le\tfrac12$, then
\begin{equation}
 \mathbb E\mathcal R_\alpha(G[S])\le\mathcal R_\alpha(G).
 \label{g-eq-normalized-expectation}
\end{equation}
Moreover, start with $G$ and repeatedly delete a uniformly chosen
remaining vertex until two vertices remain.  The normalized index of the
remaining induced graph is a supermartingale for this deletion process.
In particular, at $\alpha=\tfrac12$ the normalization is $2R(H)/|V(H)|$.
\end{corollary}

\begin{proof}
Divide \eqref{g-eq-closed-sampling} by
$R_\alpha(K_k)=k(k-1)^{1-2\alpha}/2$.  For the process statement,
condition at each step on the entire remaining graph and apply the
one-step instance.  The inequality holds for every induced graph,
including those with isolated vertices, so the conditional inequality
is valid throughout the process.
\end{proof}

\begin{proposition}[Sampling at negative exponents]
\label{g-prop-negative-sampling}
For $\alpha<0$ and $2\le k<n$,
\[
 \mathbb E R_\alpha(G[S])\le
 \frac{k(k-1)}{n(n-1)}R_\alpha(G).
\]
The coefficient is sharp, with nonempty equality precisely for matchings
together with isolated vertices.
\end{proposition}

\begin{proof}
Condition on an edge being retained.  Its endpoint degrees can only
decrease, so its positive-power weight cannot increase.  If an endpoint
has another neighbor, that neighbor has positive probability of being
omitted because $k<n$, causing a strict decrease with positive
probability.  Equality is therefore equivalent to every edge having
both endpoint degrees one.  A matching attains the coefficient.
\end{proof}

\subsection{A reduction for the Caporossi--Hansen conjecture}
\label{g-sub-CH}

For $n\ge2$, put
\begin{equation}
 B_n=\sqrt{n-1}-2+\frac2n,\qquad
 \mathcal S(G)=R(G)-\mu(G)-B_n.
 \label{g-eq-CH-surplus}
\end{equation}
The Caporossi--Hansen additive conjecture asks whether
\begin{equation}
 \mathcal S(G)\ge0
 \quad\text{for every connected graph of order }n\ge2.
 \label{g-eq-CH}
\end{equation}
The star attains equality.  The conjecture originates with
Caporossi and Hansen~\cite{g-CH2000}; its explicit formulation is also
recorded as Conjecture~4.8 in~\cite{g-FMS2003}.  The argument below proves
a reduction, not \eqref{g-eq-CH} for all graphs.

\begin{lemma}[The standard star lower bound]
\label{g-lem-star-lower}
Every connected graph of order $n\ge2$ satisfies
\begin{equation}
 R(G)\ge\sqrt{n-1}.
 \label{g-eq-star-lower}
\end{equation}
\end{lemma}

\begin{proof}
We include a self-contained spectral proof of this standard bound
\cite{g-BE1998}.  Let $m=|E(G)|$, let $A$ be the adjacency matrix, and
let $\rho$ be its largest eigenvalue.  Cauchy--Schwarz gives
\[
 R(G)\ge
 \frac{m^2}{\sum_{uv\in E(G)}\sqrt{d_G(u)d_G(v)}}.
\]
Apply the Rayleigh quotient to the vector with entries $\sqrt{d_G(v)}$.
Its squared norm is $2m$, and hence
\[
 \sum_{uv\in E(G)}\sqrt{d_G(u)d_G(v)}\le m\rho,
 \qquad R(G)\ge m/\rho.
\]
For any vertex $u$, all $n-d_G(u)-1$ nonneighbors have degree at least
one.  Therefore
\[
 \sum_{v\sim u}d_G(v)\le2m-d_G(u)-(n-d_G(u)-1)=2m-n+1.
\]
These are the row sums of $A^2$.  The maximum row-sum bound for a
nonnegative matrix gives $\rho^2\le2m-n+1$.  Finally,
\[
 m^2-(n-1)(2m-n+1)=(m-n+1)^2\ge0.
\]
Since $m\ge n-1$, all quantities used above are positive and
$R(G)\ge m/\sqrt{2m-n+1}\ge\sqrt{n-1}$.
\end{proof}

\begin{theorem}[A smallest counterexample has a cut vertex]
\label{g-thm-cut-reduction}
If \eqref{g-eq-CH} has a counterexample, every counterexample of smallest
order has a cut vertex.
\end{theorem}

\begin{proof}
Let $G$ be a counterexample of smallest order $n$.  The graph $K_2$
satisfies equality, so $n\ge3$.  Suppose that $G$ has no cut vertex.
Every $G-v$ is then connected and satisfies the conjecture by minimality.
Deleting $v$ cannot decrease distances between surviving vertices, so
\[
 W(G-v)\ge W(G)-\tau_G(v).
\]
Writing $\mu=\mu(G)$, this becomes
\begin{equation}
 \mu(G-v)\ge
 \frac{n\mu-2\tau_G(v)/(n-1)}{n-2}.
 \label{g-eq-distance-deletion}
\end{equation}
Set $\Delta_v=R(G)-R(G-v)$.  The strict failure for $G$ and the validity
for $G-v$ imply
\begin{align}
 \Delta_v
 &< B_n-B_{n-1}+\mu-\mu(G-v)\notag\\
 &\le B_n-B_{n-1}
 +\frac2{n-2}\left(\frac{\tau_G(v)}{n-1}-\mu\right).
 \label{g-eq-counterexample-decrement}
\end{align}
Now $\sum_v\tau_G(v)=n(n-1)\mu$, so summing cancels the transmission
terms and yields
\[
 \sum_v\Delta_v<n(B_n-B_{n-1}).
\]
For $n\ge3$,
\[
 n(B_n-B_{n-1})
 =\frac n{\sqrt{n-1}+\sqrt{n-2}}-\frac2{n-1}
 <\sqrt{n-1}.
\]
Indeed, $(n-1)+\sqrt{(n-1)(n-2)}>n$.  But
Corollary~\ref{g-cor-ordinary-deletion} and
Lemma~\ref{g-lem-star-lower} give
\[
 R(G)\le\sum_v\Delta_v
 <\sqrt{n-1}\le R(G),
\]
a contradiction.
\end{proof}

\begin{remark}[The scope of the reduction]
The theorem does not prove the conjecture on the class of graphs without
cut vertices by induction within that class: deleting a vertex can create
cut vertices.  It proves that a globally smallest counterexample cannot
belong to that class.  Eliminating its possible cut vertices is a separate
problem.
\end{remark}

\subsection{Exact gluing identities and two failed envelopes}
\label{g-sub-gluing}

\begin{proposition}[Gluing at one vertex]
\label{g-prop-gluing}
Let connected graphs $G_1,\ldots,G_t$, $t\ge2$, meet only at a common
root $v$, with no edges between their other vertices, and let their union
be $G$.  Suppose each $G_i$ has at least two vertices.  Write
\[
 n_i=|V(G_i)|,\quad n=1+\sum_i(n_i-1),\quad
 d_i=d_{G_i}(v),\quad d=\sum_i d_i,
\]
\[
 A_i=\sum_{u\sim_{G_i}v}d_{G_i}(u)^{-1/2},\qquad
 T_i=\sum_{x\in V(G_i)\setminus\{v\}}d_{G_i}(v,x).
\]
Then
\begin{align}
 R(G)&=\sum_iR(G_i)-L,
 &L&=\sum_iA_i\left(\frac1{\sqrt{d_i}}-\frac1{\sqrt d}\right),
 \label{g-eq-gluing-R}\\
 W(G)&=\sum_iW(G_i)+\sum_i(n-n_i)T_i.
 \label{g-eq-gluing-W}
\end{align}
Moreover,
\begin{equation}
 0\le L\le\sum_i\sqrt{d_i}-\sqrt d.
 \label{g-eq-gluing-loss}
\end{equation}
Equality in the upper bound holds precisely when every piece is a star
rooted at its center.
\end{proposition}

\begin{proof}
Every nonroot degree is unchanged by gluing, while the root degree changes
from $d_i$ to $d$ in the $i$th piece.  This proves
\eqref{g-eq-gluing-R}.  If $x\in G_i-v$ and $y\in G_j-v$ with $i\ne j$,
their distance is $d_{G_i}(x,v)+d_{G_j}(v,y)$.  Thus the cross contribution
of two pieces is $(n_j-1)T_i+(n_i-1)T_j$.  Summing over unordered pairs
of pieces gives \eqref{g-eq-gluing-W}.

All coefficients in $L$ are positive and $A_i\le d_i$, since every
nonroot degree is at least one.  This proves \eqref{g-eq-gluing-loss}.
Equality requires every neighbor of the root in every piece to have
degree one.  Connectivity then forces that piece to be a star centered
at the root.  Such stars do give equality.
\end{proof}

For two pieces, replacing each $R(G_i)$ by only its conjectural lower
bound $\mu(G_i)+B_{n_i}$ gives the gluing envelope
\begin{equation}
 \mathcal E=
 B_{n_1}+B_{n_2}-B_n+\mu(G_1)+\mu(G_2)-\mu(G)-L.
 \label{g-eq-plain-envelope}
\end{equation}
If $\mathcal E\ge0$, those lower bounds suffice.  The next example shows
that $\mathcal E$ need not be nonnegative, even though all the actual
graphs satisfy the conjecture.

\begin{example}[Two four-vertex paths]
\label{g-ex-P4-gluing}
Glue two copies of $P_4$ at endpoints.  The union is $P_7$ and
\[
 n_1=n_2=4,\quad d_1=d_2=1,\quad A_1=A_2=1/\sqrt2,
 \quad T_1=T_2=6.
\]
Consequently $L=\sqrt2-1$, $\mu(G_i)=5/3$, and $\mu(G)=8/3$.
The envelope \eqref{g-eq-plain-envelope} is
\begin{equation}
 \mathcal E=2\sqrt3-\sqrt2-\sqrt6+\frac8{21}<0.
 \label{g-eq-P4-negative}
\end{equation}
Its value is approximately $-0.0186493$.  The strict sign follows, for
example, by squaring the rational bounds
$\sqrt3<1.732051$, $\sqrt2>1.414213$, and
$\sqrt6>2.449489$.  This refutes the nonnegativity of the envelope, not
the Caporossi--Hansen conjecture.
\end{example}

To formulate a stronger unsuccessful envelope, it is useful to record
two additional lower bounds with their dependencies.

\begin{lemma}[A degree bridge and a diameter consequence]
\label{g-lem-degree-bridge}
If $G$ is connected of order $n\ge2$ and has maximum degree $\Delta$, then
\begin{equation}
 R(G)\ge
 \frac{n+(\sqrt\Delta-1)M(G)}{\sqrt\Delta+1}.
 \label{g-eq-degree-bridge}
\end{equation}
Using the established theorem $M(G)\ge D(G)/2$ of
Dvo\v r\'ak, Lidick\'y, and \v Skrekovski~\cite{g-DLS2011}, this implies
\begin{equation}
 R(G)\ge\frac{D(G)}2+\sqrt{n-D(G)+1}-1.
 \label{g-eq-diameter-bridge}
\end{equation}
\end{lemma}

\begin{proof}
Put $r=\sqrt\Delta$.  On an edge with degrees $a\le b$, set
$x=\sqrt{b/a}$, so $1\le x\le r$.  The inequality
$(x-1)(r-x)\ge0$ is equivalent to
\[
 \frac{r+1}{\sqrt{ab}}\ge\frac1a+\frac r b.
\]
Summing over edges gives $(r+1)R\ge n+(r-1)M$, because
$\sum_{uv\in E}(1/d_u+1/d_v)=n$ for a graph without isolated vertices.
This proves \eqref{g-eq-degree-bridge}.

Let $D=D(G)$ and $e=n-D$.  On a diametral geodesic, a vertex has at most
two neighbors on the geodesic; an outside vertex has neighbors in at most
three consecutive positions, since otherwise it provides a shortcut.
Counting outside vertices in either case gives $\Delta\le e+1$.
The cited bound $M\ge D/2$ and \eqref{g-eq-degree-bridge} now give
\[
 R-\frac D2\ge\frac{e}{1+\sqrt\Delta}
 \ge\frac e{1+\sqrt{e+1}}=\sqrt{e+1}-1.
\]
Only the input $M\ge D/2$ is external to this proof; the weighted
distance theorem from the corpus is not needed here.
\end{proof}

\begin{example}[Two small brooms defeat a strengthened envelope]
\label{g-ex-brooms}
Let $H_j$ consist of a hub with two pendant leaves and a path of length
$j$ attached at the hub, rooted at the other end of that path, for
$j=2,3$.  Direct edge-cut and degree calculations give
\[
\begin{array}{c|ccccccc}
 & n & T & W & M & \Delta & D & R\\ \hline
 H_2&5&9&18&3/2&3&3&1/\sqrt2+1/\sqrt6+2/\sqrt3\\
 H_3&6&14&32&2&3&4&1/\sqrt2+1/\sqrt6+2/\sqrt3+1/2
\end{array}
\]
For clarity, the Wiener values follow from summing $s(n-s)$ over tree
edges: for $H_2$ this is $4+4+6+4=18$, and for $H_3$ it is
$5+5+9+8+5=32$.

Consider the lower envelope
\[
 \mathcal L(H)=\max\left\{
 \mu(H)+B_{|H|},\ \sqrt{|H|-1},\
 \frac{|H|+(\sqrt\Delta-1)M(H)}{\sqrt\Delta+1},\
 \frac D2+\sqrt{|H|-D+1}-1\right\}.
\]
The first term is available under induction; the other three are the
proved lower bounds above.  Substitution gives
\[
 \mathcal L(H_2)=\sqrt3+\frac12,
 \qquad \mathcal L(H_3)=\sqrt3+1.
\]
Indeed, in both cases the last two terms agree with these values and
dominate the first two terms.

Glue the roots.  Both roots originally have degree one and have a
degree-two neighbor, so $L=\sqrt2-1$.  The resulting graph $G$ has
\[
 n=10,\qquad W(G)=18+32+5\cdot9+4\cdot14=151,
 \qquad \mu(G)=151/45.
\]
Even this strengthened envelope leaves the negative remainder
\begin{equation}
 \mathcal L(H_2)+\mathcal L(H_3)-L-\mu(G)-B_{10}
 =2\sqrt3-\sqrt2-\frac{37}{18}<0.
 \label{g-eq-broom-negative}
\end{equation}
Its value is approximately $-0.00566750$.  The square-checked bounds
$\sqrt3<1.732051$ and $\sqrt2>1.414213$ already prove strict negativity.
The actual surplus is positive:
\[
 \mathcal S(G)=\frac2{\sqrt6}+\frac4{\sqrt3}-\frac{55}{18}
 =\frac{6\sqrt6+24\sqrt3-55}{18}>0.
\]
For example, $\sqrt6>2.44$ and $\sqrt3>1.73$ prove the final sign.
The example identifies information lost by the lower envelope; it does
not refute any of its component lower bounds.
\end{example}

\subsubsection{A rooted consequence of the external weighted theorem}
\label{g-sub-rooted-external}

The preceding proved deletion and sampling results do not use the weighted
distance theorem in the external corpus.  For completeness, the following
derived estimate records exactly what one additional rooted specialization
would supply.  Its dependency is explicit.

\begin{proposition}[Conditional on the weighted distance inequality]
\label{g-prop-rooted-weighted}
Assume the following weighted inequality for a connected graph $H$:
\begin{equation}
 \sum_{x\ne y}p_xp_y d_H(x,y)+2\sum_xp_x^2\le M(H)+1
 \quad\text{for every probability vector }p.
 \label{g-eq-external-WD}
\end{equation}
This is the theorem stated and audited in the external manuscript
\cite{g-VMRC227}; its full proof is not reproduced here.
Root $H$ at $v$, put $m=|V(H)|-1$, $T=\tau_H(v)$, $W=W(H)$,
and $M=M(H)$.  If
\[
 T>2m,\qquad
 \frac Tm\ge\frac{2(W-T)}{m^2}+\frac2m,
\]
then
\begin{equation}
 (T-2m)^2\le
 2(M-1)\bigl((m+1)T-W-m(m+1)\bigr).
 \label{g-eq-rooted-quadratic}
\end{equation}
\end{proposition}

\begin{proof}
Give the root mass $1-t$ and each other vertex mass $t/m$, where
$0\le t\le1$.  The left side of \eqref{g-eq-external-WD} is
\[
 2+2at+2bt^2,
 \quad
 a=T/m-2,
 \quad
 b=1+1/m-T/m+(W-T)/m^2.
\]
The first assumption gives $a>0$.  The second is exactly $2b+a\le0$,
so $b<0$ and the maximizing parameter $t_*=-a/(2b)$ lies in $(0,1]$.
Applying \eqref{g-eq-external-WD} at $t_*$ yields
$a^2\le-2b(M-1)$.  Multiplication by $m^2$ gives
\eqref{g-eq-rooted-quadratic}.
\end{proof}

\subsection{The compensated deletion candidate}
\label{g-sub-CD}

\begin{conjecture}[Compensated noncut deletion]
\label{g-conj-CD}
Every connected graph $G$ of order $n\ge3$ has a noncut vertex $v$ such
that
\begin{equation}
 nR(G)-(n-2)R(G-v)-\frac{2\tau_G(v)}{n-1}
 \ge C_n^*,\qquad
 C_n^*=nB_n-(n-2)B_{n-1}.
 \label{g-eq-CD}
\end{equation}
\end{conjecture}

The notation $C_n^*$ is separate from the proved coefficient
$C_n(\alpha)$ in Theorem~\ref{g-thm-deletion}.  A useful exact expression
for this candidate threshold is
\begin{equation}
 C_n^*=n\sqrt{n-1}-(n-2)\sqrt{n-2}-4+\frac2{n-1}.
 \label{g-eq-CD-threshold}
\end{equation}

\begin{proposition}[The candidate would imply Caporossi--Hansen]
\label{g-prop-CD-implies}
If Conjecture~\ref{g-conj-CD} holds for every connected graph, then
\eqref{g-eq-CH} holds for every connected graph.
\end{proposition}

\begin{proof}
For any noncut vertex $v$, \eqref{g-eq-distance-deletion} implies
\begin{align*}
 n\mathcal S(G)-(n-2)\mathcal S(G-v)
 &=nR(G)-(n-2)R(G-v)-n\mu(G)+(n-2)\mu(G-v)-C_n^*\\
 &\ge nR(G)-(n-2)R(G-v)-\frac{2\tau_G(v)}{n-1}-C_n^*.
\end{align*}
A vertex supplied by \eqref{g-eq-CD} therefore gives
$n\mathcal S(G)\ge(n-2)\mathcal S(G-v)$.  Induction from $K_2$, whose
surplus is zero, proves the assertion.
\end{proof}

\begin{lemma}[Exact leaf decrement]
\label{g-lem-leaf-decrement}
Let $v$ be a leaf of a connected graph of order at least three, let $u$
be its neighbor, and set $d=d_G(u)\ge2$.  Then
\begin{align}
 R(G)-R(G-v)
 &=\frac1{\sqrt d}
 -\left(\frac1{\sqrt{d-1}}-\frac1{\sqrt d}\right)
     \sum_{w\in N_G(u)\setminus\{v\}}\frac1{\sqrt{d_G(w)}}
 \label{g-eq-leaf-exact}\\
 &=\sqrt d-\sqrt{d-1}
 +\left(\frac1{\sqrt{d-1}}-\frac1{\sqrt d}\right)
    \sum_{w\in N_G(u)\setminus\{v\}}
       \left(1-\frac1{\sqrt{d_G(w)}}\right).
 \label{g-eq-leaf-positive}
\end{align}
In particular, the decrement is at least $\sqrt d-\sqrt{d-1}$, with
equality precisely when $G$ is a star centered at $u$.  Also,
\[
 \tau_G(v)=\tau_G(u)+n-2.
\]
\end{lemma}

\begin{proof}
The deleted edge has weight $1/\sqrt d$.  The only surviving weights
that change are the other edges incident with $u$, where its degree drops
from $d$ to $d-1$.  This gives \eqref{g-eq-leaf-exact}; adding and
subtracting $d-1$ terms gives \eqref{g-eq-leaf-positive}.  Each summand
in its last sum is nonnegative.  Equality forces all the other neighbors
of $u$ to be leaves, and connectivity then forces a star.  For distances,
every path from $v$ to a vertex other than $v$ begins with $vu$; summing
these distances gives the transmission identity.
\end{proof}

\subsubsection{Why uncompensated deletion fails}

The superficially stronger assertion that some noncut vertex always
satisfies $\mathcal S(G)\ge\mathcal S(G-v)$ is false.
For $n\ge3$,
\[
 R(P_n)=\frac{n-3}{2}+\sqrt2,\qquad
 \mu(P_n)=\frac{n+1}{3},\qquad
 \mathcal S(P_n)=\frac n6+\sqrt2+\frac16-\sqrt{n-1}-\frac2n.
\]
The only noncut vertices of a path are its endpoints, and
\[
 \mathcal S(P_5)-\mathcal S(P_4)=\sqrt3-\frac{26}{15}<0,
 \qquad
 \mathcal S(P_6)-\mathcal S(P_5)=\frac{67}{30}-\sqrt5<0.
\]
Both signs follow immediately by squaring positive rational bounds.

The failure is not confined to small graphs.  Let $T_m$ have two hubs
joined by a path of $2m$ edges and attach $m$ leaves to each hub.  Then
$|T_m|=4m+1$, and every noncut vertex is a leaf.  Direct calculation gives
\begin{equation}
 R(T_m)=m-1+\frac{2m+\sqrt2}{\sqrt{m+1}},\qquad
 \mu(T_m)=\frac{11m^2+18m+1}{3(4m+1)}.
 \label{g-eq-doublebroom-values}
\end{equation}
For a leaf $v$,
\[
 \frac{\tau_{T_m}(v)}{4m}=m+\frac74-\frac1{4m}.
\]
Deleting a leaf preserves all surviving distances, so
\[
 \mu(T_m)-\mu(T_m-v)
 =\frac2{4m-1}
  \left(\frac{\tau_{T_m}(v)}{4m}-\mu(T_m)\right)
 \longrightarrow\frac1{24}.
\]
Meanwhile \eqref{g-eq-leaf-exact} gives
$R(T_m)-R(T_m-v)=O(m^{-1/2})$, and
$B_{4m+1}-B_{4m}=O(m^{-1/2})$.  Thus
\begin{equation}
 \mathcal S(T_m)-\mathcal S(T_m-v)\longrightarrow-\frac1{24}.
 \label{g-eq-doublebroom-failure}
\end{equation}
Every sufficiently large graph in this family defeats uncompensated
one-vertex induction.  It is not a counterexample to
Conjecture~\ref{g-conj-CD}: in fact
$\mathcal S(T_m)=m/12+O(1)$, and the compensation
$2\mathcal S(T_m)/(4m-1)$ tends to $1/24$, exactly the leading loss in
\eqref{g-eq-doublebroom-failure}.

For completeness, the Wiener calculation behind
\eqref{g-eq-doublebroom-values} is elementary.  In a tree, an edge whose
deletion gives component orders $s$ and $n-s$ contributes $s(n-s)$ to
$W$, because precisely those separated pairs use the edge.  Here the
$2m$ leaf edges contribute $8m^2$, and the path-edge cut sizes are
$m+j$, $1\le j\le2m$.  Hence
\[
 W(T_m)=8m^2+\sum_{j=1}^{2m}(m+j)(3m+1-j)
       =\frac{22m^3+36m^2+2m}{3},
\]
which proves the displayed average-distance formula.

\subsection{Leaf-transfer curvature and an auxiliary obstruction}
\label{g-sub-curvature}

\begin{proposition}[Exact curvature under leaf transfer]
\label{g-prop-leaf-curvature}
Fix a connected core graph $K$ with at least two vertices.  At each core
vertex $v$, attach $l_v\ge0$ pendant leaves.  Write $k_v=d_K(v)$,
$d_v=k_v+l_v$, and $n=|V(K)|+\sum_vl_v$.  Extend the resulting formulas
analytically to real nonnegative $l_v$.  Then
\begin{equation}
 \frac{\partial^2R}{\partial l_v^2}
 =\frac{3\sum_{w\sim_Kv}d_w^{-1/2}-d_v-3k_v}{4d_v^{5/2}}<0,
 \label{g-eq-curvature-diagonal}
\end{equation}
and, for distinct $u,v$,
\begin{equation}
 \frac{\partial^2R}{\partial l_u\partial l_v}
 =\begin{cases}
 (4d_u^{3/2}d_v^{3/2})^{-1},&uv\in E(K),\\
 0,&uv\notin E(K).
 \end{cases}
 \label{g-eq-curvature-cross}
\end{equation}
Along a transfer $l_u\mapsto l_u+t$, $l_v\mapsto l_v-t$, the total
order is fixed and
\begin{equation}
 R''(t)<0,\qquad W''(t)=-2d_K(u,v),\qquad
 (R-\mu)''(t)=R''(t)+\frac{4d_K(u,v)}{n(n-1)}.
 \label{g-eq-transfer-curvature}
\end{equation}
\end{proposition}

\begin{proof}
The exact index is
\[
 R=\sum_v\frac{l_v}{\sqrt{k_v+l_v}}
   +\sum_{uv\in E(K)}\frac1{\sqrt{d_ud_v}}.
\]
Differentiating gives \eqref{g-eq-curvature-diagonal} and
\eqref{g-eq-curvature-cross}.  Since $d_w\ge1$, the sum in the
numerator of \eqref{g-eq-curvature-diagonal} is at most $k_v$; the
numerator is therefore at most $-d_v<0$.  Along the transfer,
$R''=R_{uu}+R_{vv}-2R_{uv}<0$.

Let $h=|V(K)|$ and $\tau_K(v)=\sum_xd_K(v,x)$.  Counting core--core,
leaf--core, same-site leaf, and distinct-site leaf pairs gives the
polynomial
\[
 W=W(K)+\sum_v l_v(\tau_K(v)+h)
       +\sum_vl_v(l_v-1)
       +\sum_{u<v}l_ul_v(d_K(u,v)+2).
\]
Its second derivative along the transfer is
$2+2-2(d_K(u,v)+2)=-2d_K(u,v)$.  Since $n$ is fixed, differentiation
of $\mu=2W/[n(n-1)]$ proves the last identity.
\end{proof}

These formulas separate two competing effects.  The degree index is
strictly concave under a transfer, but the subtraction of average distance
adds a positive term proportional to the separation of the attachment
sites.  Consequently, the concavity of $R$ alone does not justify moving
all leaves to one hub when minimizing $R-\mu$.  A sufficient condition
for an entire transfer interval to have strictly concave $R-\mu$ is
\[
 \frac{4d_K(u,v)}{n(n-1)}
 <\frac1{4d_u(t)^{3/2}}+\frac1{4d_v(t)^{3/2}}
 \quad\text{throughout that interval}.
\]
Under this additional condition the minimum on the interval occurs at an
endpoint.  Without such a condition, the transfer argument remains
incomplete.

A separate tempting auxiliary comparison also fails.  Let
\[
 H(G)=\sum_{uv\in E(G)}\frac2{d_G(u)+d_G(v)}
\]
be the harmonic index.  Form a comb $Q_m$ by attaching one leaf to every
vertex of $P_m$, with $m\ge3$, so $|Q_m|=2m$.  Direct degree counting gives
\begin{align*}
 R(Q_m)&=\frac{m-3}{3}+\frac2{\sqrt6}
              +\frac{m-2}{\sqrt3}+\sqrt2,\\
 H(Q_m)&=\frac{m-3}{3}+\frac45+\frac{m-2}{2}+\frac43.
\end{align*}
Therefore
\[
 \lim_{m\to\infty}
 \frac{(2m-1)H(Q_m)}{2R(Q_m)^2}
 =\frac{5/6}{(1/3+1/\sqrt3)^2}
 =\frac{15}{8+4\sqrt3}>1,
\]
where the strict sign follows from $7>4\sqrt3$.  Thus
$H(G)\le2R(G)^2/(|G|-1)$ is false for infinitely many trees and cannot
be used to justify the compensated deletion candidate.

\subsection{Finite experiments and the remaining gap}
\label{g-sub-computation}

The numerical search concerned Conjecture~\ref{g-conj-CD}, not the
proved sampling theorems.  For a connected graph of order $n\ge3$, its
tested score was
\[
 \operatorname{score}(G)=
 \max_{v\ \mathrm{noncut}}
 \left(2R(G)+(n-2)(R(G)-R(G-v))
             -\frac{2\tau_G(v)}{n-1}-C_n^*\right).
\]
The conjecture requires this score to be nonnegative.  The computation
used NetworkX~3.7 and ordinary floating-point arithmetic, with tolerances
of order $10^{-8}$.  It was not an interval-arithmetic or formal proof.

\begin{center}
\begin{tabularx}{\linewidth}{@{}lXr@{}}
\toprule
Collection & Scope & Observed failures\\
\midrule
Connected graph atlas & all $994$ connected graphs, $3\le n\le7$ & $0$\\
Unlabeled trees & all $32{,}506$ trees, $3\le n\le16$ & $0$\\
Structured random tests & $600$ generated graphs, fixed seed $20261121$ & $0$\\
\bottomrule
\end{tabularx}
\end{center}

For the tree enumeration the counts at orders $10,\ldots,16$ were
$106$, $235$, $551$, $1301$, $3159$, $7741$, and $19320$, respectively.  Stars have exact zero
score.  Among the nonstars in the exhaustive atlas and tree checks, the
smallest reported score was approximately $0.0120192597$, attained by
$P_6$; its exact expression is
$2\sqrt2-6\sqrt5+53/5$.

The structured tests used path, star, cycle, complete, random-tree, and
augmented-tree cores, with core orders drawn from
$\{3,4,5,6,8,10,15,20\}$.  Leaf-load scales ranged from $2$ to $1000$;
subdivision lengths were generated using exponents $0,1/2,2/3,1$.
Attachments included leaves, paired attachments, path-like chains of
attachments, and some clique attachments.  The implementation's mode
named \texttt{cycle} joins the attachment vertices in a path, while each
also remains adjacent to the core vertex; it should not be read as an
unqualified cycle graph.  The random search stopped scanning vertices
as soon as it found a positive witness.  Consequently, its recorded
positive margins are lower bounds for the maximum score, not necessarily
the exact maxima.  The retained worst floating-point value was
$-4.44\times10^{-16}$ on the three-vertex star, consistent with rounding
of its exact zero score.  The full reproducibility script is included in
Appendix~\ref{r-graph-code}.

The computations motivate further study of
Conjecture~\ref{g-conj-CD}; they do not establish it.  The proved outcomes
of this part are Theorems~\ref{g-thm-deletion}, \ref{g-thm-sampling}, and
\ref{g-thm-cut-reduction}, their corollaries, and the exact transformation
identities and counterexamples above.  The outstanding mathematical step
is to prove the compensated noncut deletion statement, or to eliminate
cut vertices in a smallest Caporossi--Hansen counterexample by a different
argument.

\clearpage
\part{Zero correlations and arithmetic frontiers}\label{r-part-frontier}

\providecommand{\fCl}{\operatorname{Cl}}
\providecommand{\fden}{\operatorname{den}}

\section{A conditional zero-energy criterion for prime gaps}
\label{f-sec-zero-energy}

This part records a complete implication from an additional, explicitly
stated hypothesis on zeta-zero phases to primes in intervals of length
$o(\sqrt{x})$.  The Riemann hypothesis is assumed throughout this section.
The additional phase hypothesis is not proved here, and it is not identified
with RH or with Dirichlet GRH.  The subsequent sections prove a transfer
identity that isolates one of its missing estimates, and record exact
obstructions to several attempted transfers.  Their purpose is to retain
the mathematically checkable progress without treating an unproved input as
a conclusion.

We use $\psi(x)=\sum_{m\le x}\Lambda(m)$, with the usual von Mangoldt
function.  The nontrivial zeros are written $1/2+i\gamma$; ordinates are
counted with their multiplicities.  Standard zero counting and the truncated
explicit formula are used in their usual forms; see
\cite{f-titchmarsh}.  Unless otherwise stated, constants in $O$ and $\ll$
are absolute, and logarithms are natural.

\subsection{The precise extra hypothesis}
\label{f-sub-q-definition}

For an integer $n\ge2$, put
\[
 X=n^2,\qquad \eta=\frac1{50},\qquad
 I_n=[X-10n,X+10n].
\]
For real $T>1$ define
\[
 B_T(u)=\sum_{T<\gamma\le2T}u^{i\gamma}.
\]
For $L>0$, put $M=L\log T$ and
\begin{equation}
 P_{T,L}(n)=
 \sum_{\substack{T<\gamma\le2T\\T<\gamma'\le2T}}
 X^{i(\gamma-\gamma')}
 \exp\!\left(-\frac{(\gamma-\gamma')^2}{L^2\log^2T}\right).
 \label{f-eq-gaussian-energy}
\end{equation}
The sum is real and nonnegative, as will follow from
Lemma~\ref{f-lem-fourier-positive}.  The height $T$ is not restricted to
integers or to a predetermined dyadic grid.

\begin{hypothesis}[Gaussian phase condition $Q$]
\label{f-hyp-q}
With the preceding definitions,
\begin{equation}
 \lim_{L\to\infty}\ \limsup_{n\to\infty}\
 \sup_{n^{1/50}\le T\le n^2}
 \frac{P_{T,L}(n)}{LT\log T}=0.
 \label{f-eq-q}
\end{equation}
The supremum is over all real $T$ in the displayed interval.
\end{hypothesis}

In particular, $L$ is fixed before $n$ tends to infinity.  An estimate whose
threshold in $n$ depends on $L$ can be sufficient; interchanging the two
limits is not part of the hypothesis.  A stronger sufficient assertion
would be the following: there is one constant $C$ such that, for every
fixed sufficiently large $L$, all sufficiently large $n$ satisfy
\begin{equation}
 P_{T,L}(n)\le C T\log T
 \quad(n^{1/50}\le T\le n^2).
 \label{f-eq-q-prime}
\end{equation}
The threshold for $n$ may depend on $L$, but $C$ must not.  Dividing by
$LT\log T$ proves that \eqref{f-eq-q-prime} implies $Q$.

\begin{theorem}[Consequences of the Gaussian phase condition]
\label{f-thm-q-prime-gaps}
Assume RH and Hypothesis~\ref{f-hyp-q}.  Then, for each fixed real
$c\in(0,1]$,
\begin{align}
 \psi(x+c\sqrt{x})-\psi(x)&\sim c\sqrt{x},
 \label{f-eq-short-psi}\\
 \pi(x+c\sqrt{x})-\pi(x)&\sim\frac{c\sqrt{x}}{\log x}
 \qquad(x\to\infty).
 \label{f-eq-short-primes}
\end{align}
In addition,
\begin{equation}
 \pi((n+1)^2)-\pi(n^2)\sim\frac n{\log n},
 \label{f-eq-squares-asymptotic}
\end{equation}
and consecutive primes satisfy
\begin{equation}
 p_{k+1}-p_k=o(\sqrt{p_k}).
 \label{f-eq-small-prime-gaps}
\end{equation}
In particular, there is a prime between consecutive sufficiently large
integer squares.  No finite threshold or verification of the remaining
small squares is asserted.
\end{theorem}

The proof is given in full below.  The use of a positive Fourier kernel is
classical in the connection between zero correlations and short intervals;
related conditional frameworks include \cite{f-lpz,f-gwx}.  We make no
claim that the formulation above is historically optimal or new.

\subsection{Positive Fourier representation and its limitations}
\label{f-sub-positive-kernel}

\begin{lemma}[Gaussian Fourier identity]
\label{f-lem-fourier-positive}
For the quantities in \eqref{f-eq-gaussian-energy},
\begin{equation}
 P_{T,L}(n)=\frac{M}{2\sqrt\pi}
 \int_{\mathbb R}e^{-M^2v^2/4}|B_T(Xe^v)|^2\,dv\ge0.
 \label{f-eq-positive-integral}
\end{equation}
\end{lemma}

\begin{proof}
The elementary Fourier transform
\[
 \frac{M}{2\sqrt\pi}\int_{\mathbb R}
 e^{-M^2v^2/4}e^{irv}\,dv=e^{-r^2/M^2}
\]
applies to each term after expanding the square.  There are finitely many
zeros in the height block, so the interchange requires no limiting
argument.  The integrand is nonnegative.
\end{proof}

\begin{remark}[What positivity does not say]
\label{f-rem-psd-limitations}
The diagonal contribution $\#\{T<\gamma\le2T\}$ is not a lower bound for
$P_{T,L}(n)$: modulated off-diagonal terms can be negative.  Nor are
quadratic energies monotone under passage to larger principal submatrices.
For example, the positive definite matrix
\[
 K=\begin{pmatrix}1&\rho\\\rho&1\end{pmatrix},\qquad 0<\rho<1,
\]
has energy $(1,-1)K(1,-1)^{\mathsf t}=2-2\rho<2$, and this energy is less
than the one-coordinate energy $1$ when $\rho>1/2$.  This is also a Gaussian
kernel example: take $\rho=e^{-r^2/M^2}$ and a modulation with
$x^{ir}=-1$.  Thus positive semidefiniteness by itself gives neither the
desired arithmetic phase bound nor permission to discard an arbitrary
subset of zero ordinates from an energy estimate.

The proof below only restricts the nonnegative integral in
\eqref{f-eq-positive-integral}; it never uses the invalid principal-submatrix
monotonicity assertion.
\end{remark}

\begin{remark}[An operator-norm obstruction to a phase-blind estimate]
\label{f-rem-operator-norm}
Suppose $M=L\log T$, with $L>0$ fixed and $T$ large.  A subinterval of
length $M$ in the interior of $(T,2T]$ contains $\asymp M\log T$ ordinates,
by the standard zero-counting formula with $O(\log T)$ remainder.  On those
ordinates the unmodulated Gaussian matrix has all entries at least $e^{-1}$.
Testing on the constant vector shows that its operator norm is
$\gg M\log T$.  Modulation is conjugation by a diagonal unitary matrix and
does not change that norm.  Consequently an operator-norm argument that
ignores the particular arithmetic modulation cannot, by positivity alone,
supply the small energy required by $Q$.  This observation concerns that
proof strategy, not the logical consequences of RH.
\end{remark}

\begin{lemma}[Moving a Gaussian energy to a nearby square]
\label{f-lem-square-centre}
Fix any finite set of ordinates and write
\[
 \mathcal P_M(x)=\sum_{\gamma,\gamma'}
 x^{i(\gamma-\gamma')}e^{-(\gamma-\gamma')^2/M^2}.
\]
If $x,Y>0$ and $\epsilon=\log(x/Y)$, then
\begin{equation}
 \mathcal P_M(x)\le
 2e^{M^2\epsilon^2/12}\mathcal P_{M/2}(Y).
 \label{f-eq-square-comparison}
\end{equation}
\end{lemma}

\begin{proof}
Apply \eqref{f-eq-positive-integral} with this finite set and change
variables from $v$ to $w=v+\epsilon$.  The inequality
\[
 (w-\epsilon)^2\ge\frac{w^2}{4}-\frac{\epsilon^2}{3}
\]
follows by completing the square.  It bounds the Gaussian weight by
$e^{M^2\epsilon^2/12}e^{-M^2w^2/16}$.  The normalizing prefactor for the
kernel with width $M/2$ is half the original prefactor, giving the factor
$2$ in \eqref{f-eq-square-comparison}.
\end{proof}

For $Y$ a neighboring square, $|\epsilon|=O(x^{-1/2})$.  For fixed $L$ and
$M=L\log T$ in the relevant height range, the exponential factor in
\eqref{f-eq-square-comparison} tends uniformly to $1$.  More precisely,
$Q$ is equivalent to the analogous condition with real centers $x\to\infty$
and $x^{1/100}\le T\le x$.  To obtain that version from $Q$, set
$n=\lfloor\sqrt{x}\rfloor$.  Use the center $n^2$ if $T\le n^2$ and the
center $(n+1)^2$ otherwise.  In the latter case $T\asymp x$, so the lower
height bound for $n+1$ also holds.  Apply
\eqref{f-eq-square-comparison} with $L$ replaced by $L/2$ on the right.
After division by $LT\log T$ the factor $2$ cancels that replacement.
The converse follows by specializing the real center to a square.  Thus the
restriction to square centers does not materially weaken the problem at
this Gaussian resolution.

\subsection{From block energies to cumulative zero sums}
\label{f-sub-block-sums}

\begin{lemma}[Local block bound]
\label{f-lem-local-block}
For each fixed $L>0$ and all sufficiently large $n$,
\begin{equation}
 \int_{I_n}|B_T(u)|^2\,du
 \le 8\sqrt\pi\,\frac{X}{L\log T}P_{T,L}(n)
 \quad(n^{1/50}\le T\le n^2).
 \label{f-eq-local-block}
\end{equation}
Consequently, under $Q$, for every $\varepsilon>0$ one can first choose
$L$ and then $n_0(L,\varepsilon)$ so that
\begin{equation}
 \|B_T\|_{L^2(I_n)}\le C\sqrt\varepsilon\,n\sqrt T
 \quad(n\ge n_0,\ n^{1/50}\le T\le n^2).
 \label{f-eq-small-block}
\end{equation}
\end{lemma}

\begin{proof}
For $u\in I_n$, put $v=\log(u/X)$.  Uniformly there,
$|v|\le20/n$ for large $n$, $du=u\,dv\le2X\,dv$, and
$M\le2L\log n$.  Thus $e^{-M^2v^2/4}\ge1/2$ for large $n$, uniformly in
the allowed $T$.  Restricting the nonnegative integral in
\eqref{f-eq-positive-integral} to these $v$ gives
\eqref{f-eq-local-block}.  The definition of the ordered limit in $Q$
then supplies \eqref{f-eq-small-block}, with an absolute constant $C$.
\end{proof}

\begin{lemma}[Uniform cumulative-sum bound]
\label{f-lem-cumulative}
Put $S_t(u)=\sum_{0<\gamma\le t}u^{i\gamma}$.  Under $Q$, for each fixed
$C_0>0$,
\begin{equation}
 \sup_{n/C_0\le t\le n^2}
 \frac{\|S_t\|_{L^2(I_n)}}{n\sqrt t}\longrightarrow0.
 \label{f-eq-cumulative}
\end{equation}
\end{lemma}

\begin{proof}
Decompose $(0,t]$ into blocks
$(t/2^j,t/2^{j-1}]$ as long as their lower endpoint is at least
$n^{1/50}$.  Each is a complete block of the form used in
\eqref{f-eq-small-block}.  The endpoints need not be integers; this is why
$Q$ was stated for all real heights.  Minkowski's inequality gives a total
norm at most
\[
 C\sqrt\varepsilon\,n\sqrt t
 \sum_{j\ge1}2^{-j/2}\ll\sqrt\varepsilon\,n\sqrt t.
\]
The undecomposed bottom interval has upper endpoint below $2n^{1/50}$.
The standard estimate $N(y)=O(y\log(y+3))$ bounds its contribution in
$L^2(I_n)$ by
\[
 O\!\left(n^{1/2+1/50}\log n\right).
\]
Relative to $n\sqrt t$ this tends to zero uniformly for $t\ge n/C_0$.
First let $n\to\infty$ with the previously selected $L$ fixed, and then
let $\varepsilon\to0$.
\end{proof}

\subsection{A short-interval mean square and the pointwise recovery}
\label{f-sub-short-interval-proof}

\begin{lemma}[Uniform mean square near a square]
\label{f-lem-mean-square}
Let $K_n=[X-4n,X+4n]$.  For every fixed $0<c_0<c_1\le4$, RH and $Q$
imply, uniformly for $c_0n\le h\le c_1n$,
\begin{equation}
 J_n(h):=\int_{K_n}
 |\psi(u+h)-\psi(u)-h|^2\,du=o(Xh)=o(n^3).
 \label{f-eq-j-small}
\end{equation}
\end{lemma}

\begin{proof}
Choose a height $U=U_n\in[X/2,X]$ separated from zero ordinates by
$\gg1/\log X$.  Such a height exists by unit-interval zero counting.  The
truncated explicit formula, with this good height, gives uniformly on
$I_n$
\begin{equation}
 \psi(u)=u-\sum_{|\gamma|\le U}
 \frac{u^{1/2+i\gamma}}{1/2+i\gamma}+O(\log^2X).
 \label{f-eq-explicit-local}
\end{equation}
At a discontinuity one may first use the half-weighted function $\psi_0$;
the adjustment is $O(\log X)$ and is absorbed by the stated remainder.

Set $a=X/h\asymp n$.  For the positive ordinates up to $a$,
\[
 \sum_{0<\gamma\le a}
 \frac{(u+h)^{1/2+i\gamma}-u^{1/2+i\gamma}}{1/2+i\gamma}
 =\int_0^h(u+v)^{-1/2}S_a(u+v)\,dv.
\]
Every translated integration interval lies inside $I_n$.  Minkowski's
inequality and Lemma~\ref{f-lem-cumulative} therefore give
\begin{align*}
 \left\|\int_0^h(u+v)^{-1/2}S_a(u+v)\,dv\right\|_{L^2(K_n)}
 &\le C hX^{-1/2}\|S_a\|_{L^2(I_n)}\\
 &=o(h\sqrt a)=o(\sqrt{Xh}).
\end{align*}

For the higher positive ordinates, Abel summation gives the exact identity
\begin{align}
 \sum_{a<\gamma\le U}\frac{u^{i\gamma}}{1/2+i\gamma}
 &=\frac{S_U(u)}{1/2+iU}-\frac{S_a(u)}{1/2+ia}
 +i\int_a^U\frac{S_t(u)}{(1/2+it)^2}\,dt.
 \label{f-eq-abel-high}
\end{align}
There is a number $e_n\to0$, depending only on $c_0,c_1$, for which
$\|S_t\|_{L^2(I_n)}\le e_n n\sqrt t$ throughout the heights in this
identity.  Hence its $L^2(I_n)$ norm is at most
\[
 C e_n n\left(U^{-1/2}+a^{-1/2}+
 \int_a^Ut^{-3/2}\,dt\right)
 \ll e_n n a^{-1/2}.
\]
Multiplication by $u^{1/2}\ll n$ gives $o(n^2/\sqrt a)
=o(\sqrt{Xh})$.  Applying this bound at both $u$ and $u+h$ bounds their
difference in the same way.  The negative ordinates contribute conjugate
terms and introduce only a factor of two.  Finally, the explicit-formula
remainder has $L^2(K_n)$ norm $O(n^{1/2}\log^2X)=o(n^{3/2})$.
Combining these estimates and squaring proves
\eqref{f-eq-j-small}, uniformly in the stated range of $h$.
\end{proof}

\begin{lemma}[Monotonicity upgrades the mean square]
\label{f-lem-pointwise}
For each fixed $c_0>0$, RH and $Q$ imply
\begin{equation}
 \psi(x+h)-\psi(x)=h+o(n)
 \label{f-eq-pointwise-local}
\end{equation}
uniformly for
\[
 X-n\le x\le X+2n+1,\qquad c_0n\le h\le3n.
\]
\end{lemma}

\begin{proof}
Suppose first that an excess of at least $\varepsilon h$ occurs at some
admissible $(x,h)$, along a sequence $n\to\infty$.  Let
$\delta=\min\{\varepsilon/2,1/4\}$ and $H=(1+\delta)h$.
For every $u\in[x-\delta h,x]$, the interval $(u,u+H]$ contains
$(x,x+h]$.  Since $\psi$ is nondecreasing,
\[
 \psi(u+H)-\psi(u)-H\ge(\varepsilon-\delta)h.
\]
The interval of such $u$ lies in $K_n$ for large $n$.  It follows that
\[
 J_n(H)\ge\delta(\varepsilon-\delta)^2h^3\gg_{\varepsilon,c_0}n^3,
\]
contradicting Lemma~\ref{f-lem-mean-square}.

For a deficit of at least $\varepsilon h$, instead set
$H=(1-\delta)h$ and take $u\in[x,x+\delta h]$.  Now $(u,u+H]$ is
contained in $(x,x+h]$, and
\[
 \psi(u+H)-\psi(u)-H\le-(\varepsilon-\delta)h.
\]
The same contradiction follows.  All lengths $H$ used here lie in one
fixed interval $[c_0' n,4n]$, and their translates stay in $I_n$; these are
exactly the domains of the preceding lemma.
\end{proof}

\begin{proof}[Proof of Theorem~\ref{f-thm-q-prime-gaps}]
First take $x=n^2$ and $h=2n+1$ in
Lemma~\ref{f-lem-pointwise}.  Then
\[
 \psi((n+1)^2)-\psi(n^2)\sim2n.
\]
Prime powers of exponent at least two contribute only $O(\log^2X)$ in an
interval of length $O(\sqrt X)$ near $X$.  Indeed there are $O(1)$ squares
there; for every $k\ge3$ the corresponding interval for their bases has
length $O(X^{1/k-1/2})$, hence contains $O(1)$ integers.  Only $O(\log X)$
exponents can occur, each with weight at most $\log X$.  Removing those
powers and using $\log p\sim\log X=2\log n$ proves
\eqref{f-eq-squares-asymptotic}.

For real $x\to\infty$, put $n=\lfloor\sqrt x\rfloor$.  Then
$n^2\le x\le n^2+2n+1$.  For a fixed $c\in(0,1]$, the length
$h=c\sqrt x$ lies between fixed positive multiples of $n$ and is at most
$3n$ for large $n$.  Equation~\eqref{f-eq-pointwise-local} gives
\eqref{f-eq-short-psi}.  The same prime-power estimate proves
\eqref{f-eq-short-primes}.  Finally, for every fixed $\varepsilon>0$,
\eqref{f-eq-short-primes}, with $c=\min\{\varepsilon,1\}$, places a
prime in $(p_k,p_k+c\sqrt{p_k}]$ for all sufficiently large $k$.
This is precisely \eqref{f-eq-small-prime-gaps}.
\end{proof}

\section{A positive-kernel transfer and its unresolved arithmetic estimate}
\label{f-sec-transfer}

This section derives an identity between one zero-energy scale and the
mean square of an explicitly weighted prime Dirichlet series.  The identity
is proved; the strong upper bound subsequently requested of that mean
square is not.  In this section only, the center and the height are coupled
by $X=T^2$.  Establishing a bound at this one scale would still leave the
other heights in Hypothesis~\ref{f-hyp-q} to be treated.

\subsection{Gaussian and Cauchy kernels}
\label{f-sub-cauchy}

Let $T$ be large, $X=T^2$, $1\le a\le T/10$, and $M=2a$.  Write
$I=(T,2T]$ and
\[
 B(v)=\sum_{\gamma\in I}e^{i\gamma(\log X+v)}.
\]
For clarity the parameters in the following two energies are suppressed:
\begin{align}
 P_M&=\sum_{\gamma,\gamma'\in I}
 X^{i(\gamma-\gamma')}e^{-(\gamma-\gamma')^2/M^2},
 \label{f-eq-transfer-p}\\
 F_M&=\sum_{\gamma,\gamma'\in I}
 X^{i(\gamma-\gamma')}
 \frac{M^2}{M^2+(\gamma-\gamma')^2}.
 \label{f-eq-transfer-f}
\end{align}

\begin{lemma}[A positive comparison]
\label{f-lem-cauchy-comparison}
Both energies are nonnegative, and
\begin{equation}
 P_M\le\frac e{\sqrt\pi}F_M.
 \label{f-eq-cauchy-comparison}
\end{equation}
\end{lemma}

\begin{proof}
Fourier inversion gives
\[
 P_M=\frac M{2\sqrt\pi}\int_{\mathbb R}
 e^{-M^2v^2/4}|B(v)|^2\,dv,\qquad
 F_M=\frac M2\int_{\mathbb R}e^{-M|v|}|B(v)|^2\,dv.
\]
For $s\ge0$, $-s^2/4+s\le1$.  Therefore
$e^{-M^2v^2/4}\le e\,e^{-M|v|}$, proving the comparison of the two
nonnegative integrals.  No entrywise comparison of signed modulated
matrix sums is being used.
\end{proof}

\subsection{The exact explicit-formula identity}
\label{f-sub-series-identity}

For $a>1/2$, define the absolutely convergent Dirichlet series
\begin{equation}
 D_{X,a}(t)=\sum_{m\ge2}\Lambda(m)m^{-1/2-it}
 e^{-a|\log(m/X)|}.
 \label{f-eq-long-series}
\end{equation}
It has an infinite, decaying tail; calling it a finite Dirichlet polynomial
would require an additional truncation and error estimate.  Under RH put
\begin{equation}
 A_{X,a}(t)=\sum_{\gamma\text{ of either sign}}
 \frac{2a}{a^2+(t-\gamma)^2}X^{i(\gamma-t)}.
 \label{f-eq-a-all}
\end{equation}
This series is absolutely convergent by zero counting.

\begin{proposition}[Explicit residual and uniform domain]
\label{f-prop-explicit-transfer}
Assume RH.  For $t\in[T,2T]$, $X=T^2$, and $1\le a\le T/10$,
\begin{equation}
 A_{X,a}(t)=-D_{X,a}(t)+R_{X,a}(t),
 \label{f-eq-a-d}
\end{equation}
where the residual is exactly
\begin{align}
 R_{X,a}(t)={}&
 \frac{2aX^{1/2-it}}{a^2-(1/2-it)^2}
 -\sum_{k\ge1}
 \frac{2aX^{-2k-1/2-it}}{a^2-(-2k-1/2-it)^2}
 \notag\\[-2pt]
 &-X^{-a}\frac{\zeta'}{\zeta}(1/2-a+it).
 \label{f-eq-residual-exact}
\end{align}
Uniformly on this domain,
\begin{equation}
 R_{X,a}(t)\ll X^{-a}\log T+\frac{a\sqrt X}{T^2}
 =X^{-a}\log T+\frac aT.
 \label{f-eq-residual-bound}
\end{equation}
\end{proposition}

\begin{proof}
The Mellin transform of $w_a(y)=e^{-a|\log y|}$ is
\[
 K_a(s)=\int_0^\infty w_a(y)y^{s-1}\,dy
 =\frac{2a}{a^2-s^2},\qquad |\Re s|<a.
\]
Choose $1/2<c<a$ and set $s_0=1/2+it$.  Absolute convergence and Mellin
inversion give
\[
 D_{X,a}(t)=\frac1{2\pi i}\int_{\Re s=c}
 K_a(s)X^s\left(-\frac{\zeta'}{\zeta}(s_0+s)\right)ds.
\]
Shift the contour to the left, avoiding its poles.  The zeta pole contributes
$K_a(1-s_0)X^{1-s_0}$.  The nontrivial zeros contribute
$-K_a(\rho-s_0)X^{\rho-s_0}$, which is the negative of
\eqref{f-eq-a-all} under RH.  The trivial zeros give the displayed negative
sum in \eqref{f-eq-residual-exact}.  Finally $K_a$ has residue $1$ at
$s=-a$, yielding $-X^{-a}(\zeta'/\zeta)(s_0-a)$.  The shifted integrals
vanish on a standard sequence of rectangles: $K_a(s)=O(|s|^{-2})$ at
large height, the factor $X^{\Re s}$ decays on the left, and the usual
logarithmic-derivative bounds apply away from poles.  The residue series
are absolutely convergent.  Rearranging proves the exact identity.

For the bound, the pole term is $O(a\sqrt X/T^2)$.  In the trivial-zero
term, if $b=-2k-1/2$, then
\[
 |a^2-(b-it)^2|=|a-b+it|\,|a+b-it|\ge t^2,
\]
so its sum is negligible compared with the pole bound.  The functional
equation reflects $1/2-a+it$ to a point with real part $1/2+a\ge3/2$.
The logarithmic derivative at that reflected point is bounded by an
absolutely convergent Dirichlet series.  The gamma-factor derivative is
$O(\log T)$ uniformly for $a\le T/10$ and $T\le t\le2T$.
This proves \eqref{f-eq-residual-bound}.
\end{proof}

\begin{theorem}[Transfer to a weighted prime mean square]
\label{f-thm-mean-square-transfer}
Assume RH.  Uniformly for $X=T^2$ and $1\le a\le T/10$,
\begin{equation}
 F_{2a}=\frac a{2\pi}\int_T^{2T}|D_{T^2,a}(t)|^2\,dt
 +O(a^2\log^3T).
 \label{f-eq-transfer-mean-square}
\end{equation}
The implied constant is absolute.  In particular, for each fixed $L>0$
and $a=(L/2)\log T$, the remainder is $O_L(\log^5T)$.
\end{theorem}

\begin{proof}
Let $A_I$ denote the sum in \eqref{f-eq-a-all} restricted to $\gamma\in I$.
Write $C_a(u)=2a/(a^2+u^2)$.  Convolution of Poisson kernels gives
\[
 \int_{\mathbb R}C_a(t-\gamma)C_a(t-\gamma')\,dt
 =\frac{8\pi a}{4a^2+(\gamma-\gamma')^2}.
\]
Expanding the finite square and cancelling the common factor $X^{-it}$
therefore gives the exact identity
\begin{equation}
 F_{2a}=\frac a{2\pi}\int_{\mathbb R}|A_I(t)|^2\,dt.
 \label{f-eq-poisson-exact}
\end{equation}

The standard bound of $O(\log(|t|+3))$ zeros in a unit interval implies
\[
 |A_{X,a}(t)|+|A_I(t)|\ll\log T\qquad(T\le t\le2T).
\]
Summing the Poisson tails of the discarded ordinates also gives
\begin{equation}
 |A_{X,a}(t)-A_I(t)|\ll\log T
 \left(\frac a{a+t-T}+\frac a{a+2T-t}\right).
 \label{f-eq-boundary-tail}
\end{equation}
For example, subdivide each tail into unit intervals and compare the
resulting sum with the integral of $a/(a^2+u^2)$.  Ordinates farther than
$T$ from this height interval contribute $O(a\log T/T)$ and are absorbed
by the right-hand side.  Hence
\begin{align*}
 \int_T^{2T}\big||A_{X,a}(t)|^2-|A_I(t)|^2\big|\,dt
 &\ll a\log^2T\log(1+T/a)\\
 &\ll a\log^3T.
\end{align*}
The same tail argument, now for the fixed finite set $I$, gives
\[
 \int_{\mathbb R\setminus[T,2T]}|A_I(t)|^2\,dt\ll a\log^2T.
\]

Finally, \eqref{f-eq-a-d}, \eqref{f-eq-residual-bound}, and
$|A_{X,a}|\ll\log T$ imply
\[
 \int_T^{2T}\big||D_{X,a}(t)|^2-|A_{X,a}(t)|^2\big|\,dt
 \ll a\log^2T.
\]
For instance, use $\big||A-R|^2-|A|^2\big|\le2|A||R|+|R|^2$,
$X=T^2$, and $a\ge1$; the term $TX^{-a}\log^2T$ is then harmless.
Insert these three estimates into \eqref{f-eq-poisson-exact} and retain
the prefactor $a/(2\pi)$.  The resulting error is $O(a^2\log^3T)$.
\end{proof}

\subsection{The off-diagonal term that is still needed}
\label{f-sub-off-diagonal}

Put
\[
 b_m=\Lambda(m)m^{-1/2}e^{-a|\log(m/X)|}.
\]
Absolute convergence of the series justifies the exact expansion
\begin{align}
 \int_T^{2T}|D_{X,a}(t)|^2dt
 ={}&T\sum_{m\ge2}b_m^2
 \notag\\
 &+2\Re\sum_{h\ge1}\sum_{m\ge2}b_mb_{m+h}
 \int_T^{2T}e^{it\log(1+h/m)}dt.
 \label{f-eq-prime-pair-expansion}
\end{align}
For fixed $L>0$, $X=T^2$, and $a=(L/2)\log T$, partial summation with
the RH estimate $\psi(y)=y+O(\sqrt y\log^2 y)$ gives
\begin{equation}
 \sum_m b_m^2\sim\frac{\log X}{a}=\frac4L,
 \qquad
 \sum_m mb_m^2\ll\frac{X\log X}{a}.
 \label{f-eq-diagonal-evaluation}
\end{equation}
For the first relation, the main integral is
$\int (\log y)y^{-1}e^{-2a|\log(y/X)|}\,dy$; putting
$v=\log(y/X)$ gives $(\log X)/a$ up to a negligible endpoint error.
The contribution of higher prime powers and the prime-counting remainder
are negligible on the window of relative width $1/a\asymp1/\log T$.

A sufficient unproved estimate at this scale is
\begin{equation}
 \int_T^{2T}|D_{T^2,a}(t)|^2dt
 \le C\frac{T\log T}{a},\qquad a=\frac L2\log T,
 \label{f-eq-missing-mean-square}
\end{equation}
where $C$ is independent of fixed $L$ and the threshold in $T$ may depend
on $L$.  Theorem~\ref{f-thm-mean-square-transfer} and
Lemma~\ref{f-lem-cauchy-comparison} would then imply
$P_M\ll T\log T$ at $X=T^2$.  The oscillatory factors in
\eqref{f-eq-prime-pair-expansion} and the products
$\Lambda(m)\Lambda(m+h)$ are the missing arithmetic information.
Separations $h\asymp X/T=T$ already affect the oscillation over this
time interval, while the coefficient window has width $X/a\gg T$.

The general Hilbert-inequality mean-value bound
\cite{f-montgomery-vaughan} gives only
\begin{equation}
 \int_T^{2T}|D_{X,a}(t)|^2dt
 =T\sum_m b_m^2+O\!\left(\sum_m mb_m^2\right)
 \ll(T+X)\frac{\log X}{a}.
 \label{f-eq-hilbert-loss}
\end{equation}
At $X=T^2$ this misses \eqref{f-eq-missing-mean-square} by a factor of
order $T$.  The termwise absolute bound in the zero representation instead
gives $|D_{X,a}|\ll\log T$, hence only
$F_M\ll aT\log^2T\asymp L T\log^3T$.  Neither estimate proves $Q$.

\subsection{Why character averaging does not recover the principal term}
\label{f-sub-character-average}

Choose a prime $q\asymp T$ and define
\[
 D_\chi(t)=\sum_{m\ge2}b_m\chi(m)m^{-it}.
\]
Orthogonality of all Dirichlet characters modulo $q$ gives
\begin{equation}
 \frac1{\varphi(q)}\sum_{\chi\bmod q}\int_T^{2T}|D_\chi(t)|^2dt
 =\sum_{\substack{m\equiv k\pmod q\\(mk,q)=1}}
 b_mb_k\int_T^{2T}e^{it\log(k/m)}dt.
 \label{f-eq-character-average}
\end{equation}
This identity is exact.  Apply the separated-frequency mean-value
inequality to the sequence $\log(r+jq)$ in each reduced residue class.
The inverse neighboring frequency gap is $O(1+m/q)$, so
\begin{align}
 \frac1{\varphi(q)}\sum_{\chi\bmod q}\int_T^{2T}|D_\chi(t)|^2dt
 &\ll(T+1)\sum_m b_m^2+q^{-1}\sum_m mb_m^2
 \notag\\
 &\ll\frac{T\log X}{a}.
 \label{f-eq-character-average-bound}
\end{align}
Thus an average over characters really is of the desired diagonal order.
This part of the argument needs no GRH character estimate.

However, nonnegativity gives for the principal character only
\[
 \int_T^{2T}|D_{\chi_0}(t)|^2dt
 \le\varphi(q)\left\{
 \frac1{\varphi(q)}\sum_\chi\int_T^{2T}|D_\chi(t)|^2dt\right\}.
\]
The factor $\varphi(q)\asymp T$ restores the loss.  The distinction between
$D$ and $D_{\chi_0}$ is harmless here: their difference is supported on
powers of $q$ and is $O(\log T/T)$ when $q\asymp T$, $X=T^2$, and
$a\ge1$.  The unresolved issue is the extraction of a single character
from the average, not those omitted coefficients.  No argument in this
investigation removes that factor using scalar GRH character-sum bounds.

\section{Exact domain checks for two further transfer attempts}
\label{f-sec-domain-checks}

This section concerns the hypotheses of proposed applications, not an audit
of the accepted RH premise.  A source theorem can be correct while not
being applicable to the continuous shifts or quadratic energies required
above.  The checks below rule out the literal substitutions attempted in
this investigation; they do not rule out a new theorem with a larger
domain.

\subsection{Angular infinity types, radial shifts, and logarithmic metrics}
\label{f-sub-theta-domain}

The source files cited here were inspected at the fixed public mirror
commit in \cite{f-lean-source}.  In
\texttt{IsUnitaryNormOf.lean}, the predicates \texttt{HasInf} and
\texttt{HasInfQ} require an infinity component of the form
\begin{equation}
 \chi_\infty(k)=\prod_\phi
 \left(\frac{\phi(k)}{|\phi(k)|}\right)^{n_\phi},
 \qquad n_\phi\in\mathbb Z,
 \label{f-eq-angular-type}
\end{equation}
with $n_\phi=2a_\phi$ for prescribed rational weights in
\texttt{HasInfQ}.  The \texttt{IsMu0} host in \texttt{CharHost.lean}
requires that fixed angular type together with global triviality and
specified local finite-place conditions.

For positive rational $k$, the right-hand side of
\eqref{f-eq-angular-type} is $1$.  A nonzero radial norm twist at the
infinite places has the form $k^{it[E:\mathbb Q]}$ and cannot be
identically $1$ on all positive rationals unless $t=0$.  For example,
requiring this simultaneously for $k=2,3$ would make $\log2/\log3$
rational if $t\ne0$, contradicting unique factorization.  Thus the direct
insertion of arbitrary real shifts $it$ does not satisfy this particular
host's infinity-type hypothesis.  The availability of more general Hecke
characters at a definitional level does not change that conclusion about
the theorem being applied.

The predicate \texttt{TaylorLogMetric.IsChi}\footnote{Declared in
\texttt{IsDiscriminantTwelfthRootOfIdealLattice.lean}.} is defined on
norm-one fractional ideals with a fixed principal-unitary ideal law.
Radial twists are invisible on a norm-one domain, so this does not supply
the missing continuous parameter.  Also,
\texttt{TaylorLogMetric.th22} bounds a Galois mean of logarithms of norm
squares.  A bound on a logarithmic mean alone does not bound an individual
quadratic energy: even two positive numbers $e^A,e^{-A}$ have mean
logarithm zero while their maximum is unbounded.

The inspected class-sum Cauchy--Schwarz statements retain explicit
cardinality and class-number factors.  They do not state cancellation of
the arithmetic phases in \eqref{f-eq-prime-pair-expansion}.  Finally, the
``ratio'' in \texttt{KernelRatioClosedWithCharacters.lean} is the
algebraic-jet kernel-rationality property
\texttt{StandardLatticesGL3.kRat}; it is not a ratios asymptotic for shifted
$L$-functions.  These are concrete type and conclusion mismatches, not
claims that no different use of the source could yield a stronger estimate.

\subsection{The support restriction in a recent pair-correlation theorem}
\label{f-sub-pair-correlation-support}

Wang's Theorem 2.2 in \cite{f-wang} fixes
$0<\lambda<\theta<1$, a short height window $H=T^\theta$, and a fixed
even smooth test function supported in $[-\lambda,\lambda]$.  Its
error depends on that test function.  In the normalization
\[
 u=(\gamma-\gamma')\frac{\log T}{2\pi},
\]
the modulation at $X=n^2$ translates Fourier frequency to
\begin{equation}
 \alpha=\frac{\log X}{\log T}=\frac{2\log n}{\log T}.
 \label{f-eq-frequency-shift}
\end{equation}
For the heights required by $Q$, one has $1\le\alpha\le100$, outside
the strict support band of that theorem.  In addition, the Gaussian with
$M=L\log T$ has Fourier width of order $1/(L\log^2T)$ in this
normalization.  It is neither a fixed test function nor compactly
supported.  The changing-test error would have to be controlled even if
the support problem were overcome.  Therefore this theorem cannot be
substituted directly for the missing phase hypothesis.  The wider
ratios-conjecture approach in \cite{f-gwx} also uses additional inputs;
it is not a statement that GRH alone supplies the desired correlations.

\section{A determinant investigation of the Gieseking volume}
\label{f-sec-gieseking}

This section records a different attempted passage from recent
irrationality methods to a named constant.  The positive results here are
exact determinant identities and a conditional norm criterion.  The
numerical tables are exploratory.  No proof that the Gieseking volume is
irrational is claimed.

\subsection{The constant and the precise arithmetic target}
\label{f-sub-volume-constant}

Let $\chi_{-3}$ be the nonprincipal character modulo $3$, and put
\begin{equation}
 L=L(2,\chi_{-3})=
 \sum_{k\ge0}\left(\frac1{(3k+1)^2}-\frac1{(3k+2)^2}\right),
 \qquad \xi=\frac L4.
 \label{f-eq-l-constant}
\end{equation}
For the Clausen function
$\fCl_2(\theta)=\sum_{k\ge1}\sin(k\theta)/k^2$, the regular ideal
tetrahedron volume, also the Gieseking volume, is
\begin{equation}
 V=\fCl_2(\pi/3)=\frac{3\sqrt3}{4}L=3\sqrt3\,\xi.
 \label{f-eq-volume-relation}
\end{equation}
The geometric identification and related arithmetic context are discussed
in \cite{f-cdt}.  The series identity can be verified directly: splitting
the two series into residue classes modulo $6$, the even part of the
$L$-series is $-L/4$, so its odd part is $5L/4$.  The sine series is
$\sqrt3/2$ times the odd part minus the even part, giving
\eqref{f-eq-volume-relation}.

Calegari, Dimitrov, and Tang prove that $1,\zeta(2),L(2,\chi_{-3})$
are linearly independent over $\mathbb Q$ \cite{f-cdt}.  This implies
$L\notin\mathbb Q$ but does not imply $V\notin\mathbb Q$.  If instead
$V=A/B\in\mathbb Q$, with $B\ne0$, then
\begin{equation}
 \xi=\frac{A\sqrt3}{9B},\qquad
 \xi^\sigma=-\xi,\qquad
 \xi^2=\frac{A^2}{27B^2},
 \label{f-eq-pure-quadratic-target}
\end{equation}
where $\sigma$ is the nontrivial automorphism of $\mathbb Q(\sqrt3)$.
The required exclusion is therefore a pure quadratic value, not just a
rational value of $L$.

Litsey's November 2026 preprint \cite{f-litsey} gives a Loewner-determinant
framework with explicit estimates for $L(2,\chi_{-3})$ and reports formal
verification of parts of its argument.  Its stated scope does not include
irrationality of the Clausen volume in \eqref{f-eq-volume-relation}.
We do not certify that preprint or its formal development here.  Its
reported rates are used below to diagnose what a proposed extension would
still need, while all elementary identities and the norm criterion are
proved independently.

\subsection{An explicit rational determinant family}
\label{f-sub-determinant-family}

For $k\ge1$, define the rational numbers
\begin{align}
 g(k)&=\frac1{3k-2}-\frac1{2k-1}+\frac1{6k},
 \label{f-eq-g}\\
 g'(k)&=-\frac3{(3k-2)^2}+\frac2{(2k-1)^2}-\frac1{6k^2}.
 \label{f-eq-g-prime}
\end{align}
The prime on $g'$ denotes the displayed derivative of the rational
function, evaluated at $k$.  Define a rational symmetric matrix on
nonnegative integer nodes by
\begin{equation}
 M[t,u]=
 \begin{cases}
 \displaystyle\frac1{6(u-t)}\sum_{k=t+1}^{u}g(k),&t<u,\\[6pt]
 \displaystyle\frac16\sum_{k=1}^{t}g'(k),&t=u,\\[6pt]
 M[u,t],&t>u.
 \end{cases}
 \label{f-eq-rational-matrix}
\end{equation}
For integers $r\ge1$, $m\ge0$, let
\begin{align*}
 C&=\{0,1,\ldots,r-1\},\\
 R&=C\cup\{r,r+2,\ldots,r+2m-2\},\\
 Q&=C\cup\{r+1,r+3,\ldots,r+2m-1\},
\end{align*}
where the added sets are empty when $m=0$.  Put
\begin{equation}
 P_{r,m}(x)=\det\bigl(M[t,u]+x\,\mathbf1_{t=u}\bigr)_{t\in R,u\in Q}
 \in\mathbb Q[x].
 \label{f-eq-determinant-polynomial}
\end{equation}
Its degree is at most $r$, since only the $r$ common nodes contain the
variable.  The geometry used for the principal diagnostic table is even
$r$ with $m=\lceil r/4-1/2\rceil$.

\begin{proposition}[Connection to the digamma Loewner matrix]
\label{f-prop-digamma-matrix}
Let $\psi_\Gamma(z)=\Gamma'(z)/\Gamma(z)$ denote the digamma function,
and define
\[
 F_\chi(z)=-\frac13\psi_\Gamma(z+1/3)
 +\frac12\psi_\Gamma(z+1/2)-\frac16\psi_\Gamma(z+1),
 \qquad \omega=-\frac16.
\]
For its Loewner matrix
\[
 \mathcal L_{F_\chi}[t,u]=
 \begin{cases}
 (F_\chi(u)-F_\chi(t))/(u-t),&t\ne u,\\
 F_\chi'(t),&t=u,
 \end{cases}
\]
one has, at nonnegative integer nodes,
\begin{equation}
 \omega\mathcal L_{F_\chi}[t,u]=M[t,u]+\xi\mathbf1_{t=u}.
 \label{f-eq-loewner-identity}
\end{equation}
\end{proposition}

\begin{proof}
The digamma recurrence gives
$F_\chi(k)-F_\chi(k-1)=-g(k)$.  Its differentiated version gives
$F_\chi'(k)-F_\chi'(k-1)=-g'(k)$.  These prove the off-diagonal
identity and reduce the diagonal identity to
$\omega F_\chi'(0)=\xi$.
The trigamma series gives
\[
 \psi_\Gamma'(1/3)-\psi_\Gamma'(2/3)=9L,
 \qquad
 \psi_\Gamma'(1/3)+\psi_\Gamma'(2/3)=\frac{4\pi^2}{3}.
\]
Together with $\psi_\Gamma'(1/2)=\pi^2/2$ and
$\psi_\Gamma'(1)=\pi^2/6$, these yield
$F_\chi'(0)=-3L/2$, and multiplication by $-1/6$ gives $L/4$.
\end{proof}

\subsection{Parity and reflection do not supply the other conjugate}
\label{f-sub-parity}

\begin{proposition}[Failure of a uniform parity identity]
\label{f-prop-no-parity}
The determinant family \eqref{f-eq-determinant-polynomial} does not obey
a general identity $P_{r,m}(-x)=\pm P_{r,m}(x)$.  In particular,
\begin{equation}
 P_{2,0}(x)=x^2-\frac7{36}x-\frac1{1296}.
 \label{f-eq-parity-counterexample}
\end{equation}
\end{proposition}

\begin{proof}
Here the matrix is
\[
 \begin{pmatrix}x&1/36\\1/36&x-7/36\end{pmatrix}.
\]
Taking its determinant gives \eqref{f-eq-parity-counterexample}, which
has both a nonzero odd coefficient and nonzero even coefficients.
\end{proof}

Additional exact rational checks give the following ratio of the
coefficient of $x^{r-1}$ to the leading coefficient.  Each is nonzero:
\begin{center}
\begin{tabular}{r r l}
\hline
$r$&$m$&$[x^{r-1}]P_{r,m}/[x^r]P_{r,m}$\\
\hline
2&0&$-7/36$\\
4&1&$-1029911/3175200$\\
6&1&$-2956587752267/3895779888000$\\
8&2&$-12927501768602788895963/8681618633814846720000$\\
\hline
\end{tabular}
\end{center}
These finite checks do not exclude a new construction with a specially
designed parity symmetry.  They do exclude treating such a symmetry as an
identity already possessed by these undeleted determinants.

There is also an exact reflection check.  Set
$F(z)=\psi_\Gamma(z+1/3)-\psi_\Gamma(z+2/3)$.  The digamma reflection
formula gives
\begin{equation}
 F(-z)=-F(z)+\pi\{\cot\pi(z+2/3)-\cot\pi(z+1/3)\}.
 \label{f-eq-reflection}
\end{equation}
Differentiating, with the minus sign from $-z$ retained, gives
\begin{equation}
 F'(-z)=F'(z)+\pi^2\{
 \csc^2\pi(z+2/3)-\csc^2\pi(z+1/3)\}.
 \label{f-eq-reflection-derivative}
\end{equation}
At every integer $j$ the two cosecant squares coincide, so
\begin{equation}
 F'(-j)=F'(j).
 \label{f-eq-reflection-integers}
\end{equation}
Thus reflecting integer nodes does not simply replace the constant on the
derivative diagonal by its negative.  A two-contour or doubled determinant
would need a new proof of its arithmetic denominator, analytic decay, and
nonvanishing; the reflection identity alone supplies none of those three.

\subsection{The exact number-field norm requirement}
\label{f-sub-norm-barrier}

\begin{proposition}[A sufficient conjugate-decay criterion]
\label{f-prop-norm-criterion}
Let $P_r\in\mathbb Q[x]$ have degree $d_r=O(r)$, and suppose positive
integers $D_r$ satisfy
\[
 D_rP_r\in\mathbb Z[x],\qquad
 \log D_r\le(\Lambda+o(1))r^2.
\]
Suppose, for infinitely many $r$, that $P_r(\xi)\ne0$ and
\begin{equation}
 |P_r(\xi)|\le e^{-(D_+-o(1))r^2},\qquad
 |P_r(-\xi)|\le e^{-(D_--o(1))r^2}.
 \label{f-eq-two-decays}
\end{equation}
If $D_++D_->2\Lambda$, then $3\sqrt3\,\xi$ is not rational.
Here the assertion concerns the positive constant $\xi$ in
\eqref{f-eq-l-constant}.
\end{proposition}

\begin{proof}
Assume for a contradiction that $3\sqrt3\,\xi=A/B\in\mathbb Q$.
By \eqref{f-eq-pure-quadratic-target}, $\xi\in\mathbb Q(\sqrt3)$,
$\xi^\sigma=-\xi$, and one can choose a fixed positive integer $b$
such that $b\xi$ is an algebraic integer.  Then
\[
 Z_r=b^{d_r}D_rP_r(\xi)
\]
is an algebraic integer: write it as an integer polynomial in $b\xi$
using the factors $b^{d_r-j}$ in its $j$th coefficient.  When
$P_r(\xi)\ne0$, its field norm is a nonzero integer, so
\[
 1\le|N_{\mathbb Q(\sqrt3)/\mathbb Q}(Z_r)|
 =b^{2d_r}D_r^2|P_r(\xi)P_r(-\xi)|.
\]
Taking logarithms, using $d_r=O(r)$, and applying
\eqref{f-eq-two-decays}, gives
\[
 0\le\{2\Lambda-D_+-D_-+o(1)\}r^2,
\]
which is impossible along the asserted infinite set of indices.
\end{proof}

The denominator and positive-embedding rates reported for the relevant
common-node construction in \cite{f-litsey} are
\begin{equation}
 \Lambda_\chi=4.09592776731\ldots,
 \qquad
 D_+=\frac{580530026928040140331}{125\cdot10^{18}}
 =4.644240215424321\ldots.
 \label{f-eq-published-rates}
\end{equation}
Their use in the direct norm argument would require a negative-conjugate
decay rate strictly larger than
\begin{equation}
 2\Lambda_\chi-D_+=3.54761531919\ldots.
 \label{f-eq-needed-conjugate-rate}
\end{equation}
It would also require infinitely many nonzero values at the quadratic
number in \eqref{f-eq-pure-quadratic-target}.  A rational-point
nonvanishing theorem for a family of determinants with a bounded number
of common nodes deleted does not automatically prove that quadratic-point
statement.  These are separate missing steps.

The factor $2$ in the denominator cost can sometimes be reduced by
cancellation in the norm polynomial.  Write
\begin{equation}
 P(x)=E(x^2)+xO(x^2),\qquad
 Q_P(y)=E(y)^2-yO(y)^2.
 \label{f-eq-norm-polynomial}
\end{equation}
Then $Q_P(x^2)=P(x)P(-x)$.  If $\xi^2=u/v\in\mathbb Q$ in lowest
terms and $D_P^{\mathrm N}=\fden(Q_P)$ is the least positive integer
clearing its coefficients, a nonzero value satisfies
\begin{equation}
 |P(\xi)P(-\xi)|\ge
 (D_P^{\mathrm N}v^{\deg Q_P})^{-1}.
 \label{f-eq-rational-norm-lower}
\end{equation}
This follows by evaluating the integer polynomial
$D_P^{\mathrm N}Q_P$ at $u/v$ and clearing the denominator
$v^{\deg Q_P}$.  Accordingly, a proved asymptotic bound for
$\log D_P^{\mathrm N}$ could improve the direct criterion.  It must be
computed or proved; it cannot be inferred from the smallness of the one
real embedding.

\subsection{Exploratory conjugate and border computations}
\label{f-sub-numerical-tables}

Define the finite-size decay indicators
\[
 d_\pm(r,m)=-\frac1{r^2}\log|P_{r,m}(\pm\xi)|.
\]
The following values were recomputed from
\eqref{f-eq-rational-matrix}--\eqref{f-eq-determinant-polynomial}, using
the trigamma expression for $\xi$.  The precision column is the number
of working decimal digits.  These are ordinary high-precision evaluations,
not interval-certified inequalities, asymptotic limits, or proofs of
eventual nonvanishing.

\begin{center}
\begin{tabular}{r r r r r}
\hline
$r$&$m$&Precision&$d_+(r,m)$&$d_-(r,m)$\\
\hline
12&3&200&4.369672777293&0.527530842951\\
24&6&288&4.515953548401&0.467340154958\\
36&9&432&4.566849912358&0.455555250719\\
48&12&576&4.592732360927&0.448394597748\\
\hline
\end{tabular}
\end{center}
The observed negative-conjugate values are far below the sufficient rate
in \eqref{f-eq-needed-conjugate-rate}.  This is evidence about this
construction at these sizes, not a theorem excluding a better construction
or a different proof.

We also varied the border length $m$.  For each polynomial let
$D_{r,m}^{\mathrm N}=\fden(Q_{P_{r,m}})$ be computed exactly over the
rationals, and define
\begin{equation}
 s_{r,m}=\frac{\log D_{r,m}^{\mathrm N}
 +\log|P_{r,m}(\xi)P_{r,m}(-\xi)|}{r^2}.
 \label{f-eq-border-score}
\end{equation}
A successful asymptotic norm argument of this kind would need a negative
limit superior of the corresponding quantity, together with nonvanishing;
the additional fixed denominator $v^{\deg Q_P}$ in
\eqref{f-eq-rational-norm-lower} contributes only $O(r)$ to the logarithm.
All scores in the following finite test are positive.  Longer borders can
improve analytic decay while increasing denominator cost even more.
For $m$ outside the published geometry, the constants in
\eqref{f-eq-published-rates} are not being asserted.

\begin{center}
\small
\begin{longtable}{r r r r r r}
\hline
$r$&$m$&$r^{-2}\log D_{r,m}^{\mathrm N}$&$d_+$&$d_-$&$s_{r,m}$\\
\hline
\endfirsthead
\hline
$r$&$m$&$r^{-2}\log D_{r,m}^{\mathrm N}$&$d_+$&$d_-$&$s_{r,m}$\\
\hline
\endhead
4&0&3.887090&2.325204&0.279232&1.282653\\
4&1&6.364314&3.792965&0.899450&1.671900\\
4&2&10.122258&5.591023&1.715981&2.815254\\
4&4&20.352230&10.351985&5.149711&4.850534\\
4&8&53.549689&24.830090&17.387696&11.331903\\
\hline
6&0&4.361997&2.546725&0.176330&1.638943\\
6&1&5.983557&3.533842&0.498131&1.951584\\
6&3&10.237426&5.932201&1.665649&2.639576\\
6&6&21.488762&10.827640&5.079352&5.581770\\
6&12&54.785916&25.335433&17.439684&12.010800\\
\hline
8&0&4.664819&2.672253&0.128559&1.864007\\
8&2&7.005661&4.232433&0.556096&2.217132\\
8&4&10.591857&6.140690&1.567331&2.883837\\
8&8&21.254303&10.987105&5.053344&5.213853\\
8&16&55.485942&25.636858&17.481746&12.367337\\
\hline
\end{longtable}
\end{center}

In these fifteen exact rational calculations,
$\log D_{r,m}^{\mathrm N}=2\log\fden(P_{r,m})$: cancellation in the
norm polynomial did not lower the denominator.  This is a finite
observation only.  The displayed decimals use 200 working digits for
their analytic part, with exact rational coefficients for the denominator.

\subsection{Reproducible formulas and computational method}
\label{f-sub-reproducibility}

The following Python code specifies the matrix, exact polynomial, and
norm-polynomial computations used above.  It uses \texttt{mpmath} for
the exploratory real evaluations and \texttt{sympy} for exact rational
arithmetic.  The Schur complement in the exact-polynomial routine was
invertible in the tested cases; a singular border can instead be handled
by computing the determinant polynomial directly.  Nothing in the
conditional propositions depends on these computations.

\begin{lstlisting}[language=Python,basicstyle=\ttfamily\scriptsize,breaklines=true]
import math
import mpmath as mp
import sympy as sp

def nodes(r, m):
    common = list(range(r))
    return (common + [r+2*j for j in range(m)],
            common + [r+2*j+1 for j in range(m)])

def matrix_mp(r, m, x):
    R, C = nodes(r, m)
    G, H = [mp.mpf(0)], [mp.mpf(0)]
    for k in range(1, r+2*m):
        G.append(G[-1] + 1/mp.mpf(3*k-2)
                 - 1/mp.mpf(2*k-1) + 1/mp.mpf(6*k))
        H.append(H[-1] - 3/mp.mpf(3*k-2)**2
                 + 2/mp.mpf(2*k-1)**2 - 1/(6*mp.mpf(k)**2))
    return mp.matrix([
        [H[t]/6+x if t == u else
         (G[u]-G[t])/(6*(u-t)) for u in C] for t in R])

def exact_poly(r, m):
    R, C = nodes(r, m)
    G, H = [sp.Rational(0)], [sp.Rational(0)]
    for k in range(1, r+2*m):
        G.append(G[-1] + sp.Rational(1,3*k-2)
                 - sp.Rational(1,2*k-1) + sp.Rational(1,6*k))
        H.append(H[-1] - sp.Rational(3,(3*k-2)**2)
                 + sp.Rational(2,(2*k-1)**2)
                 - sp.Rational(1,6*k*k))
    A = sp.Matrix([
        [H[t]/6 if t == u else
         (G[u]-G[t])/(6*(u-t)) for u in C] for t in R])
    x = sp.Symbol('x')
    if not m:
        return sp.Poly((-A).charpoly(x).as_expr(), x)
    border = A[r:,r:]
    S = A[:r,:r] - A[:r,r:]*border.inv()*A[r:,:r]
    return sp.Poly(border.det()*(-S).charpoly(x).as_expr(), x)

def norm_poly(P):
    y = sp.Symbol('y')
    E = sum(P.nth(j)*y**(j//2)
            for j in range(0, P.degree()+1, 2))
    O = sum(P.nth(j)*y**((j-1)//2)
            for j in range(1, P.degree()+1, 2))
    return sp.Poly(E**2-y*O**2, y)

def denominator(P):
    return sp.ilcm(*[c.q for c in P.all_coeffs()])

for r in (12,24,36,48):
    m = math.ceil(r/4-0.5)
    mp.mp.dps = max(200,12*r)
    xi = (mp.polygamma(1,mp.mpf(1)/3)
          - mp.polygamma(1,mp.mpf(2)/3))/36
    dp = -mp.log(abs(mp.det(matrix_mp(r,m,xi))))/r**2
    dm = -mp.log(abs(mp.det(matrix_mp(r,m,-xi))))/r**2
    print(r, m, mp.nstr(dp,18), mp.nstr(dm,18))

# Border table: use 200 decimal digits and, for each r in (4,6,8),
# m in (0, ceil(r/4-.5), r//2, r, 2*r).
# Compute P=exact_poly(r,m), DN=denominator(norm_poly(P)), and
# s=log(DN)/r**2-dp-dm.  All denominators are exact integers.
\end{lstlisting}

GRH improves some error terms in prime-distribution estimates that occur
in denominator analysis.  Such an improvement does not by itself change
the leading quadratic denominator constant for a fixed valuation profile,
nor does it bound $P_{r,m}(-\xi)$ or prove quadratic-point nonvanishing.
Those are the additional mathematical inputs needed for this proposed
norm strategy.  The tables and exact parity test identify that gap; they
do not fill it.

\section{A finite diagnostic for an Erd\H{o}s--Straus transfer}
\label{f-sec-es-diagnostic}

One attempted use of a small quadratic nonresidue required a divisor of
$K^2$ in one of two specified classes modulo $q$.  The finite test
$p=61$, $q=7$, $K=(p+q)/4=17$ already shows why the nonresidue condition
alone cannot guarantee the tested divisor condition: the divisors of
$K^2=289$ have residues
\[
 \{1,17,289\}\pmod7=\{1,3,2\},
\]
whereas the two required residues in that test were $4$ and $5$.
Explicitly, the tested channels were
$d\equiv-K$ or $d\equiv-Kp^{-1}\pmod q$ with $d\mid K^2$;
here $p^{-1}\equiv3\pmod7$, giving those two residues.
The statement being refuted is this proposed sufficient-condition
transfer, not the Erd\H{o}s--Straus conjecture.  Indeed the same prime has
the explicit valid decomposition
\[
 \frac4{61}=\frac1{16}+\frac1{366}+\frac1{2928}.
\]
This example is retained only as a short, reproducible diagnostic.  No
general obstruction to Erd\H{o}s--Straus follows from it.

\subsection*{Status of the frontier results}

Theorem~\ref{f-thm-q-prime-gaps} is a proved conditional implication with
the extra hypothesis \eqref{f-eq-q} displayed explicitly.  The positive
kernel comparison, residual formula, and transfer theorem are established
on their stated domains.  Estimate~\eqref{f-eq-missing-mean-square}, the
full phase condition $Q$, the required conjugate-decay and nonvanishing
inputs for the volume, and any new conclusion on the named open problems
remain unproved in this investigation.  The determinant numbers are
exploratory computations, while the parity and reflection identities and
the finite divisor example are exact calculations.


\clearpage
\appendix
\section{Reproducibility of the compensated-deletion experiments}\label{r-graph-code}
The following Python program records the local-deletion calculation and
the finite searches used for the compensated-deletion candidate in
Part~\ref{r-part-graph}. It uses Python 3 and NetworkX 3.7. The options
\texttt{--atlas}, \texttt{--trees 16}, and \texttt{--adversarial 600}
select the respective tests. The finite computations use floating-point
arithmetic and are evidence about a conjecture, not a proof of it.

For a proposed noncut vertex the program evaluates
\[
 2R(G)+(n-2)\bigl(R(G)-R(G-v)\bigr)
 -\frac{2\tau_v}{n-1}
 -\left(nB_n-(n-2)B_{n-1}\right),
 \qquad B_n=\sqrt{n-1}-2+\frac2n.
\]
The adversarial mode may stop after finding one positive witness. Its
reported positive value is then a positive numerical witness,
not necessarily the maximum over all noncut vertices. A
putative failure is examined against every inequivalent noncut vertex.
The retained run has 600 cases and a smallest recorded value of about
$-4.44\cdot10^{-16}$ on the three-vertex star, a roundoff-level equality
case. The random seed is fixed in the source below.

\begin{lstlisting}[language=Python]

import networkx as nx
import math, random, json, argparse, time

def R(G):
    d=G.degree()
    return math.fsum(1/math.sqrt(d[u]*d[v]) for u,v in G.edges())

def B(n):
    return math.sqrt(n-1)-2+2/n

def threshold(n):
    s,t=math.sqrt(n-1),math.sqrt(n-2)
    return 2*s+(n-2)/(s+t)-4+2/(n-1)

def delta_R(G,v):
    ds=G.degree(); ns=set(G[v]); dv=ds[v]
    removed=math.fsum(1/math.sqrt(dv*ds[u]) for u in ns)
    changed={tuple(sorted((u,w))) for u in ns for w in G[u] if w != v}
    extra=math.fsum(1/math.sqrt((ds[u]-(u in ns))*(ds[w]-(w in ns)))-1/math.sqrt(ds[u]*ds[w]) for u,w in changed)
    return removed-extra

def check(G, full=False, existence_only=False):
    n=len(G); rr=R(G); cuts=set(nx.articulation_points(G)); C=threshold(n)
    seen=set(); vals=[]
    for v in sorted(G, key=lambda w: (G.degree(w)!=1, G.degree(w))):
        if v in cuts: continue
        # Equal open neighborhoods imply equal deletion and distance data.
        sig=frozenset(G[v])
        if sig in seen: continue
        seen.add(sig)
        tau=sum(nx.single_source_shortest_path_length(G,v).values())
        delta=delta_R(G,v)
        val=2*rr+(n-2)*delta-2*tau/(n-1)-C
        vals.append((val,v,delta,tau))
        if existence_only and val > 1e-8: return vals[-1]
    ans=max(vals)
    return (ans,vals) if full else ans

def loaded_core(core, loads, lengths=None, block_mode="leaf"):
    h=len(core); G=nx.Graph();G.add_nodes_from(range(h)); nxt=h
    for i,(u,v) in enumerate(core.edges()):
        length=1 if lengths is None else lengths[i]
        last=u
        for j in range(length-1):
            G.add_edge(last,nxt);last=nxt;nxt+=1
        G.add_edge(last,v)
    for u,load in enumerate(loads):
        verts=list(range(nxt,nxt+load));nxt+=load
        G.add_edges_from((u,w) for w in verts)
        if block_mode=="clique":
            G.add_edges_from((x,y) for x in verts for y in verts if x<y)
        elif block_mode=="pairs":
            G.add_edges_from((verts[j],verts[j+1]) for j in range(0,len(verts)-1,2))
        elif block_mode=="cycle" and len(verts)>1:
            G.add_edges_from(zip(verts[:-1],verts[1:]))
    return G

def audit_atlas():
    worst=(float("inf"),None);cnt=0
    for i,G in enumerate(nx.graph_atlas_g()):
        n=len(G)
        if n<3 or not nx.is_connected(G):continue
        score=check(G)[0];cnt+=1
        if score<worst[0]:worst=(score,(i,n))
        assert score > -1e-8,(i,score)
    print("atlas",cnt,"worst",worst,flush=True)

def audit_trees(max_n):
    for n in range(3,max_n+1):
        worst=(float("inf"),None);cnt=0
        for G in nx.nonisomorphic_trees(n):
            score=check(G)[0];cnt+=1
            if score<worst[0]:worst=(score,list(G.edges()))
            assert score > -1e-8,(n,score,list(G.edges()))
        print("trees",n,cnt,"worst",worst[0],flush=True)

def adversarial(seed_count=500):
    rng=random.Random(20261121);worst=(float("inf"),None)
    for it in range(seed_count):
        h=rng.choice([3,4,5,6,8,10,15,20]);typ=rng.randrange(6)
        if typ==0:core=nx.path_graph(h)
        elif typ==1:core=nx.star_graph(h-1)
        elif typ==2:core=nx.cycle_graph(h)
        elif typ==3:core=nx.complete_graph(h)
        elif typ==4:core=nx.random_labeled_tree(h,seed=it)
        else:
            core=nx.random_labeled_tree(h,seed=it)
            for j in range(h//2):u,v=rng.sample(range(h),2);core.add_edge(u,v)
        scale=rng.choice([2,5,10,30,100,300,1000])
        load_mode=rng.randrange(5)
        if load_mode==0:loads=[scale]*h
        elif load_mode==1:loads=[rng.randrange(0,2*scale+1) for _ in range(h)]
        elif load_mode==2:loads=[max(0,round(scale*rng.random()**4)) for _ in range(h)]
        elif load_mode==3:loads=[rng.choice([0,1,scale,2*scale]) for _ in range(h)]
        else:loads=[0]*h;loads[rng.randrange(h)]=h*scale
        total=h+sum(loads)
        exponent=rng.choice([0,.5,2/3,1]);mult=rng.choice([.2,.5,1,2,5])
        base=max(1,round(mult*total**exponent/max(1,len(core.edges()))))
        lengths=[max(1,round(base*rng.choice([.5,1,1,1,2]))) for _ in core.edges()]
        mode=rng.choice(["leaf","leaf","leaf","pairs","cycle"])
        if it%20==0 and scale<=30:mode="clique"
        G=loaded_core(core,loads,lengths,mode)
        ans=check(G,existence_only=True);score=ans[0]
        desc=dict(iteration=it,n=len(G),hub_count=h,core=list(core.edges()),loads=loads,lengths=lengths,mode=mode,ans=ans)
        if score<worst[0]:
            worst=(score,desc);print("new worst",json.dumps(worst),flush=True)
        if score < -1e-7:
            with open("ch_compensated_counterexample.json","w") as f:json.dump(dict(desc=desc,edges=list(G.edges())),f)
            print("FAIL",json.dumps(desc),flush=True);return
        if it%25==0:print("progress",it,"n",len(G),"worst",worst[0],flush=True)
    with open("ch_compensated_adversarial_results.json","w") as f:json.dump(dict(count=seed_count,worst=worst),f)
    print("ALL PASSED",seed_count,json.dumps(worst),flush=True)

if __name__=="__main__":
    p=argparse.ArgumentParser();p.add_argument("--atlas",action="store_true");p.add_argument("--trees",type=int,default=0);p.add_argument("--adversarial",type=int,default=0);a=p.parse_args()
    if a.atlas:audit_atlas()
    if a.trees:audit_trees(a.trees)
    if a.adversarial:adversarial(a.adversarial)


\end{lstlisting}

\section{Coverage and interpretation of the consolidated results}\label{r-coverage}
This is a single mathematical manuscript, rather than a concatenation of
earlier PDFs. The two preceding prime-product notes are consolidated in
Part~\ref{r-part-prime}: repeated definitions and analytic inputs have
been merged, while the one-third proof, its sharpness examples,
all-density and vanishing-density results, compression arguments,
exact-cardinality and window variants, two-prime estimates, residue-graph
theorem, and finite certificates are retained. The source-premise record
and the attributed Mills-type consequence are retained as well.

Part~\ref{r-part-graph} contains the complete deletion and sampling
proofs developed subsequently, the exact gluing calculations, the
minimal-counterexample reduction, and the unproved compensated-deletion
target. Part~\ref{r-part-frontier} contains the conditional Gaussian
prime-gap implication, the positive-kernel transfer, the scope checks on
attempted analytic inputs, the Gieseking determinant investigation, and
the finite Erd\H{o}s--Straus transfer diagnostic. Numerical diagnostics
are not promoted to asymptotic assertions.

The complete mathematical conclusions are those stated with their
hypotheses in the theorem, lemma, proposition, and corollary environments.
The additional Gaussian correlation hypothesis, the compensated graph
deletion candidate, the full two-prime covering problem, and the
Gieseking irrationality goal remain separate open tasks within this
investigation. In particular, this compilation does not assert a proof
of any of those remaining tasks.

\clearpage
\phantomsection
\addcontentsline{toc}{section}{References}
\begin{thebibliography}{99}\small\raggedright

\bibitem{p-JMV}
D. Jao, S. D. Miller, and R. Venkatesan,
\emph{Expander graphs based on GRH with an application to elliptic curve
cryptography}, J. Number Theory \textbf{129} (2009), 1491--1504.
\href{https://arxiv.org/abs/0811.0647}{arXiv:0811.0647};
\href{https://arxiv.org/html/0811.0647v2}{primary full text}.

\bibitem{p-MV}
H. L. Montgomery and R. C. Vaughan,
\emph{Multiplicative Number Theory I: Classical Theory},
Cambridge Studies in Advanced Mathematics 97, Cambridge University
Press, 2007. Theorem 13.7.

\bibitem{p-KST}
O. Klurman, I. E. Shparlinski, and J. Ter\"av\"ainen,
\emph{On the generation of multiplicative groups by small primes},
preprint, September 2026.
\href{https://arxiv.org/abs/2609.25393}{arXiv:2609.25393}.

\bibitem{p-Hughes}
S. D. Hughes, \emph{Few small primes suffice to cover unit groups},
preprint, October 2026.
\href{https://arxiv.org/abs/2610.04678}{arXiv:2610.04678}.

\bibitem{p-BP}
A. R. Booker and C. Pomerance,
\emph{Squarefree smooth numbers and Euclidean prime generators},
Proc. Amer. Math. Soc. \textbf{145} (2017), 5035--5042.
\href{https://arxiv.org/abs/1607.01557}{arXiv:1607.01557}.

\bibitem{p-MS}
M. Munsch and I. E. Shparlinski,
\emph{On smooth square-free numbers in arithmetic progressions},
J. London Math. Soc. (2020), doi:10.1112/jlms.12297.
\href{https://arxiv.org/abs/1710.04705}{arXiv:1710.04705}.

\bibitem{p-Vu}
V. H. Vu, \emph{Structure of large incomplete sets in abelian groups},
Combinatorica \textbf{30} (2010), 225--237.
\href{https://arxiv.org/abs/math/0610158}{arXiv:math/0610158}.

\bibitem{p-HB}
D. R. Heath-Brown, \emph{A mean value estimate for real character sums},
Acta Arith. \textbf{72} (1995), 235--275.
\href{https://doi.org/10.4064/aa-72-3-235-275}{doi:10.4064/aa-72-3-235-275}.
Theorem 1.

\bibitem{p-Cameron}
P. J. Cameron, \emph{The random graph revisited},
in European Congress of Mathematics, Barcelona 2000, Vol. II,
Birkh\"auser, Basel, 2001, pp. 267--274.
\href{https://www.math.uni-bielefeld.de/~rehmann/ECM/cdrom/3ecm/pdfs/pant3/camer.pdf}{Primary manuscript}.

\bibitem{p-EOS1987}
P. Erd\H{o}s, A. M. Odlyzko, and A. S\'ark\"ozy,
On the residues of products of prime numbers,
Period. Math. Hungar. 18 (1987), 229--239.
\url{https://www.renyi.hu/~p_erdos/1987-22.pdf}.

\bibitem{p-Szabo2024}
B. Szab\'o, On the existence of products of primes in arithmetic progressions,
Bull. London Math. Soc. 56 (2024), 1227--1243.
\url{https://doi.org/10.1112/blms.12990}.

\bibitem{p-Banks2026}
W. D. Banks, Products of two elements of Granville's random model in residue classes,
arXiv:2609.37434 (2026).
\url{https://arxiv.org/abs/2609.37434}.

\bibitem{p-Montgomery1971}
H. L. Montgomery, Topics in Multiplicative Number Theory,
Lecture Notes in Mathematics 227, Springer, 1971.

\bibitem{p-SS2017}
I. E. Shparlinski and A. V. Sutherland,
Finding elliptic curves with a subgroup of prescribed size,
Int. J. Number Theory 13 (2017), 133--152.
\url{https://arxiv.org/abs/1403.7887}.

\bibitem{p-KLM2026}
N. Kandhil, A. Languasco, and P. Moree,
Pair correlation of zeros of Dirichlet L-functions: a possible path towards
the conjectures of Chowla, Elliott-Halberstam and Montgomery,
Math. Ann. 394 (2026), article 43.
\url{https://doi.org/10.1007/s00208-026-03383-y}.

\bibitem{p-Saito2026}
K. Saito, Transcendency of variants of Mills' constant,
Ramanujan J. 70 (2026), article 69; arXiv:2508.16068v3, revised 7 December 2025.
\url{https://arxiv.org/abs/2508.16068}.

\bibitem{p-ST2022}
K. Saito and W. Takeda,
Topological properties and algebraic independence of sets of prime-representing constants,
Mathematika 68 (2022), 429--453.
\url{https://arxiv.org/abs/2112.14383}.

\bibitem{g-BE1998}
B.~Bollob\'as and P.~Erd\H os,
\emph{Graphs of extremal weights},
Ars Combinatoria \textbf{50} (1998), 225--233.
\url{https://combinatorialpress.com/article/ars/Volume%20050/volume-50-paper-18.pdf}.

\bibitem{g-CH2000}
G.~Caporossi and P.~Hansen,
\emph{Variable neighborhood search for extremal graphs. 1. The AutoGraphiX system},
Discrete Mathematics \textbf{212} (2000), 29--44.
\url{https://doi.org/10.1016/S0012-365X(99)00206-X}.

\bibitem{g-FMS2003}
O.~Favaron, M.~Mah\'eo, and J.-F.~Sacl\'e,
\emph{The Randi\'c index and other Graffiti parameters of graphs},
MATCH Communications in Mathematical and in Computer Chemistry
\textbf{47} (2003), 7--23; see Conjecture~4.8.
\url{https://match.pmf.kg.ac.rs/electronic_versions/Match47/match47_7-23.pdf}.

\bibitem{g-DLS2011}
Z.~Dvo\v r\'ak, B.~Lidick\'y, and R.~\v Skrekovski,
\emph{Randi\'c index and the diameter of a graph},
European Journal of Combinatorics \textbf{32} (2011), 434--442.
\url{https://doi.org/10.1016/j.ejc.2010.12.002}.
Author manuscript: \url{https://kam.mff.cuni.cz/~bernard/pub/index.pdf}.

\bibitem{g-VMRC227}
Shaitetsu,
\emph{A proof audit of Fajtlowicz's average-distance conjecture:
weighted distances, the harmonic index, and the extremal star},
research manuscript, 21 November 2026, entry~227 of the
Verifiable Mathematical Reasoning Corpus.
\url{https://github.com/QEEDDD/Verifiable-Mathematical-Reasoning-Corpus}.
This is the external source of the weighted inequality explicitly assumed
in Proposition~\ref{g-prop-rooted-weighted}; its proof is not included here.

\bibitem{f-titchmarsh}
E.~C. Titchmarsh,
\emph{The Theory of the Riemann Zeta-Function},
second edition, revised by D.~R. Heath-Brown,
Oxford University Press, Oxford, 1986.

\bibitem{f-montgomery-vaughan}
H.~L. Montgomery and R.~C. Vaughan,
Hilbert's inequality,
\emph{Journal of the London Mathematical Society} (2)
\textbf{8} (1974), 73--82.
\url{https://doi.org/10.1112/jlms/s2-8.1.73}.

\bibitem{f-lpz}
A. Languasco, A. Perelli, and A. Zaccagnini,
An extended pair-correlation conjecture and primes in short intervals,
arXiv:1311.0597, 2013.
\url{https://arxiv.org/abs/1311.0597}.

\bibitem{f-gwx}
H. Gou, V.~Y. Wang, and M.~W. Xu,
Landau's problems through ratios of $L$-functions,
preprint, arXiv:2611.18825v1, November 2026;
also available as Hexagon:2611.00423v1.
\url{https://arxiv.org/abs/2611.18825}.
\url{https://hexagonmath.org/2611.00423v1}.

\bibitem{f-wang}
B. Wang,
Simple critical zeros and distinct zeros of the Riemann zeta-function in
short intervals,
preprint, arXiv:2609.07918v1, September 2026.
\url{https://arxiv.org/abs/2609.07918}.
The support and fixed-test restrictions discussed here are those of
Theorem~2.2 in the inspected version.

\bibitem{f-lean-source}
SINF-KEN,
\emph{rh-lean-check}, public source mirror,
commit \href{https://github.com/SINF-KEN/rh-lean-check/tree/e7c8972c701d3396947808ab20cc16dc14623d33}{\texttt{e7c8972c701d}}.
\url{https://github.com/SINF-KEN/rh-lean-check}.
The present discussion cites only the domains and conclusions of the
specific source declarations named in the text.

\bibitem{f-cdt}
F. Calegari, V. Dimitrov, and Y. Tang,
The linear independence of $1$, $\zeta(2)$, and $L(2,\chi_{-3})$,
arXiv:2408.15403, 2024.
\url{https://arxiv.org/abs/2408.15403}.
Author-hosted primary text:
\url{https://www.math.uchicago.edu/~fcale/papers/L2chi.pdf}.

\bibitem{f-litsey}
S. Litsey,
Irrationality measures from L\"owner determinants,
preprint, Hexagon:2610.00401v3, 10 November 2026.
\url{https://hexagonmath.org/pdf/2610.00401v3}.
The constants used in the diagnostic are the rates stated for the
$L(2,\chi_{-3})$ construction in the inspected version; no independent
verification of the entire preprint is claimed here.


\end{thebibliography}
\end{document}
